\documentclass[11pt]{article}
\usepackage[T1]{fontenc}
\usepackage[utf8]{inputenc}
\usepackage{lmodern}
\usepackage[margin=2.5cm]{geometry}
\usepackage{amsmath,amssymb,amsthm}
\usepackage{enumitem}
\usepackage{microtype}
\usepackage[colorlinks=true,linkcolor=blue,urlcolor=blue]{hyperref}
\usepackage{xurl}
\usepackage{listings,xcolor}
\lstset{basicstyle=\ttfamily\small,breaklines=true,columns=fullflexible,frame=single,keepspaces=true,upquote=true,literate={≤}{{$\le$}}1 {≥}{{$\ge$}}1 {−}{{-}}1 {×}{{$\times$}}1}
\providecommand{\tightlist}{\setlength{\itemsep}{0pt}\setlength{\parskip}{0pt}}
\setlength{\parindent}{0pt}
\setlength{\parskip}{0.6em}
\allowdisplaybreaks

\title{Almost-linear-time lossy min-cost flow with prescribed demands}
\date{}

\begin{document}
\maketitle

The proof uses a finite-word decremental path gate for the selector. Section 9 gives the short-threshold cover, its integer implementation and the numerical accounting of its graph routines.

\hypertarget{numerical-model-and-conversion-to-a-one-sided-answer}{%
\section{Numerical model and conversion to a one-sided answer}\label{numerical-model-and-conversion-to-a-one-sided-answer}}

Let \(N=m+n+2\), with \(m\) the original arc count. Thus \(N\) charges the initial scan and the zero records for discarded arcs. In the reduction below all arc counts are bounded by \(m\). For the analysis, take \(W\) from the theorem on the positive-capacity arcs and write \(B=1+\lceil\log_2 N\rceil+\lceil\log_2(2+NW/\delta)\rceil\). The algorithm determines its working precision from normalized quantities. A polynomial in \(B\) is admissible since any fixed power of \(\log N\) is \(N^{o(1)}\).

Use input references with exact sign and zero tests, exact order comparisons of cost magnitudes, and windowed signed-exponent and mantissa access. More explicitly, the reference for a positive nonzero value has an exact floor binary exponent. It can be reused as a scale or compared to another exponent. For an integer window size \(L\), a difference of two referenced exponents, including specified integer offsets, is returned exactly when in \([-L,L]\) and otherwise flagged with its direction. The input service also provides enclosing dyadic intervals for bounded normalized data and for a positive ratio clipped at a specified dyadic bound. For a requested dyadic mesh \(t\), ask for endpoints on the mesh \(t/4\) with width at most \(t/2\). Such endpoints can enclose a value at a mesh boundary. The endpoints use \(O(1+\log_2(2+\text{range bound})+\log_2(2+1/t))\) bits, and only polynomial work in the requested precision and exponent window is charged. Taking the lower endpoint and clamping it to zero when the quantity is nonnegative gives a stored dyadic lower approximation with error at most \(t/2<t\). This construction uses no exact comparison of a real quantity to a mesh boundary in the calling algorithm. Reading full exponent fields, full significands, or evaluating real comparison requests from an arbitrary raw encoding is charged separately if precision access is unavailable. A global power of two can be returned by an original exponent reference and an integer offset. We use that representation also in the output on nonnegative-cost arcs. This convention matters if the entire input capacity scale has an arbitrarily large negative binary exponent. All local numerical algorithm operations after access use finite rational words, or explicitly rounded binary arithmetic defined below.

The enclosing-interval requests can also be implemented from the exponent windows and mantissas. For a ratio of two positive inputs, the normalized mantissas are in \([1,2)\), so their ratio lies in \([1/2,2]\). A window of size \(O(1+\log_2(2+\text{range bound})+\log_2(2+1/t))\) suffices to identify ratios certainly below \(t/4\) or certainly above the clipping bound, when one is used. Return the interval \([0,t/4]\) in the first case and enclose the clipping bound on the mesh in the second. In the remaining window, divide enclosing mantissa intervals with a positive lower denominator, apply the known power-of-two shift, and clip the endpoints. For a ratio with a power of two the denominator mantissa is exactly one. Taking a sufficiently large constant times the window size in mantissa bits makes the resulting interval narrower than \(t/4\). Rounding its ends outward then gives width at most two mesh units of size \(t/4\). Negate and exchange endpoints for signed normalized data. Zero is handled by its sign and zero test.

Here is the reduction. Its additive subroutine takes an exactly feasible bounded dyadic box LP with opposing-sign binary columns and with unary or zero columns allowed. It returns a dyadic feasible-box vector whose balance error in max norm and upward objective error are each at most its specified tolerance \(\eta\). Its required bit bound is \(O(B)\) per input entry for the calls in this section. For a matrix \(A\) write \(\|A\|_{\infty\to\infty}=\max_v\sum_e|A_{ve}|\), with zero for an empty matrix.

For each assembled additive LP compute \(\widehat B\) as one plus \(\lceil\log_2 N\rceil\) plus the maximum bit length of explicitly expanded numerators and denominators of \(A,b,c,u,\eta\). If \(\widehat B>N\), solve the bounded rational LP exactly by Grötschel--Lovász--Schrijver Theorem 6.4.12. Round its optimum downward inside its box on a binary mesh at most \(\eta/[4(1+m_{\rm box})(1+\|A\|_{\infty\to\infty}+\|c\|_\infty)]\), where \(m_{\rm box}\) is its variable count. Each coordinate changes by at most the mesh, so both additive errors are at most \(\eta\). Listing the problem, finding the exact optimum and rounding cost \((N+\widehat B)^{O(1)}=\widehat B^{O(1)}\) in this case, with an absolute exponent. Otherwise use §§2--7. The same exact rational solve and rounding procedure is the scheduled fallback of §7.3. \href{https://www.zib.de/userpage/groetschel/pubnew/paper/groetschellovaszschrijver1988.pdf\#page=192}{Grötschel--Lovász--Schrijver, Theorem 6.4.12}

Discard zero-capacity arcs and recover their output as zero. If no arcs remain, feasibility implies zero demands and the result follows. Put \(z_e=u_e\) if \(c_e<0\), and \(z_e=0\) otherwise. Let \(D\) be diagonal with entries \(-1\) if \(c_e<0\), and \(+1\) otherwise. Set \(U=2^{1+\max_e\lfloor\log_2 u_e\rfloor}\). Thus \(\max_e u_e\le U\le 2\max_e u_e\), and we need only a reference for its exponent. The original imbalance matrix is \(A\), with at most two entries and absolute entries at most one per column, including after combining a self-loop into a unary. In the variables \(f=z+UDy\), the true normalized equation is \(A'y=b\), with \(A'=AD\) and \(b=(d-Az)/U\). The box is \(0\le y_e\le u_e/U\le1\). With \(a_e=|c_e|\), the cost is \(c^\top z+Ua^\top y\). Let \(C^*\ge0\) be the optimum of \(a^\top y\) in this normalized problem.

Take \(\tau\) to be the largest power of two at most \(\min\{1/256,\delta/(256U)\}\). Compare the referenced exponent of \(\delta\) to that of \(U\). If their difference is negative and falls below the requested window, double the window until the difference is returned exactly. A flagged positive difference suffices immediately for the cap \(1/256\). Because \(U\) is an exact power of two, the returned difference gives the floor exponent of \(\delta/U\). This uses window size \(O(B)\) and fixed polynomial work in \(B\) without taking \(B\) as input. Take \(h\) to be the largest power of two at most \(\tau/(64N)\), so \(h>\tau/(128N)\). Use the lower approximation with \(t=h\) on \(u_e/U\) and on each gain, giving \(0\le\bar u_e\le u_e/U\) and \(0\le\bar\gamma_e\le1\). For \(\overline{d/U}\) use the signed lower endpoints. All errors are at most \(h\). By true feasibility \(\|d/U\|_\infty\le m\), so this is a bounded request. Form \(\bar A\) from the exact tails and the single rounded gain on each arc. On a self-loop its sole entry is \(\bar\gamma_e-1\), with error at most \(h\). Set \(\bar A'=\bar A D\) and \(\bar b=\overline{d/U}-\bar A\bar z\), where \(\bar z_e=\bar u_e\) on negative-cost arcs and zero elsewhere. For each column, the gain change contributes at most \(h\) in max norm when its coefficient is at most one in magnitude. Replacing \(z_e/U\) by \(\bar z_e\) contributes a further at most \(h\), since the absolute entries of \(\bar A\) are at most one. The rowwise triangle inequality gives
\[
 \|b-\bar b\|_\infty\le(1+2m)h,\qquad
 \|(A'-\bar A')y\|_\infty\le mh \quad (0\le y\le u/U).
\]
For an optimal true \(y^*\), put \(\bar y_e^*=\min(y_e^*,\bar u_e)\). Each coordinate decreases by at most \(h\), so \(\|\bar A'(\bar y^*-y^*)\|_\infty\le mh\). Equation \(A'y^*=b\) and the preceding two bounds give rounded residual at most \((1+4m)h<\tau/8\). The transformed cost can only decrease because \(a_e\ge0\). Delete rounded zero widths, holding those normalized coordinates at zero. For every vector in the rounded box the preceding two bounds give \(\|A'y-b\|_\infty\le\|\bar A'y-\bar b\|_\infty+(1+3m)h\). These estimates apply irrespective of the absolute scale of \(U\). If all capacities are far below \(\delta\), the cap in the definition of \(\tau\) is used and the normalized mantissas require only \(O(\log N)\) places.

\textbf{Cost screening.} Sort the positive \(a_e\) of the remaining columns by exact magnitude. Consecutive equal magnitudes form one cost level, and the levels are in strictly increasing order. A prefix \(S\) contains all zero-cost columns and the first \(j\) complete positive levels, where \(j=0\) is allowed. For a prefix solve
\[
\begin{split}
 \min\ &\sum_v(p_v^++p_v^-),\\
 &\bar A'_S y+r^+-r^-+p^+-p^-=\bar b,\\
 &0\le y\le\bar u_S,\qquad 0\le r^\pm\le\tau,\qquad 0\le p^\pm\le4N .
\end{split}
\]
It is exactly feasible at \(y=r^\pm=0\) with \(p^\pm\) the positive and negative parts of \(\bar b\), since \(\|\bar b\|_\infty\le2m+h\). Use tolerance \(\tau/64\) for both errors. Check the returned box bounds, the objective bound \(\sum_v(p_v^++p_v^-)\le\tau/8\), and the direct residual bound \(\|\bar A'_S y-\bar b\|_\infty<2\tau\) by exact dyadic arithmetic. Pass precisely when these checks hold and retain \(y\) as the witness. For a valid additive answer, the objective bound implies the residual bound since \(|r_v^+-r_v^-|\le\tau\), \(|p_v^+-p_v^-|\le\sum_v(p_v^++p_v^-)\), and the enlarged equation has error at most \(\tau/64\). A box vector of residual at most \(\tau\) extends to an exact feasible point with \(p^\pm=0\), by placing each signed residual in the appropriate \(r\)-coordinate. The optimum penalty is then zero, and the additive answer passes with objective at most \(\tau/64\).

The full prefix therefore passes by the clipped optimal vector. Test \(j=0\) first. A pass gives the desired answer already by \(a^\top y=0\), with the demand bound from the residual check and the rounding estimate above. If there are no positive levels, this is the full prefix and necessarily passes on a valid additive answer. Otherwise, on rejecting \(j=0\), test and store the full prefix, then bisect between the rejecting zero prefix and the passing full prefix. On each pass replace the passing endpoint and its witness, and on each rejection replace the rejecting endpoint. This uses \(O(\log N)\) calls and needs only the two endpoint certifications, even if tolerance-region answers vary. At adjacency let \(\sigma>0\) be the level added to the passing prefix and let \(y^0\) be its retained witness. Then \(a^\top y^0\le m\sigma\). Restrict \(\bar y^*\) to the rejecting prefix. Its residual is at most \(\tau/8+\sum_{a_e\ge\sigma}\bar y_e^*\), by the column bound on \(\bar A'\). As a residual of at most \(\tau\) would have certified a pass, this sum exceeds \(7\tau/8\). Consequently
\[
 C^*\ge\sigma\sum_{a_e\ge\sigma}\bar y_e^*>(7/8)\sigma\tau.
\tag{1}
\]

\textbf{Truncation and the strict inequality.} Let \(K_{\rm c}\) be the least power of two at least \(8N/\tau\). Put \(\theta=h\), \(\eta=\theta\tau/64\), and let \(h_2\) be the largest power of two at most \(\eta/N\). For each remaining column use the lower approximation with \(t=h_2\) for \(\min(a_e/\sigma,2K_{\rm c})\), assigning exact zero when \(a_e=0\). Denote these approximations by \(\widetilde a_e\) and retain \(H=\{e:\widetilde a_e\le K_{\rm c}\}\). The lower bound property keeps every \(a_e\le K_{\rm c}\sigma\). If \(e\in H\), its true clipped ratio is less than \(K_{\rm c}+h_2<2K_{\rm c}\), so no clipping occurred and \(0\le a_e/\sigma-\widetilde a_e<h_2\).

There is a vector on \(H\) with rounded residual less than \(2\tau\) and true normalized cost at most \(C^*\). To prove existence, if \(C^*\ge m\sigma\), take \(y^0\), whose support is included in \(H\). If \(C^*<m\sigma\), restrict \(\bar y^*\) to \(H\). The omitted coordinates have total at most \(C^*/(K_{\rm c}\sigma)<\tau/8\). Their removal adds at most that amount to the rounded residual and decreases the cost. The algorithm uses the same LP in both cases and does not compare \(C^*\) to \(m\sigma\). In particular, it can verify exact feasibility by checking \(y^0\)'s support, rounded box and residual before solving
\[
 \min\widetilde a^\top y,\qquad
 \bar A'_H y+r^+-r^-=\bar b,\qquad
 0\le y\le\bar u_H,\quad 0\le r^\pm\le2\tau
\]
at tolerance \(\eta\). The positive and negative parts of \(\bar b-\bar A'_H y^0\) extend \(y^0\) to feasible \(r^\pm\) in the specified boxes. The cost witness above also extends to these boxes and has rounded objective at most \(C^*/\sigma\). Pad the result by zeros. Its residual is at most \(2\tau+\eta\) and, since all \(y_e\le1\),
\[
 a^\top y/\sigma\le C^*/\sigma+\eta+mh_2\le C^*/\sigma+2\eta.
\]
Set \(y'=(1-\theta)y\) exactly. By (1), \(C^*>0\) and \(2\eta\le\theta C^*/\sigma\). Thus \(a^\top y'\le(1-\theta)(C^*+2\sigma\eta)\le(1-\theta^2)C^*<C^*\), using \(0<\theta<1\). Its true normalized residual is at most \(2\tau+\eta+m\theta+(1+3m)h<3\tau\le\delta/U\), because \(\|\bar A'(y'-y)\|_\infty\le m\theta\). It stays in the rounded box, which lies inside the true normalized box. The earlier zero-cost case uses \(y'=y\) and the same residual allowance. In either case set \(y'_e=0\) on omitted coordinates, including those with zero rounded width. Return \(Uy'_e\) for \(c_e\ge0\) and \(u_e-Uy'_e\) for \(c_e<0\). In particular an omitted negative-cost coordinate returns \(u_e\). The interval \(0\le y'_e\le u_e/U\) gives the exact capacity bounds in either sign case. Multiplying the normalized balance by \(U\) gives the original residual bound. The cost identity \(c^\top f=c^\top z+Ua^\top y'\) gives \(c^\top f\le c^\top z+UC^*=\mathrm{OPT}\). These records specify exact values even at a near cancellation at a bound. Expansion into a requested conventional encoding is charged separately.

Here \(\log(1/\tau)=O(B)\) by \(U\le2W\), and all normalized magnitudes, meshes, and tolerance words are \(O(B)\)-bit dyadics. The clipped requests and exact ordering avoid logarithms of inverse nonzero costs or of \(1/(1-\gamma_e)\). Their exponent comparisons need only a window \(O(B)\), even when a scale or cost exponent is referenced instead of written. Each additive call uses its own \(\widehat B\) from the assembled words. The analytic parameter \(B\) may be much larger than every \(\widehat B\). In §§4--7 all numerical bounds for a single assembled call hold with \(1+\widehat B\) in place of \(B\), since they use only the assembled LP and its tolerance. All computations here apart from the additive calls cost \(N\operatorname{poly}(B)\), including the explicit \(O(\log N)\) copy and verification scans, and each call has \(O(N)\) variables and rows. Output a single reference for \(U\) and one record per input arc, in input order. A discarded zero-capacity arc receives a zero record. Each remaining arc receives its dyadic \(y'_e\) and a sign indicator. A negative-cost record with \(y'_e=0\) represents \(u_e\). There are \(O(N)\) records of \(O(B)\) bits apart from the original input references. Later failure allowances are conditional on the call input, so adaptive binary search does not change the probability guarantee.

\hypertarget{a-simultaneous-selector-for-priced-two-entry-images}{%
\section{A simultaneous selector for priced two-entry images}\label{a-simultaneous-selector-for-priced-two-entry-images}}

\hypertarget{atoms-transfers-and-conformal-splitting}{%
\subsection{Atoms, transfers, and conformal splitting}\label{atoms-transfers-and-conformal-splitting}}

A signed traversal of a column with entries \((-a,b)\), \(a,b>0\), takes mass \(x\) from its first end, delivers \(xb/a\) at its second, and uses coefficient \(x/a\) at raw price \(\ell x/a\). The reverse uses coefficient \(-x/b\), gain \(a/b\), and price \(\ell x/b\). Along a path use the actual successive masses. A signed list's raw price counts absolute traversal prices before cancellation. A named resource is counted once per occurrence even for a small nonzero coefficient. The resource may be an input lifetime, an edge in a particular forest, or an endpoint incidence, and one atom can have multiple occurrences of one name. In the selector the input is a finite list of atoms \(U_f\) with actual compact lifts, positive raw certificates \(\ell_f\), and images that are zero, unary, or \(-a_fe_u+b_fe_v\) on distinct rows. Use \(r\) for the number of eligible image labels and \(M\) for a structural master size. Drop unused image labels from the selector's local data structures. The listed combinatorial input has size polynomial in \(M\). Each resource has input occurrence sum at most \(C_0\), including per-label incidences, and each atom contributes at most \(J\) occurrences to any name. Let \(\beta\) bound the explicit bit height of current rational prices and image coefficients. Prices and gains are unrestricted.

The following conformal operation will also be used on actual compact lists. Orient each nonzero signed constituent by its sign. Treat a unary as an edge to an unconstrained private leaf, replacing it by its negative when its sign requires that orientation. In the nonnegative cone preserving balance at every row outside the specified anchors, decompose the original nonnegative coefficients into extreme rays. A ray has connected support because distinct components can be scaled separately. If the support has \(h\) edges, \(v\) vertices and \(t\) unconstrained vertices, one-dimensional nullity at the constrained rows gives \(h\le v-t+1\). Connectivity gives \(h\ge v-1\). In particular \(t\le2\). At \(t=2\) the support is a tree. Its only leaves can be the unconstrained rows because a constrained leaf with positive flow violates balance. It is a directed path, giving opposite signs at its two endpoints. A private unary leaf drops out of the image. At \(t\le1\) the true image is zero or unary. In particular, with no private leaves, a nonzero unary ray is a directed cycle with an optional attached tail to its anchor, by the same dimension bound and the leaf test. Count resources in the original macro occurrences, with private leaves used only for the proof. The price sums of all the conformal pieces add exactly to the starting raw price.

For a list of \(p\) nonzero occurrences this decomposition uses at most \(p\) rays, counting each zero-image macro column as its own ray. One construction works on the current positive support and the slice where the coefficients sum to one. Find a direction in the homogeneous equality space with coordinate sum zero if one exists. Move until the first zero, repeating until an extreme point of the slice is found. Subtract its largest conformal multiple from the original nonnegative remainder, eliminating a positive coefficient, then repeat. Gaussian elimination on the local rational matrix suffices. Each produced direction has bit height polynomial in \(p\) times the input bit height by clearing denominators and using minors. The residual coefficient vector can use a common denominator, multiplied at each subtraction by the denominators of the new direction ratios. Thus the whole operation uses a fixed polynomial number of bit operations and its output height is linear in the starting bit height times a polynomial in \(p\). Each old occurrence is used at most once per output ray, so contributes at most \(p\) output visits to its resources. Existing path handles are treated as macro columns since they cancel at their internal vertices even if paths meet away from the macro endpoints.

\hypertarget{short-certificates-that-tolerate-both-gain-orientations}{%
\subsection{Short certificates that tolerate both gain orientations}\label{short-certificates-that-tolerate-both-gain-orientations}}

Assume \(r\ge2\) for binary demands, and put \(L=4\lceil\log_2(2r)\rceil+6\). For a set of binary names \(T\), let \(F(T)\) consist of demands having no path of at most \(L\) names in \(T\setminus\{f\}\) that supplies the demand's full signed amount at one of its endpoints for price at most \(L\ell_f\). The amount delivered to the other endpoint is whatever the path transports. Membership is decreasing under enlargement of \(T\).

We claim \(|F(T)\cap T|\le5r\). Otherwise peel vertices of degree at most five in the multigraph of those names, leaving a nonempty graph of minimum degree at least six. Fix a BFS tree in a component. For its balls \(B_i\) of sizes \(s_i\), every nontree edge within \(B_i\) gives a fundamental simple cycle of length at most \(2i+1\), including a two-edge parallel cycle. Set \(h=\lceil\log_2(2r)\rceil\). If all fundamental cycles in \(B_h\) were vertex-disjoint, each \(B_{i+1}\) would contain at most \(s_{i+1}/2\) nontree edges. All edges incident to \(B_i\) lie in \(B_{i+1}\), so then \(3s_i\le |E(B_{i+1})|\le 3s_{i+1}/2-1\). This gives \(s_h>2^h\), impossible. Two intersecting cycles give a connected bridgeless graph with at most \(4h+2\) edges and cycle rank two, since there are precisely two nontree edges with respect to the fixed tree. Its columns are dependent. In any nonzero dependence \((x_e)\), choose \(f\) to maximize \(|x_e|\ell_e\). Its exact replacement on the other edges has price at most \((L-1)\ell_f\).

Deleting \(f\) leaves a unicyclic graph. The replacement's support after removing zero coefficients therefore has at most one cycle. Conformally split outside its two prescribed endpoints as in §2.1. Any nonzero unary ray must use the sole cycle in its fixed directed orientation, and all such rays have the same sign at their anchors. Indeed a gaining cycle outputs excess at its entry and can only supply an outward tail, whereas a losing cycle can only be fed by an inward tail. A balanced cycle has no nonzero unary output. Choose the endpoint whose required sign is opposite this unary sign, or either endpoint if there is no unary. At this endpoint the sum of amounts of rays with the target sign is at least its target magnitude \(t\). These rays must be open paths to the other endpoint. Their total raw price is at most \((L-1)\ell_f\), so one has unit price at most \((L-1)\ell_f/t\). Scale it to \(t\), reversing and changing the global sign if needed when starting at the positive endpoint. This contradicts the criticality of \(f\) and proves the claim. In particular no gap from zero in a cycle residual is needed.

\hypertarget{sampling-and-exhaustion}{%
\subsection{Sampling and exhaustion}\label{sampling-and-exhaustion}}

Here is the sampling estimate for any decreasing \(F\) satisfying \(|T\cap F(T)|\le d\). A uniform sample of \(s\) distinct names with \(256d\le s\le m_0/4\), from \(m_0\) demands, has
\[
 \Pr(|F(T_s)|>m_0/2)\le 2e^{-s/64}.
\tag{2}
\]
In a uniform sample order let \(X_i\) indicate that the \(i\)-th name is in \(F(T_{i-1})\), with self-exclusion as above. In the reverse reveal, \(X_i=1\) exactly when the last exposed name of \(T_i\) is one of its at most \(d\) critical selected names. In the last \(b_s=\lceil s/2\rceil\) positions its conditional probability is at most \(2d/s\le1/64\). Iterating the conditional bound \(\mathbb E[2^{X_i}\mid\text{prior reveals}]\le1+1/64\) and using Markov's inequality gives probability at most \((1+1/64)^{b_s}2^{-b_s/8}\le e^{-s/64}\) that the block contains at least \(b_s/8\) failures. In the forward reveal, as long as \(|F(T_{i-1})|>m_0/2\), each conditional probability is at least \(1/4\), since at most \(s\le m_0/4\) names have been exposed. Define the indicator to be one after this condition first fails. For the modified indicators the corresponding conditional moment \(\mathbb E[2^{-X_i}\mid\text{prior reveals}]\) is at most \(7/8\). The probability that their sum in the block is at most \(b_s/8\) is at most \((7/8)^{b_s}2^{b_s/8}\le e^{-s/64}\), using \(\log 2<3/4\). On \(|F(T_s)|>m_0/2\) the modification never occurred, by monotonicity. On that event either the original block sum is at least \(b_s/8\) or the modified sum is at most \(b_s/8\). The two bounds give (2).

If \(S\) has uniform size \(t\), \(512d\le t\le m_0/4\), let \(\zeta=2^{-16}\). For a fixed set of at most \(\zeta t\) deleted positions in a uniform sample order the remainder is a uniform sample of size at least \(t/2\). There are at most \(\exp(h(\zeta)t)\) such sets, with binary-choice entropy in natural units \(h(\zeta)<1/256\). Thus (2) gives
\[
 \Pr(\exists X\subseteq S,\ |X|\le\zeta t,\ |F(S\setminus X)|>m_0/2)\le 2e^{-t/256}.
\tag{3}
\]
The set \(X\) is allowed to depend on all previous routes.

In a stage on \(m_0>0\) remaining binary demands take the common probability
\[
 p=\min\{1,\ \mathrm{roundup}_{2^{-\mathbb Z_{\ge0}}}\max(A r/m_0,\ cJ\log(2M)/C_0)\},
\]
where rounding above one means one. In forming the dyadic probability use the integer bit length of \(2M\) for its logarithmic factor. Use \(A\ge8192\) and large absolute \(c\). If \(p>1/16\), all current demands are affordable in the support budget as literals. Indeed \(m_0=O(r)\) if the first term forced the cleanup, and \(m_0=O(rJ\log(2M))\) if the second did, using \(2m_0\le r C_0\). Otherwise draw independent indicators \(I_e\) with this \(p\), using bounded fair bits. Verify \(m_0p/2\le t=\sum I_e\le2m_0p\) and for each resource
\[
 \sum_e I_e m_{xe}/p\le2C_0,
\tag{4}
\]
where \(m_{xe}\le J\) is its atom multiplicity. The expectation is at most \(C_0\). Each random summand is at most \(C_0/(c\log(2M))\), and its variance is at most that bound times its mean. Bernstein at excess \(C_0\), followed by a union bound over the polynomial number of named resources, proves a high constant probability for the simultaneous check. A binomial size bound and (3) apply because the sample is uniform conditional on its size before (4) is tested. The accepted size is \(O(r(1+J\log(2M)))\). Equations (3) and (4) need only a union bound.

Keep all sampled names in the stage support and give each \(P/p\) traversal tokens with \(P=\lceil4H/\zeta\rceil\), where \(H\) bounds the actually returned path length. Process demands in a fixed order. If the demand's own name was sampled and has a token left, assign the single-literal list with coefficient one. Otherwise use the gate in §2.4 to match either endpoint and take its aggregated list. Debit one token per name on either kind of list and remove exhausted names from later searches. Defer only on gate failure. At most \(Hm_0\) tokens are spent, so at most \(Hm_0p/P\le\zeta t\) names are exhausted. Let \(X\) be this set at the end of the stage. A deferred demand had no available literal and failed the gate. It lies in \(F(S\setminus X)\), since an \(L,L\) witness in the final available set would have been present at its query. On (3), at least half are served. Verify this progress and all path certificates, then use independent fresh stages on the deferred list. All assigned paths keep their names, including exhausted ones. There are at most \(O(\log M)\) successful stages. By (4) each stage spends at most \(2P C_0\) occurrences on a resource. This proves simultaneous primary load \(O(C_0 H\log(2M))\), and total selected support \(O(r(1+J\log(2M))\log(2M))\). Resource checks and progress verification ensure this bound on every accepted outcome. Repeat a failed attempt with a cap chosen in §7. If there are no binary demands, skip the sampling and query stages.

\hypertarget{price-respecting-search-by-an-undirected-state-graph}{%
\subsection{Price-respecting search by an undirected state graph}\label{price-respecting-search-by-an-undirected-state-graph}}

Fix the branch parameter \(a\) and let \(X_*=M^{12}\), \(\lambda_*=\lceil\log_2(2+X_*)\rceil\), and \(L_*=4\lceil\log_2(4M^2)\rceil+6\). Take \(e(a)=2^{C_2a^2}\) with a sufficiently large absolute integer \(C_2\), and use
\[
          H=4L_*\lambda_*^{e(a)} .
\]
The selector's input guard bounds the actual number of eligible row labels by \(2M^2\), so \(L=4\lceil\log_2(2r)\rceil+6\le L_*\) on each binary call. These parameters are declared before the grid is built. Require \(H^2\le M\), \(\beta\le M^5\), an initial graph of size at most \(X_*\), and the same \(X_*\) bound for each cover reference size from §9. Section 7.1 proves eventual adequacy of these guards uniformly in the regime \(\widehat B\le N\).

For a demand with price \(\ell_f\), measure mass in units of \(\ell_f\) and start with \(a_f/\ell_f\) or \(b_f/\ell_f\). Use the same grid at each row, with points \(2^i j/Q\) for \(Q\le j<2Q\), \(Q\) a power of two between \(10^5H\) and \(2\cdot10^5H\). Write \(q=1+1/Q\). Consecutive ratios lie between \(1+1/(2Q)\) and \(q\). Use a range of binary exponents \([-C H\beta,C H\beta]\), with absolute \(C\) and a few extra exponents for endpoint rounding. It contains every prefix of any \(L\)-hop witness, since every actual input ratio has logarithmic magnitude \(O(\beta)\). There are \(G_0=O(\beta H^2)\) grid points. At an available binary \((-a,b)\) on \(u,v\) with price \(\ell\), insert a unit undirected edge between \((u,x)\) and \((v,y)\) if
\[
 q^{-4}\le y/(xb/a)\le q^4,\qquad \max\{x\ell/a,y\ell/b\}\le4L .
\]
There are \(O(1)\) candidates \(y\) per \(x\), found by a clipped exponent and mantissa comparison. A path of price at most \(L\ell_f\) from the required endpoint gives a grid path by rounding its true normalized prefix amounts down. Its individual prices are at most \(L\) before rounding. Its price eligibility is symmetric under reversing an individual traversal.

Build the entire simple state graph of a sample that passes the size and resource checks. Each adjacency has a queue of the sampled names that generate it, and a name has an inverse list of its adjacencies. On exhaustion remove the name from its queues and delete the adjacencies whose queues become empty. Use the deterministic decremental cover in §9 at target radius \(L\). For a demand, take the rounded start state of one endpoint and find its covering cluster. Scan the grid states of the opposite endpoint for a member of that cluster. Repeat at the other orientation if needed. Upon the first positive test ask for one short path, which by the bound \(X_*\) has at most \(H\) edges. An \(L\)-hop price-respecting witness projects to a state path within the covered radius, so its destination passes the test. Every returned edge has a live name in its adjacency queue. The cover is deterministic for the online sequence, including deletions caused by earlier reported walks. \href{https://arxiv.org/pdf/2109.05621v1}{Chuzhoy v1, Theorems 1.2 and 3.3 and §§3.4.1--3.4.3}

Lift the reported walk using actual rational gains from the true mass. The ratio of an actual normalized mass to its state mass stays within \(q^{4H+1}\) in both directions by the edge tests, hence within two. Each raw price is at most \(8L\ell_f\). The exact signed list therefore costs at most \(8HL\ell_f\). Starting at the target's positive end requires a global sign change of the lift. Combine repeated names and drop canceled coefficients before storing the list and debiting tokens. Keep the unfolded traversal price in the certificate. Exact image and true price checks are on the actual list. The graph has \(n'=rG_0\) labels, including isolated states. One attempted stage initializes and may delete \(m'=O(|S|G_0)\) links. Each demand uses at most \(2G_0\) membership tests and at most one path query. The total cover work from §9, including all initial builds and updates in each initial connected component, is
\[
 O_a\big((n'+m')X_*^{c_1/a}\log^{O_a(1)}(2M)\big).
\]
For small initial components the separate BFS probes cost \(O_a(1)\) each. Each large-component membership probe uses \(O_a(\log^{O_a(1)}(2M))\) operations on cover records. A single positive probe per demand entails one path query and conversion with the same logarithmic-power bound. If the list size is \(b_{\rm list}\), then \(r,|S|\) and the number of attempted demands are \(O(b_{\rm list})\). With packed exponent and significand indices the enumeration and exact filtering of candidate adjacencies take \(O((r+|S|)G_0)\) tests on rational words of height polynomial in \((1+\beta)(1+H)\), with absolute degree. Lifting uses products and price sums of \(O(H)\) rational words, so its bit height also has that bound. Sum all membership and path costs over demands and add (A) over components, including array construction and disposal. A trial rejected in the preliminary size or resource scan requires only its own draw and scan. A trial that reaches the cover pays the entire first construction and every deletion, even if a progress or validation check later rejects it. The resulting per-attempt bound is
\[
 (1+\mathrm{list\ size})\,M^{O(1/a)}\log^{O_a(1)}(2M)\,\operatorname{poly}(\beta)
\tag{5}
\]
with absolute exponent of \(\beta\). In fact \(G_0=O((1+\beta)(H+1)^2)\), and the state graph factor \(X_*^{c_1/a}\) is bounded independently of \(\beta\). State IDs use \(O(\log M)\) bits. Each attempt that constructs a cover pays its first build. An accepted path emits only exact data formed from \(O(H)\) original columns. Its length, price factor and primary load bound are logarithmic powers in \(M\), independently of the precision. Neither state nor cover records enter a child chart.

\hypertarget{exact-unary-completion-and-projections}{%
\subsection{Exact unary completion and projections}\label{exact-unary-completion-and-projections}}

Let \(P_f\) be an assigned exact primary lift, taken as zero for a unary or zero image. The difference \(D_f=U_f-P_f\) has unary or zero image, with certificate \(L_f'=\ell_f+L(P_f)\). If the image is zero, use \(P_f\) as the replacement. The physical difference can still appear in a kernel list. At each row sort the nonzero unary images by decreasing exact efficiency \(|(AD_f)_v|/L_f'\). Divide ordered ranks into groups of sizes \(1,1,2,4,\ldots\), truncating the last group. The first rank uses itself. In each later group, independently choose a uniform pivot from the preceding group. For the recipient \(D_f\), use the exact signed multiple of that pivot matching its image. Its price is at most \(L_f'\). This includes arbitrarily small nonzero residuals and uses no tolerance comparison. The support acquires the chosen pivots' central original atoms, at most \(O(r\log(2M))\), in addition to the primary support. Each final replacement stores at most \(2H+1\) original-atom occurrences and has certificate at most \(\ell_f+2L(P_f)\).

For any resource the sum of its visits in all pivot candidates is at most its input load plus its primary load. The expected added visits in subsequent groups are at most twice that sum, since group sizes grow by at most two. The contribution from one pivot draw is bounded by \(C_0 J(1+H)\), since the anchor incidence count is at most \(C_0\). For independent nonnegative variables bounded by \(D_0\), the inequality \(e^{x/D_0}\le1+(e-1)x/D_0\) on \([0,D_0]\) gives an upper tail \(O(\sum\mathbb E X+D_0\log M)\) with failure \(M^{-c_1}\) for any fixed \(c_1\) by raising the constant. The rank-one part is deterministic and is covered by the candidate visit sum. Apply this to each named resource, verify all loads, and cap retries. Total load is \(C_0\log^{O_a(1)}M\) for \(J=\log^{O_a(1)}M\). This argument uses actual resource multiplicities, irrespective of whether pivot scaling is positive or negative. For a uniform rank in a group draw \(\lceil\log_2(\text{size})\rceil\) bits, accepting only valid ranks. Cap retries at \(C\log(2M)\) per draw and reject the entire batch on any exhaustion. The union bound over the polynomial number of groups keeps the per-batch success probability above \(1/2\).

For validation, aggregate visits first by original atom name, and then scan each atom's resource list once. In a unary group add the use of its pivot constituents multiplied by the number of recipients. Each constituent list is short in original atom names. This preserves (5) even when the resource lists have different lengths. Exact image and price checks can likewise use the two-entry images and the original certificates, before expanding any physical paths for a snapshot's own indices.

These steps prove the full-image selector with no gain or price range factor in support, stored list length, raw stretch, or named resource load. For a projected call first conformally split relative to all old and new anchors, then restrict each image to the new ones. It is still of the allowed type. Use the same selector on these images and their original lifts and price certificates. The full-image call uses the occurrence guarantee for its stored replacement lists. In particular a projected unary lift may have an image at an old row outside the projection. The returned original factors are the actual lifts in either call.

\hypertarget{frozen-common-marginals-and-the-physical-forests}{%
\section{Frozen common marginals and the physical forests}\label{frozen-common-marginals-and-the-physical-forests}}

A chart consists of oppositely signed binary columns, genuine unaries, and possible zero columns. It has a signed protected unit coordinate \(Z_v\) of price one and score zero at each row, in the chart's common row units. Prices are positive, immutable on each lifetime, and the incidence degree is \(D=\log^{O_a(1)}M\). At a nonbottom level keep each zero input column directly as a singleton kernel witness with its own price and score, inserting and erasing it on the column's lifetime edits. Remove these coordinates from the geometric construction and the compressed child. For a comparison vector the zero-column part is a sum of those singleton witnesses, and the remaining part has the same image and disjoint cost terms. An envelope \(\Xi\ge1\) bounds the ratio of every absolute entry to its column price. Let \(m_i\le M\) be this level's structural capacity and put \(k\asymp m_i^{1/a}\), a power of two. At every nonbottom level the guard \(b_i\ge2\) of §7.1 implies \(k\ge2\). In an epoch of at most \(m_i/k\) edits we keep fixed copies of the initial binary paths and use new logical marks for all edit endpoints. The chart may acquire new columns and rows during the epoch. Throughout §§3--4, a physical coordinate or resource means a lifetime or incidence of this chart. Its representation in earlier charts is handled separately by the public-copy bound of §5.2.

Write \(\mathcal L\) for a factor \(\log^{O_a(1)}(2M)\) whose exponent may increase by a fixed amount or by a factor depending on the fixed branch \(a\). Values called constant-degree polynomial in local size and precision have absolute degrees, independent of \(a\).

\hypertarget{a-common-gauge-short-lag-and-terminal-policies}{%
\subsection{A common gauge, short lag and terminal policies}\label{a-common-gauge-short-lag-and-terminal-policies}}

Normalize the price of each column to one by dividing by its price and multiply row \(v\) by \(r_v=\min_{f\ni v}\ell_f/|A_{vf}|\). The resulting magnitudes are at most one. For this geometry only, let \(U_v\in[1,Z]\) be the cheapest genuine unary price to remove one normalized unit. This includes the protected cap, so \(Z\le\Xi\) in these units. Starting at \(z^0_v=U_v\), perform synchronous sweeps
\[
 z^{j+1}_v=\min\left(U_v,\ \min_{e=vw}\frac{1+b z^j_w}{a}\right)
\]
with both signed orientations \((-a,b)\) of each binary in local tail-head order. Break ties by lifetime order. The sweep map is coordinatewise monotone and its first application gives \(z^1\le z^0=U\), so repeated applications keep the sequence decreasing. Each term stays at least one because \(U_v\ge1\) and \(0<a\le1\). Partition time into blocks of \(H_0=10(k+1)^2\) sweeps. Each row drops by a factor greater than \(8/7\) at most \(O(1+\log Z)\) times, because its values stay in \([1,Z]\). With \(r_0\) rows, a block having more than \(\lfloor r_0/k\rfloor\) such rows accounts for more than \(r_0/k\) of these drops. Thus there are \(O(k(1+\log Z))\) blocks before one with at most \(\lfloor r_0/k\rfloor\) exceptional rows. Test the block with exact comparisons of its first and last values, mark those rows, and use its first vector. Multiplying the normalized rows by it gives the frozen gauge. Throughout one sweep a stored rational is the result of a single affine map of one previous value or a unary choice. Its bit height grows by at most a constant times the normalized input bit height per sweep. The sweep count is polynomial in \(k\) and in the bit height of \(Z\), with absolute degrees.

The forward and reverse offers of a path of at most \(H_0\) hops between stable rows imply, with its true gauged unit price \(L_P\) and transfer \(\tau_P>0\),
\[
 7/8\le L_P+\tau_P,\qquad (7/8)\tau_P\le1+L_P .
\tag{6}
\]
Indeed if the saved index is \(j\) and the path has \(h\le H_0\) hops, following it and using \(z^j\) at its endpoint gives an upper bound on \(z^{j+h}\) at its start. At a stable start this latter value is at least \(z^{j+H_0}\ge(7/8)z^j\). Divide by the saved start value when converting to gauged amounts. The reverse path costs \(L_P/\tau_P\) per unit, giving the second bound. Put \(s_e=\max(p_e,q_e)\) for the two gauged magnitudes of a price-one binary. On stable-stable binaries with \(1/s_e\le1/16\), use length \(d_e\) obtained by rounding \(1/s_e\) upward. Use the largest power-of-two mesh at most the smallest \(1/s_e\) in this set, skipping the construction when the set is empty. Then \(1/s_e\le d_e\le2/s_e\), and the mesh height is polynomial in the bit heights just bounded. On a metric hop the smaller gauged magnitude is at least \((7/8)s_e-1\) by (6) for one hop. Thus the unit price in either direction is at most \(2d_e\). Apply (6) to each prefix of a stable metric walk. Writing \(L_j\) for its accumulated price starting at unit input, the next increment is at most \(2d_{j+1}(1+L_j)/(7/8)\). Consequently \(1+L_P\le\prod_e(1+16d_e/7)\le\exp((16/7)\sum_e d_e)\). For total metric length at most a sufficiently small absolute \(c_0\), this gives \(L_P=O(\sum d_e)\) and, by (6), all prefix transfers in \([1/2,2]\). The same applies to any contiguous subwalk meeting the hop and length bounds.

The action attaining the minimum in the first sweep of the saved block is a fixed policy on all rows. If it is a binary, \(z^{j+1}_v\le z^j_v\) gives \(p-q=a z_v^j-b z_w^j\ge1\) for its oriented gauged entries. Per unit input it costs at most \(1-\tau\), with \(\tau=q/p<1\). A selected unary has \(z^{j+1}_v=U_v\), so \(U_v=z^j_v\) by \(z_v^{j+1}\le z_v^j\le U_v\). Its price per unit in the frozen gauge is \(U_v/z^j_v=1\). For a binary traversal with entering mass \(x\), the price is at most \(x-x\tau\), so path prices sum to at most the loss of mass along the path. A directed cycle of gain \(\tau<1\) gives true unit absorption by taking the once-around vector divided by \(1-\tau\), at cost at most one. For a compact default, remove one arc of each policy cycle, put a size-\(k\) portal grid on the remaining forest, and terminate the policy at its first mark, an available unary, or an unmarked short cycle. A portal grid on any tree is built bottom-up by marking a vertex whenever its unmarked descendant count reaches \(k\), charging a disjoint set of at least \(k\) vertices at every mark. Removing all marks leaves components of size less than \(k\). Restoring one removed cycle arc within a policy component joins at most two such pieces. Policy paths to their first stop and the unmarked cycles have \(O(k)\) hops and constant handle count, representing the omitted cycle arc as a real literal. For an edge with both endpoints unmarked, any start using it belongs to its portal-free piece. For a terminal edge into a mark, starts are in the incident unmarked pieces, at most \(D\) pieces of size \(O(k)\). In particular the default costs at most one per unit and at most \(O(kD)\) direct starts use a given edge. \href{https://raw.githubusercontent.com/openai/math/adc7f1241b42e322a6451854ab7e4b4c146bf78a/preprints/Almost-Linear-Time-Maximum-Cardinality-Matching-in-Sparse-General-Graphs-September-24-2026/main.pdf\#page=58}{OpenAI, Lemma 7.6}

Distances in the stable metric graph are \(+\infty\) between components, and a distance label given by a minimum over no reachable offers is \(+\infty\). Write \(R_k(x)\) for the distance to the \(k\)-th stable vertex from \(x\), counting \(x\), or \(+\infty\) if the component has fewer than \(k\) vertices. Compute all \(k\)-nearest lists by multi-source Dijkstra. Insert all seeds with their initial offsets, shifting those offsets by a common constant if necessary to make them nonnegative. The heap records an origin, an arrived-at vertex, and its tentative distance, ordered by distance and then the fixed origin order. Expand a popped record only if its origin is new at its vertex and that vertex has retained fewer than \(k\) origins. Store the origin and predecessor, and enqueue its distance plus each incident length at the adjacent vertex. If an origin is among the first \(k\) at an endpoint of a shortest path, its rank cannot be worse at any prefix vertex, since better origins at a prefix extend along its remaining suffix with the same tie preference. Thus its true shortest label can propagate all along that path unless it has already settled with the same label. The nonnegative lengths and the heap order ensure that this is the label accepted for the origin. There are at most \(k\) expansions per vertex and at most \(O(m_i k)\) heap records in total, giving \(O(m_i k\log(2Mk))\) operations on distances. When computing shifted offers later, remove the common offset from the resulting labels. For the volume exit, in increasing \(R_k\) choose centers with disjoint finite \(k\)-nearest lists. A skipped \(x\) is within \(2R_k(x)\) of an earlier accepted center by a list intersection. There are \(O(m_i/k)\) centers. Build one shortest-path forest to these centers and stop at the first mark, adding its portal grid. Its offered distance is at most \(2R_k(x)\), and its realized path has at most \(k+1\) hops.

For an exceptional exit, at every stable neighbor of an exceptional row give the offer of taking that final literal hop, with offset the true unit-input price \(1/p\) at its stable end. A multi-source shortest-path forest with these offsets gives the label \(D_E(x)\). Each offset is a rational of local height and distance additions use only the common dyadic metric and one source offset. Truncate at the first portal, using a size-\(k\) grid and marking the exceptional targets. If a seed is reached, the final hop terminates at its marked exceptional row. It can have unrestricted delivered mass, while paying the true unit price. The unit prices of these excursions are \(O(D_E(x))\) whenever a gate below allows them, by (6) for the short stable prefix. A losing offer with no use can be omitted. The functions \(R_k,D_E\) are 1-Lipschitz on each finite component where they are finite, since they are an order statistic of distances and a minimum of shifted distances respectively.

\hypertarget{a-small-number-of-physical-copies-for-all-metric-scales}{%
\subsection{A small number of physical copies for all metric scales}\label{a-small-number-of-physical-copies-for-all-metric-scales}}

Here is a proof of the needed metric cover with arbitrary dyadic lengths. Let \(h=\lceil\log_2(2M)\rceil\). Use cutoffs \(t_j=2^{8hj}\) and a second sequence multiplied by \(2^{4h}\), for integer \(j\). A band \([t_j,t_{j+1})\) will serve distances in \([2^{2h}t_j,2^{-2h}t_{j+1})\). These serving intervals in the two systems cover every positive distance. For each system process occupied bands in increasing order. Before a band the union-find classes are the components of shorter edges. Contract them for that band, drop induced loops, and retain the exact input IDs on quotient edges joining distinct classes. Only incident quotient vertices need materialization. Across all bands the number of materialized incidences and quotient labels is \(O(m_i)\). Sort occupied exponents and find band indices by integer comparisons, skipping empty bands.

Round each quotient edge's length divided by \(t_j\) upward to an integer. This is at most twice its original scaled length and is at most \(2^{8h}+1\). For each radius \(R=2^s,\ 0\le s\le8h\), produce \(O(\log(2M))\) partitions by giving each quotient label \(c\) an independent geometric delay \(X_c\), with \(\Pr[X_c\ge t]=2^{-t}\) for integers \(t\ge0\). A vertex chooses the maximal \(R X_c-\operatorname{dist}(c,v)\), with a fixed tie order. Reject any delay above \(\Gamma_0=C\log(2M)\) rounded up. The winning center is within \(\Gamma_0 R\) because the vertex's own score is nonnegative. Center dominance persists along its shortest path, so a cell has an in-cell shortest-path tree. At every vertex check that in at least one partition its runner-up is behind by more than \(2R\), with the check automatic if there is only one reachable origin. This ensures its closed radius-\(R\) ball stays in the cell. For a proposed winner conditional on all other delays, the winner event is an integer lower threshold on its own delay. Exceeding this threshold by three forces margin greater than \(2R\), at conditional probability at least \(1/8\) of the winner event. Summing over possible winners proves a single-partition probability at least \(1/8\) before capping. Independent partitions and a union bound on delays and margins over all band vertices and radii show all checks pass with probability at least \(3/4\) for absolute constants. Generate capped delays with fair bits. Find each winner and runner-up by the same truncated Dijkstra proof above with two origins and initial offsets \(-RX_c\) and positive edge lengths. This reproduces the verified padding argument of the matching preprint using only \(O(\log M)\)-bit quotient integers. The total work per attempt over occupied bands, radii and partitions is \(m_i\log^{O(1)}M\), apart from the original rational distance bookkeeping. \href{https://raw.githubusercontent.com/openai/math/adc7f1241b42e322a6451854ab7e4b4c146bf78a/preprints/Almost-Linear-Time-Maximum-Cardinality-Matching-in-Sparse-General-Graphs-September-24-2026/main.pdf\#page=54}{OpenAI, Lemma 7.2 and its proof}

For each system and each pair \((s,\text{partition index})\), extend the cell trees of each band to a spanning forest of its quotient components. Keep their original edges in a single physical forest for this index, as the bands are processed. Inductively the forest spans each old class using edges of length below \(t_j\) and contracted joins are acyclic. A path using one quotient cell tree visits every old class at most once. Its within-class detours have total length at most \(Mt_j\). For endpoints of original distance \(d\) served by this band, a shortest path uses no edge \(\ge t_{j+1}\), and its rounded quotient length is at most \(2d/t_j\). If the classes coincide the within-class bound alone suffices. Otherwise they occur in the quotient, and the least dyadic \(R\ge2d/t_j\) is in the tested range. Use a partition padding the first endpoint. The physical path length is at most
\[
 2\Gamma_0 R t_j+Mt_j\le (8\Gamma_0+1)d .
\]
Preprocess each assembled forest in one root traversal, storing component IDs, depths, exact dyadic root distances, and ancestors at distances \(1,2,4,\ldots\) in hops. Obtain the ancestor at distance \(2^j\) by applying the previous jump twice. For a same-component pair, use the jumps to bring the deeper vertex to the shallower depth, then lift both in descending powers, keeping them below their common ancestor until their parents agree. If the vertices coincide after the first step, that vertex is their common ancestor. The path length is the sum of the two root distances minus twice the common ancestor's root distance. For each ordered pair of interest scan the indices once and choose the first with path length at most \(\Lambda\operatorname{dist}(x,y)\), taking \(\Lambda=8\Gamma_0+1\) and skipping distinct-component pairs in a forest. For \(x=y\) use the empty path. A serving band skipped for lack of edges can only have the same-class case for the query, for which the earlier forest edges suffice. We use only the \(O(m_i k)\) ordered pairs and their distances from the \(k\)-nearest lists. The preprocessing and scans take \(m_i k\log^{O(1)}(2M)\) word operations. The resulting fixed path \(P(x,y)\) uses one handle in one of \(O(\log^2(2M))\) physical forest copies. A path of metric length strictly below \(R_k(x)\) has fewer than \(k\) hops because each of its distinct vertices is that close to \(x\). Mark a portal grid in each assembled forest for later compression. The cutoffs, original path summaries, and exact threshold comparisons have polynomial bit work in the available local rational height. None of the forests replicates as the exponent span grows.

\hypertarget{shared-genuine-absorbers}{%
\subsection{Shared genuine absorbers}\label{shared-genuine-absorbers}}

To certify short priced absorption, try one LP per stable row \(x\), minimizing raw price for image \(-e_x\) on the stable ball \(B(x,R_k(x)/4)\), or its full metric component if it has fewer than \(k\) rows. Use all chart columns internal to the ball, including its genuine unaries. Separate the two signs of each variable and impose raw cost at most the genuine cap certificate at \(x\) to bound the polytope. The displayed unit equation is in the frozen gauge. At an optimum a dependence on the positive signed coefficients has zero price derivative: otherwise a small step of the decreasing sign stays nonnegative, preserves the image, and stays in the price sublevel while decreasing the price. A zero-derivative dependence gives a short nontrivial feasible segment on both sides of the optimum by positivity of the supported coefficients. Therefore the optimum vertex selected in §7.4 has independent supported signed columns. Every other connected component has zero requirement, and removing its supported columns would reduce the price unless it is already empty. A pure tree of binary columns cannot source a unary, since its positive annihilating potential is obtained by propagating the gain ratios. The relevant component has at least \(v-1\) binary edges and at most \(v\) total columns on \(v\) rows. Conservation permits leaves only at \(x\) or a private unary leaf of §2.1, and degree-two balanced vertices have one incoming and one outgoing traversal. The support is a path to a true unary or a directed cycle with an optional tail from \(x\). In the latter case \(-e_x\) forces cycle loss and the once-around vector is scaled by the entering mass divided by \(1-\tau\). Keep exact rational coefficients even if \(1-\tau\) is tiny. These are explicit bounded feasible LPs on \(O(Dk)\) columns with rational height polynomial in \(k\) and the current input height. Exact rational optimization with a vertex optimum is used as specified in §7.4.

For sufficiently small absolute price \(h_0>0\), an absorber of price \(h\le h_0\) from this LP is contained in \(B(x,C_1h)\), for an absolute \(C_1\). Its tail costs at most \(h\), so (6) keeps every tail prefix's mass at least \(7/8-h\). The once-around unit cycle cost is at most \(h/(7/8-h)\) since entering mass is at least that much and \(1-\tau\le1\). Apply (6) to its prefixes as well. Every hop's price per unit is at least \(1/s_e\), which excludes nonmetric edges for small \(h_0\), and the uniform lower bounds on mass show that the sum of rounded lengths is \(O(h)\). This proof uses the \(H_0\) lag on all stable walks of the requisite hop count, without assuming that an absorber already uses metric edges.

Let \(H_x\) be the optimal offer. Keep candidates \(H_x\le\min(h_0,R_k(x)/C_2)\) for a sufficiently large absolute \(C_2\). In increasing \(H_x\) accept each unless an earlier accepted candidate lies within \(2C_1 H_x\). Accepted support balls \(B(x,C_1 H_x)\) are pairwise disjoint, and every skipped candidate is within \(2C_1 H_x\) of one of no greater price. This test uses the nearest lists. Set \(D_A(x)=\min_v(\operatorname{dist}(x,v)+H_v)\) over accepted owners \(v\). A deterministic shortest-path forest with these offsets realizes the distance part and each owner wins at itself. At a candidate row \(D_A(x)\le(2C_1+1)H_x\). Attach at the owner a private copy of its support, keeping one closing chord of a cycle as a literal. The supports are disjoint, so the number of physical copies is bounded by a constant on top of the routing forest. For a route from an eligible state \(x\) of \(D_A(x)\le T r\) and \(R_k(x)>K r\), with \(K\) much larger than \(T\), the forest path to its owner has fewer than \(k\) hops. When \(Tr\) is small it can be traversed at price \(O(\operatorname{dist}(x,v))\) with transfer at most two. The complete hypothetical genuine absorber costs \(O(D_A(x))\) and is supported within \(O(D_A(x))\) metric distance from \(x\). The actual route may truncate its tail at the first portal. Give the forest with open cycles a portal grid and keep the cycle as an open path and a literal when used. If a grid mark lies on a cycle that a route would meet, truncate on first reaching such a mark. This still costs at most the complete loss-paid cycle price, since the full loss normalizer is at least one. Each path tail or once-around piece has \(O(k)\) expanded hops and a constant number of handles. No whole retained absorber is a changing primitive of the child.

\hypertarget{route-law-price-and-exact-finite-catalogue}{%
\subsection{Route law, price, and exact finite catalogue}\label{route-law-price-and-exact-finite-catalogue}}

Take \(T=C_3(1+\Lambda)\), \(K=C_4T\) with sufficiently large absolute constants, take the top \(r_\star\) dyadic and comparable to \(c/(K+T+\Lambda)\) for small absolute \(c>0\), and take a power of two \(G\) at least \(64\) and a sufficiently large fixed power of \(K+T+\Lambda\). All these and \(1/r_\star\) are logarithmic powers. Let \(S\) include the exceptional and initial portal and volume marks plus the endpoints of all edits to this point in the epoch. Initially \(O((m_i/k)\log^{O(1)}M)\) such rows suffice. Distinguish these logical marks from later Steiner ports that are used only to represent selected paths. For the logical distance \(D_S\), store a mark's offer only at itself and at rows strictly nearer than \(R_k\) of that mark. If \(D_S(x)\le8r\) could pass while \(R_k(x)>K r\), the mark has \(R_k\ge R_k(x)-8r>8r\) and its offer is present. Thus the stored minimum is the true eligible one, and the true nearest-mark distance is 1-Lipschitz. Marks absent from the frozen chart have no such offers.

Use a common shift \(\omega\in[1,G)\) uniform and shared priorities on all stable rows. To choose a common first scale, take the minimum \(t_{\min}\) of \(r_\star\), the least metric edge length divided by eight if there is an edge, and the positive finite values of \(R_k(x)/K,D_E(x)/8,D_A(x)/T\). The least edge length also bounds every positive ball-growth distance from below and, after division by eight, every positive \(D_S(x)/8\), including after future mark edits. If there are no positive finite labels or metric edges, use \(t_{\min}=r_\star\). Let \(j_0\) be the greatest integer with \(G^{j_0+2}\le t_{\min}\). This is an exponent calculation using the floor binary exponent of the rational \(t_{\min}\) and the integer \(\log_2 G\). The first scale gives singleton balls and precedes all positive gate thresholds and the top, for every shift. Ordinary stages have \(r=\omega G^j<r_\star\) for successive integers \(j\ge j_0\), followed by one dummy top. At each ordinary stage the current state \(x\) of an active start \(v\) is the preceding priority winner in the ball about \(v\), or \(v\) initially. It lies within \(r/G\) of \(v\). Test, in fixed order, a mark stop at distance zero, volume if \(R_k(x)\le K r\), exceptional or mark distance if \(D_E(x)\le8r\) or \(D_S(x)\le8r\), and absorption if \(D_A(x)\le T r\). Use their actual fixed routes. Portal truncations of the default, volume, exceptional and absorber routes use the epoch's fixed stops. New logical marks act through the mark gates. Every logical original start uses the zero route, and exceptional starts always use the zero route to themselves. If no gate fires, move to the priority winner of \(B(v,r)\) using the canonical ordered path from §3.2. At the dummy top take the truncated stationary policy. When a stable prefix exists before an exceptional last hop, that hop enters the marked exceptional row and ends the route.

A moving state has \(R_k(x)>K r\). The chosen detour length is at most \(2\Lambda r<R_k(x)\), and the representative's distance to \(v\) is still at most \(r\). Thus moves use fewer than \(k\) hops and all updated winners before stopping lie in the \(k\)-nearest list of their start. There are at most \(k\) changes of winner. A mark-distance path also lies below \(R_k(x)\) by the choice of \(K\), its cost \(O(\Lambda r)\). The volume and exceptional paths have been truncated on the fixed forests, and cost \(O(Kr)\) and \(O(r)\) per unit including a final exceptional hop. Absorber routes use the bound in §3.3. The ordinary stable part has at most \(k(k-1)\) nonterminal hops plus at most \(3k+3\) stable terminal hops. In particular (6) applies to the concatenation and each of its prefixes. Its geometric length sum is \(O((K+T+\Lambda)r_\star)\), chosen small for (6). A final exceptional hop, a loss-normalized cycle, or the stationary policy is charged by its own unit-input certificate. The raw price of the normal prefix through scale \(r\) is \(O(\Lambda r)\), and all complete unit routes cost \(O(1)\) including the top.

\textbf{Finite priorities.} For each start test the distinct closed balls at radii strictly below its \(R_k\). They are complete distance-tie prefixes of its stored list and include the singleton start. A tie at finite radius \(R_k\) is outside this test and can have more than \(k\) vertices. For each frozen metric edge test all unions of one such prefix from each endpoint. The pre-horizon argument uses only these unions. Their count is \(O(m_i(k+1)^2)\) and each has size at most \(2k\). Sample \(S_p=\lceil C k\log(2M)\rceil\) independent uniform priority permutations. Verify that every full stored list in distance order with the fixed tie order has at most \(C'\log(2M)\) successive record minima per permutation, and that for every tested union \(U\) and \(x\in U\) the frequency of \(x\) winning \(U\) is at most \(4/|U|\). The first check succeeds with high constant probability by independent uniform insertion ranks. For a list of length \(j\), \(\mathbb E 2^{\#\mathrm{records}}=\prod_{t=1}^j(1+1/t)=j+1\). A change of winner when extending a complete ball prefix must include a new record in this refined order. The second check follows by binomial tails with mean \(S_p/|U|\) and a union bound. Exact scans of all tested sets take \(m_i k^{O(1)}\log^{O(1)}M\) work. Generate each permutation by sorting sufficiently many random key bits and reject ties. Caps give finite work. Use the empirical law, uniform over the verified table, independently of \(\omega\). The two checks hold deterministically for the accepted table. Its entire size is internal.

\textbf{Price of a pair of routes.} In price-one gauged column notation let a stable metric chord be \((-p,q)\) and \(s=\max(p,q)\). Its rounded length \(d\) satisfies \(sd\le2\). Couple its two arms by the same index and shift. For comparison define the horizon as the first radius where at least one original start has \(R_k(v)\le r\), or the dummy top, whichever is earlier. Before it both hypothetical centers, even on stopped arms, are determined by tested small prefixes. At an ordinary horizon an active arm at a start with \(R_k(v)\le r\) must pass the volume gate if no nearer mark stop intervenes. At \(r\ge d\) before the horizon, the centers differ if the union winner is outside the intersection. The tested relative frequency bounds this probability by
\[
 4(\log |B(u,r+d)|-\log |B(u,r-d)|).
\]
The proof brackets intersection and union by these two balls and uses \(1-x\le-\log x\). For \(r<d\) use the bound one. The product of radius and shift density is at most \(G/(G-1)\) per band. Summing and integrating the difference of translates of the monotone log count gives at most \(2d\log M\), since every jump contributes on an interval of length \(2d\). Thus the expected sum of disagreeing pre-horizon radii is \(O(d\log(2M))\). Let \(r_m\) be the stage just after the first agreeing update beyond the last such disagreement, or the horizon if sooner. It is at most \(G^2\) times the last disagreement radius, up to \(O(G^2d)\) for the start or top. This gives \(\mathbb E r_m=O(G^2d\log(2M))\).

Here are the deterministic price deductions from that meeting scale.

\begin{itemize}
\tightlist
\item
  A mark, exceptional, or volume gate on one arm at radius \(r\le r_m\) forces a stop of the other by radius \(O(G\max\{r,d\})\), possibly at the top. Write \(Ar\) for its gate threshold, with \(A=8\) or \(A=K\). At the other arm's stage of radius \(r'\), the centers are within \(d+r/G+r'/G\). The label is 1-Lipschitz, so if this stage is reached at the first \(r'\ge 2\max\{r,d\}\), its label is at most \(Ar+d+r/G+r'/G\le Ar'\). The stored mark offer is valid whenever the volume gate fails, by the offer rule above. Any earlier gate has already stopped the arm. With a logical-start stop, the first label is zero and the same bound uses \(r=0\). If the top intervenes first, then \(r_\star\le 2G\max\{r,d\}\); its constant unit price is \(O(r_\star^{-1}r_\star)\). An ordinary horizon with no prior meeting is covered by the volume gate argument. At the dummy horizon both default prices are covered at \(r_m=r_\star\).
\item
  If one arm absorbs before a common active state, use its hypothetical complete genuine absorber, even if its realized leg terminated at a portal. The earlier stop \(r_0\le r_m/G\) is within \(O((r_m+r_0)/G)\) metric distance of the meeting row using that start's hypothetical center. The canonical path there costs \(O(\Lambda r_m/G)\) in units with transfer at most two, assuming the meeting volume gate fails. The transported offer costs \(O(T r_0+\Lambda r_m/G)\) and lies in the small stable ball at the meeting row. With the chosen separation of \(G,K,T\) this is an admissible candidate below \(T r_m/(2C_1+1)\), so \(D_A\le T r_m\) and the other arm terminates there unless it already stopped. Its support is inside the candidate LP ball, since its radius is below \(R_k/4\), and it is below both price cutoffs by \(K\gg T\) and the small top.
\item
  If both arms enter a common state \(x\), their chord and signed prefixes form a raw vector \(V\) of price \(L_V\le 1+O(\Lambda s r_m)\) with physical unary image \(\rho e_x\) and \(|\rho|=O(s)\). Identical later transitions and terminations use the same actual copies. The shared part is traversed with the actual combined coefficient whose magnitude is \(|\rho|\). A zero coefficient cancels the entire raw suffix. Otherwise \(-V/\rho\) is a genuine offer of price \(L_V/|\rho|\) at \(x\), supported within \(O(\Lambda r_m/G+d)\) of \(x\). At a meeting strictly before the horizon both starts are within the preceding radius of the common state, so the chord and all its prefix supports lie within \(O(\Lambda r_m/G)\) of it. At a later common stage \(r\) before any stop, pass to this offer by a canonical path of length \(O(\Lambda(r+r_m)/G)\). Its transfer is bounded, and the resulting offer lies in \(B(\text{state},R_k/4)\) if the volume gate fails. At the first eligible ordinary radius that is a sufficiently large constant times \(\max(r_m,L_V/|\rho|)\), its price fits the candidate cutoff and implies \(D_A\le T r\). If the horizon intervenes both arms are still at the same current state before the horizon gate and take the same action there. If the top or an earlier geometric stop intervenes, its price is bounded at its own radius with the factor \(O((1+r_\star^{-1})\mathcal L)\). Thus the common suffix bill is at most \(\mathcal L(L_V+|\rho|r_m)\).
\end{itemize}

The transported offer has a precise support and sign. If the current row is \(x'\), the old offer \(Y\) has image \(-e_x\), and the transport path \(Q\) has image \(-e_{x'}+\tau_Q e_x\), use \(Q+\tau_Q Y\). In the common-prefix case \(Y=-V/\rho\). Its image is \(-e_{x'}\), its price is at most \(L_Q+\tau_Q L(Y)\), and its physical support is the union of the path and offer supports. All rows of the old prefix offer are within \(O(\Lambda r_m/G)\) of the meeting state, including the two endpoints of the chord. The old absorbed-arm offer has support within \(O(T r_0)\) of that arm's old state. Hence their transported supports have radius \(O((\Lambda(r+r_m)+T r_m)/G)\) about the current state at an eligible stage \(r\ge r_m\). If its volume gate fails, \(R_k(x')>Kr\). This puts the support inside \(R_k(x')/4\). Choosing \(r\) large enough compared to \(L_V/|\rho|\) puts the price below the cutoff \(R_k(x')/C_2\) and below \(Tr/(2C_1+1)\). The top cutoff keeps this price below \(h_0\). For a truncated cycle in an actual route use its ordinary entering mass up to the first fixed portal. Use the loss normalizer only for a complete turn.

At a first ordinary horizon \(r_h\) with \(R_k(v)\le r_h\), a still-active other arm passes the volume threshold by the first available radius \(r\ge\max(r_h,d)\), since its state then has \(R_k\le r_h+d+r/G\). This also covers an arm that had already stopped before the horizon. An earlier top is charged by its fixed unit price. These implications use (6) only on stable paths. The constants can be chosen in order \(C_1,C_2,C_3,C_4\), then \(c\) sufficiently small and \(G\) sufficiently large, since every inherited radius term in a transported offer is multiplied by \(1/G\). The combined expected bill is \(\mathcal L\) in the price-one units, using \(sd\le2\) and the bound on \(\mathbb E r_m\). A stable chord outside the metric has \(s=O(1)\), and pays the two unit route bounds. For a chord with one exceptional endpoint, let \(p\) be its gauged magnitude at the stable end \(v\). Its direct seed gives \(D_E(v)\le1/p\). If an ordinary radius reaches \(1/p\), the first such radius is at most \(G/p\), and at its current state \(D_E\le1/p+r/G\le8r\). The arm stops there or earlier, at total unit price \(O((K+T+\Lambda)G/p)\). Otherwise its last ordinary radius is below \(1/p\) but at least \(r_\star/G\), so multiplying the constant unit price at the top by \(p\) costs \(O(G/r_\star)\). In either case the price contribution is \(\mathcal L\). The exceptional start itself has zero route. For two exceptional ends both routes are zero. A genuine price-one unary has gauged magnitude at most one because \(z_v\le U_v\). Zero-image columns have no arms. This proves the price bound for every frozen column.

\textbf{Finite realization.} A start's threshold calendar includes the distinct nearest-list growth distances, the gate thresholds at states in its nearest list, including each current mark offer minimum divided by eight, and the top cutoff. There are \(O(k)\) such thresholds. Omit infinite thresholds and treat zero thresholds as satisfied wherever their gate is examined. Each positive finite \(t\) supplies at most one breakpoint \(t/G^j\in[1,G)\). For a positive rational \(b\), the integer \(\lfloor\log_G b\rfloor\) is the floor of its floor binary exponent divided by \(\log_2 G\). Find the binary exponent from numerator and denominator bit lengths and a comparison with the candidate power of two. Apply this to \(b=t\) for the breakpoint. Use the union of the two chord-arm breakpoint sets and the endpoints \(1,G\), and evaluate each open interval at its rational midpoint \(\omega\).

For a fixed midpoint and threshold \(t>0\), compute \(j=\lfloor\log_G(t/\omega)\rfloor\) by the same comparisons. The first stage satisfying the threshold has index \(\max(j_0,j)\) if \(t/\omega=G^j\), and \(\max(j_0,j+1)\) otherwise. Include that index in the calendar and include its successor for each ball-growth event. Group events at an equal index. Starting at \(j_0\), visit their sorted distinct indices only until the top threshold and then take the dummy top if the arm is still active. Gate predicates at a fixed state change only at their threshold indices. Between threshold crossings and the listed successors, neither a ball nor its winner changes, and every gate at the current state has already been checked. This reproduces the successive-stage route. A ball larger than the stored list is never queried, since at any stage with \(R_k(v)\le r\) the active old state has \(R_k(x)\le r+r/G\) and stops by the volume gate or an earlier mark stop. Breakpoint values of \(\omega\) have zero measure.

For every priority index and interval keep the chord and the two exact prefixes with a shared suffix combined in its actual raw coefficient, or the two separate routes. Every atom uses \(O(\log M)\) handles by the record check. Its probability is the interval's exact length divided by \(S_p(G-1)\). Products and gains only refer to physical paths of \(O(k)\) hops per handle, including genuine cycle recipes. Their exact rational work and calendar arithmetic are polynomial in local \(k,D\), precision and logarithms with absolute degrees. The compact catalogue has \(O(k^2\log(2M))\) atoms per frozen column and fixed polylogarithmic many physical forest copies.

Fix one raw-copy literal \(\widetilde{1}_f\) for every frozen column, and use a linear projection \(\pi\) summing copy coefficients onto each physical coordinate. The unit route from each row \(v\) is a negative unit at \(v\) plus its terminal image, or a true absorbing unit. In physical units let \(P^S_v\) be its mean raw vector, converting out of the frozen gauge. Its law depends only on that start, the common table and shift, and the logical marks. The stop on a marked start makes \(P^S_v=0\) for \(v\in S\). For the coupled atom \(U_f\) the signed arms are exactly \(P^S A_f\) in expectation, so
\[
   \mathbb E U_f=\widetilde{1}_f+P^S A_f,
   \qquad
   \mathbb E L(U_f)\le\mathcal L\ell_f .
\tag{7}
\]
Every individual image is supported on \(S\) and every intermediate imbalance on a raw forest handle is at its two ends. The raw suffix cancellation in (7) occurs before taking \(\pi\) or clipping. This is why a priority table and forest tie rule are shared across all chords.

\textbf{Resource locality.} If a normal, nearby-mark, or complete cheap-absorber support visited at radius \(r\) reaches a row \(w\), its distance from the current state is at most \(O((\Lambda+T)r)\). Failure of the volume gate means \(R_k(x)>Kr\). Our choice \(K\gg\Lambda+T\) yields \(\operatorname{dist}(v,w)<R_k(w)\) for the original start \(v\), by the 1-Lipschitz bound and \(\operatorname{dist}(v,x)\le r/G\). Hence at most \(k\) original starts can ever use a fixed such row or incident resource. Current states reached before an exit have the same bound by the preceding normal move, or are their own original start. Build reverse state lists by scanning \((v,x)\) in the stored \(k\)-nearest lists of \(v\) and retaining only \(\operatorname{dist}(v,x)<R_k(x)\), together with each initial state. This lists every possible actual start of a state, even under future gate changes, and keeps at most \(k\) per state. Portal-truncated stationary, net and exceptional exits add \(O(kD)\) eligible exit states per resource. A terminal literal can be incident to a marked vertex with at most \(D\) adjacent pieces of size \(O(k)\). Thus one physical resource belongs to at most \(k^{O(1)}\mathcal L\) affected chord groups. Each group has \(k^{O(1)}\mathcal L\) atom occurrences, each containing only logarithmic-power many handles. One can store actual inverse resource lists by scanning their physical occurrences, in \(m_i k^{O(1)}\mathcal L\operatorname{poly}(B)\) work. A new logical mark affects at most \(k\) state offers and each such state at most \(k\) starts. Their incident frozen groups, and groups found by hole indices, are the only ones that need regeneration. These locality counts and the initial portal count depend on \(m_i,k,D,\log M\) and remain valid under adaptive choices of the marks.

\hypertarget{exact-comparison-through-births-deaths-and-resets}{%
\section{Exact comparison through births, deaths, and resets}\label{exact-comparison-through-births-deaths-and-resets}}

\hypertarget{comparisons-on-old-logical-marks}{%
\subsection{Comparisons on old logical marks}\label{comparisons-on-old-logical-marks}}

In a settled snapshot let \(S\) be the logical marks, and \(D_H\) the diagonal restriction killing every raw copy of each physically dead lifetime in \(H\). Keep the following comparison for \textbf{every} physical live \(z\) with \(\operatorname{supp}Az\subseteq S\):
\[
 z=\sum_f c_f C_f+\pi W,\qquad
 \sum_f |c_f|\,L(C_f)+L(W)\le K_S\|\ell z\|_1 .
\tag{8}
\]
Here \(C_f\) are installed exactly conserved signed lists, with exact scores and price certificates. The aggregate raw remainder \(W\) uses the current forest copies and selected or born literals, with divergence in each forest only at promoted copies of rows. The comparison is an existence assertion about \(z\). The data structure maintains the list and score heap, support, and ports. This formulation comes before identifying or connecting the copy rows.

A clipped path is split at its dead edges. At most its two live end portions can have original switch endpoints. Put the other surviving fragments into a forest remainder with endpoints promoted at holes in all affected copies. Do this with the signed coefficients actually transported in the original path, obtained by prefix products. Remove dead literals. A list of \(p\) handles thus has a direct part of at most \(2p\) handles. The price of its end portions and all interior fragments sums to its surviving raw price. The interior fragments in a conceptual remainder need not be stored per demand. Affected handles can be split by scanning their \(O(k)\) physical edges, or by a forest hole index. This rule works through any number of cuts, since each retained subpath of an original eligible handle has at most the original number of hops.

At initialization and at each full support reset, clip the current catalogue (7), move the interior fragments aside, and conformally split each direct lift into opposing binaries, true unaries and zero images on \textbf{all} logical marks. Run §2 on those actual signed atoms and their true certificates, including the original lifts and resource lists. For a direct lift \(X\) with returned replacement \(R\) store \(X-R\) as one physical kernel list, using all the returned signed factors of its conformal pieces. Zero images can supply a kernel list themselves. This installs their exact scores from frozen lifetime scores. Add \(R\) and the interior fragments to \(W\), along with the literal born part of the current input. If \(z\) balances off \(S\), so does its surviving frozen restriction because the born part has image on \(S\). Extend the frozen restriction by zero on dead coordinates. Its weighted sum of identities (7) then has \(P^S Az_{\rm frozen}=0\) before clipping. Applying \(D_H\) and projecting gives exactly its live vector. This proves (8) with \(K_S=\mathcal L\) by the selector and the price sum in (7). Each chosen \textbf{original} piece promotes the endpoints of its actual path or literal factors in their copies. Thus every used raw forest path in \(W\) has promoted ends. The literal births are supported on hubs. The marginal identity cancels in a common raw copy space. A stored kernel difference balances after physical identification.

\hypertarget{an-incremental-snapshot}{%
\subsection{An incremental snapshot}\label{an-incremental-snapshot}}

Take the earlier snapshot as base, with logical marks \(S_0\), and let \(H\) be the deaths and \(S_1\) the new mark set at the end of an interval of \(E\) edits. Put \(T=S_1\setminus S_0\). The cut anchors and \(T\) each have \(O(E)\) labels. Process these two calls.

\begin{enumerate}
\def\labelenumi{\arabic{enumi}.}
\tightlist
\item
  \textbf{Full cut repair.} Apply \(D_H\) to the inherited (8) for a comparison vector supported on survivors and balanced off \(S_0\). For a hit \(C_f\), let \(Y_f\) consist of the interior fragments just described and \(X_f=D_H C_f-Y_f\). Its old physical image was zero. The full image of its direct part is supported on cut anchors because the removed traversals and \(Y_f\) contribute only at these anchors. Split \(X_f\) conformally there, and replace its pieces using §2 on the full image. Install \(X_f-R_f\) with its exact score, even when the central demand has died and the surviving difference can still be useful. The new remainder receives \(Y_f+R_f\). An unhit \(C_f\) keeps its actual list, coefficient data and score. The old \(W\) clips at old ports and cut ports. Every previously selected live factor can remain until the window ends. This gives an exact inherited comparison on the surviving physical rows at an extra factor \(\mathcal L\).
\item
  \textbf{New-row repair.} Update the common \(S_1\) catalogue on all affected frozen chord groups. Find them by the reverse state lists and death visits of §3.4. Always use surviving frozen versions as the central chords and freeze the rest of the geometric paths before clipping. Split every clipped new atom as \(Y+X\) by keeping only end fragments in \(X\). Its full image is on \(S_1\). Conformally split at \(S_1\), and project each piece's image to \(T\). Apply the \textbf{same} selector on \(T\) with the full lift certificates. Let \(R\) denote the sum of the returned lifts for \(X\). Then \(X-R\) has image on \(S_0\) and uses only surviving parent lifetimes. The residual can have nonzero image there, so its role is in the inherited comparison, rather than in the kernel heap. Promote the actual factors of \(Y+R\) as needed, with \(Y\) already having cut endpoints. Current interval births stay literal on their new prices and scores.
\end{enumerate}

To verify the second step for general \(z\) balanced off \(S_1\), apply (7) with \(P^{S_1}\). Let \(F\) be the surviving frozen part of \(z\) and \(H_{\rm all}\) the cumulative epoch deaths. The restriction satisfies \(\pi D_{H_{\rm all}}=D_{H_{\rm all}}^{\rm phys}\pi\), with \(D_{H_{\rm all}}^{\rm phys}\) zeroing the corresponding physical coordinates. The current marks include the images of the complementary born part, so \(\sum_f F_f P^{S_1} A_f=0\) in the raw copy vector space. Clip every resulting atom using \(D_{H_{\rm all}}\). Write \(\lambda_t\) for its signed coefficient \(F_f\) times its catalogue probability. Then in that space the surviving frozen literal vector is
\[
 \widetilde F=\sum_{t\ {\rm affected}}\lambda_t(Y_t+R_t)
   +\left(\sum_{t\ {\rm affected}}\lambda_t(X_t-R_t)
           +\sum_{t\ {\rm unchanged}}\lambda_t D_{H_{\rm all}} U_t\right),
\]
Unchanged atoms avoid interval deaths, and their clipped images still lie on \(S_0\). After applying \(\pi\) the bracketed vector is supported on the base's surviving lifetimes and balances off \(S_0\). Apply the full cut repair comparison to that vector, without constructing it. Its norm is bounded by the raw total of the affected direct lists and replacements plus the unchanged atoms. Together with (7) this proves \(K_{S_1}\le\mathcal L(1+K_{S_0})\). The projected replacements need no new indexed kernel lists.

For two successive adaptive intervals the copy restriction for their union is the product of their individual restrictions. The selected base at the second interval can be the immediate snapshot or its admissible coarser ancestor. In either case the bracketed argument uses the actual earlier marks of that base, and the full cut acts on precisely that base's lists. A reset applies (7) anew with the current marks and all epoch deaths. It can retire selected-only ports after the new child graph is prepared and the old references have been fenced. Physical reprices use distinct names throughout this calculation, even when their input columns have the same image.

\hypertarget{local-indexing-and-clock-bounds}{%
\subsection{Local indexing and clock bounds}\label{local-indexing-and-clock-bounds}}

For each frozen group maintain the uncut catalogue, inverse descriptions of eligible state visits, and actual resource indices for the stored catalogue and for local kernel lists. Maintain dictionaries for marks and their offered distances at reachable states, current lifts, and heap scores, using path copying for snapshot ancestors. Exact score and price data along a live frozen path can be found from an \(O(k)\) scan, or a cached product-and-sum summary. One may index all physical resources of the short path handles. A cut finds affected lists before any raw record is retired. Across \(E\) edits the catalogue group count is \(E k^{O(1)}\mathcal L\) by §3.4. If \(C_{\rm old}\) bounds previous simultaneous visits to a resource, locating hit inherited lists costs \(O(E C_{\rm old})\) indexed visits. With \(p=\mathcal L\) macro occurrences in a hit list or a catalogue atom, the direct selector input from it has at most \(2p\) macro occurrences before splitting. It contributes at most \(2p\) conformal pieces, each with at most \(2p\) macros. Each stored macro can have \(O(k)\) physical resource occurrences. Interior fragments belong to the forest remainder with already promoted hole endpoints. Even a path cut into \(\Theta(k)\) disconnected fragments has at most two direct end portions. At a label each new binary or unary incidence has a first or last physical incidence, or a literal unary, in the list. Its load is bounded by the sum over at most \(D\) physical incidences.

Consequently the inputs of §2 have \(C_0=k^{O(1)}\mathcal L\) and \(J=\mathcal L\). Their aggregate explicit list size, including resource scans, is \(E k^{O(1)}\mathcal L\) for repairs and \(m_i k^{O(1)}\mathcal L\) for the all-mark calls. At an all-mark call there are \(O((m_i/k)\mathcal L)\) image labels. The full cut repair uses \(O(E)\) cut labels. The projected repair uses \(O(E)\) new labels after conformity on all marks. Factors of projected pieces still use their full lifts and resource lists, including their old endpoints. Each selected original piece promotes only the ends of its \(O(p)\) macros. This gives \(O(E\mathcal L)\) such promotions in either repair and \(O((m_i/k)\mathcal L)\) in an all-mark call. The finite priority table's multiplicity is included in list sizes and resource loads. Stored list lengths and occurrence bounds increase by a logarithmic power per snapshot ancestor for the full cut call. Inherited kernel list counts and lengths refer only to the full cut and all-mark calls.

Choose \(b_i\asymp m_i^{1/a^2}\) a power of two and use inner spacings \(s_j=b_i^j\). In each support window keep a root snapshot and the most recent snapshot at each spacing. At a tick multiple of \(s_j\), reconstruct level \(j\) from the last higher-spacing snapshot, using the interval ending at this tick. At coinciding spacings use the coarsest reconstruction first and take its result at all finer levels. A tick with no coarser update uses the level-one parent of the finest level, whose spacing is \(b_i\). Thus any repair interval at spacing \(s_j\) spans at most \(b_i s_j\) elemental edits up to the constant transaction size, and each ancestor chain has \(O(a^2)\) steps. There are at most \(\lfloor Q/s_j\rfloor\) calls at that spacing in a prefix of \(Q\) ticks, apart from the one root initialization. This follows by counting multiples. Only the final lowest snapshot for a settled tick is queried. The interval snapshots use persistent roots. Later physical support can retain selections beyond those in an earlier comparison ancestor. Holes and mark offers are applied from the selected ancestor's edit log. No comparison uses a dead lifetime, and every surviving previously selected raw factor remains available through the window.

Reset supports after a completed window of length comparable to \(m_i/(k b_i)\), or at an epoch boundary. The reset keeps all epoch holes and logical marks and does a new all-mark call. The work of a completed window's reset and retirements is charged to that window. Epochs end after a power-of-two number of ticks comparable to \(m_i/k\). One can align window boundaries and spacings by rounding down to integral powers of two. Each elementary edit transaction consists of a constant number of lifetime operations, with admission of a row and its cap indivisible and protected caps valid in every receiving state. Counting transactions or their operations changes the estimates by a constant. A larger inherited batch is processed as elementary transactions with a query fence until settlement.

Over \(Q\) ticks after a level initialization, the sum of repair interval lengths is \(O(Qb_i)\) times a logarithmic power. This follows by multiplying \(b_i s_j\) by the counted calls and summing the at most \(O(a^2)\) levels. Roots of windows need only one scan and support replacement per elapsed window. The hierarchy yields \(O(E\mathcal L)\) selections per interval, \(O(Qb_i\mathcal L)\) actual support additions and removals over the prefix, and \(O((m_i/k)\mathcal L)\) active support size within one window. Local work after initialization is \(Q k^{O(1)}b_i^{O(1)} M^{O(1/a)}\mathcal L\operatorname{poly}(B)\), and initialization has the same per-record bound times \(m_i\). These follow from (5) and the index counts, including scans to verify resources. Local coefficients only amplify their incoming bit height linearly times a polynomial in the short list dimension per ancestor, since their operations are short products and sums, conformal decompositions and exact one-row divisions. On \(O(a^2)\) ancestors this is a logarithmic-power multiplier. Fixed raw path summaries are recomputed from immutable chart words. These bit factors do not enter emitted record counts. The spacing formula is the one in the proof of Lemma 8.3, with selected real factors replacing its ordinary support chords. \href{https://raw.githubusercontent.com/openai/math/adc7f1241b42e322a6451854ab7e4b4c146bf78a/preprints/Almost-Linear-Time-Maximum-Cardinality-Matching-in-Sparse-General-Graphs-September-24-2026/main.pdf\#page=66}{OpenAI, Lemma 8.3}

\hypertarget{compilation-to-a-child-with-bounded-degree}{%
\subsection{Compilation to a child with bounded degree}\label{compilation-to-a-child-with-bounded-degree}}

Promote the initial portal grids and the logical marks in all their copies. Promote in their own copy the endpoints of a selected raw path and all literal attachments used by a selected factor. A retained literal is keyed by its copy and lifetime, so several selections of the same literal use one record. Represent in each live forest the Steiner subtree of its promoted ports by also marking branch vertices and compressing the two-port pieces to unique paths. Off this subtree a binary tree flow with no outside attachment has zero motion by stripping leaves and using the nonzero entry at each one. On a two-port piece conservation makes the flow a scalar multiple of its transported path: choosing the signed amount on its first edge fixes the only possible successive amount at each interior row. Each compressed path has \(O(k)\) edges because it avoids the initial grid in its interior.

Store at each physical-copy vertex its port or branch record, its current segment pointer if it is an interior vertex, or an absent flag. Each segment stores its ordered constituent list. For an interior promotion use the pointer and scan the old segment to make the two new segments. From a vertex outside the represented subtree scan the unrepresented piece up to its attachment, stopping at a represented vertex. Such a piece has at most one attachment since the underlying graph is a forest, and has \(O(k)\) vertices since it avoids the initial grid. In the absence of an attachment, install an isolated port. Otherwise install the joining path, splitting the segment at the attachment when it was an interior vertex. Update the pointers of the changed pieces by walking their old and new constituent lists. A local promotion or cut changes \(O(D)\) compressed records and takes \(O(kD)\) scans and pointer operations, per copy before dictionary overhead. Before a frozen cut, promote both ends, then remove their now-single-edge segment. In each surviving part the newly promoted endpoint replaces every former port reached across the cut for purposes of deciding which of its edges separate ports. Thus no pruning cascade is necessary. Keep all unexpired logical and cut marks through the epoch. Keep the private closing chords of used cycles as individual literals with promoted ends, so a hole in a cycle is an ordinary cut or a literal death.

At initialization and at a support reset, construct the represented subtrees in bulk from the final promoted sets. At a reset these sets keep the fixed grid, logical and cut ports and include the selected ports of the new all-mark call. Root each live tree and count its promoted vertices and the promoted descendants of each node. Retain an edge exactly when both sides contain a promoted vertex. Keep singleton promoted vertices, retain the branch vertices of the resulting subtrees, and walk each maximal two-port piece once to build its compressed record. Each retained branch vertex has at least three incident subtree edges, so there are at most as many branch vertices as promoted vertices. A scan of all live copies and literal attachments takes \(m_i\mathcal L\) operations apart from the local word arithmetic. Compare old and final records in dictionaries by copy and endpoint labels and their immutable lift data, keeping identical records and replacing changed ones with fresh lifetimes. Construct the final hub attachments from the final endpoint records and publish in the order in §5.2. Even if every old or final record changes, their total number is \(O((m_i/k_i)\mathcal L)\) by §4.3. This is an \(O(b_i\mathcal L)\) child edit charge per tick of the completed window. The bulk scan and arithmetic cost at most the all-mark local work per reset, which the same completed window pays with the \(k_i b_i\) multiplier already allowed in §4.3.

For a segment choose a dyadic column rescaling \(S\) so its price certificate lies in \((1/2,1]\). In the common chart units write its image as \((-a,b)\) at its ports. Write \(\widehat a,\widehat b\) for rounding toward zero on a fixed dyadic mesh \(\vartheta\le1/128\). Its exact child lift is
\[
   S+(a-\widehat a)Z_u+(\widehat b-b)Z_v .
\tag{9}
\]
Its assigned child price can be 2, bounding its exact raw price since the two adjustments cost at most \(2\vartheta\). Use the opposite cap corrections in the comparison of the unrounded \(S\). Their cost is a constant multiple of the original traversed price, using the lower normalization \(1/2\). If a rounded endpoint is zero this projected column is a unary with \(O(1)\) parent handles in its exact lift. If both are zero, put the lift in a separate local heap keyed by the structural segment lifetime, and use its opposite cap corrections in the child comparison. This heap updates when its segment changes. Equation (8) uses the actual uncompiled paths. Retain selected input literals or born literals directly at their live dyadic values. Every protected child unit column maps to the true parent's cap at its image row, with price one and score zero.

For an ordinary conserved comparison, the total image over all copies of a physical row is zero after (9) and the local differences. Join copies to one hub per physical row using unit nulls, which have zero physical lift. Give nulls a dyadic price between \(1/(8\Xi_i)\) and \(1/(4\Xi_i)\), where \(\Xi_i\) is the parent's entering cap-unit envelope. If the hub has many new incidences use a binary trie on \textbf{local slot indices}. Its depth is \(O(\log M)\) and its degree is bounded. Assign each incident copy or born literal its own leaf. Give trie nodes protected caps in the same inherited units. Sending copy imbalances to the hub conserves the child vector and has price at most \(O(\log M)\) times the predecessor raw price. The endpoint mass of a path or literal is bounded by \(\Xi_i\) times its price, and its hub route has at most \(O(\log M)\) nulls. This argument applies also to the bounded extra cap corrections. A null lift and the other lifts are priced contractions. Their incidences satisfy \(A_i P_i=I_\phi A_{i+1}\), where \(P_i\) is the exact linear lift and \(I_\phi\) identifies child copies with their mapped parent row.

The next envelope is at most \(8\Xi_i\). A rounded end of a scaled path is bounded by the original entering or leaving raw incidence at that end. Relative to its certified price its efficiency is bounded by \(\Xi_i\). A unary lift has its image bounded by the total priced incidence. True caps have efficiency one, and the above null range gives the bound for unit nulls. A literal and a null have immutable dyadic data. Each new row comes with its mapped protected cap as one elementary admission. Give each receiving instance one permanent row with only its protected unit coordinate, mapping it to any parent cap live at the start of the parent epoch. This keeps the receiving chart nonempty even when no factor is selected. This coordinate is zero in every conserved comparison by its row equation, and adds one row and one column to the charged initialization. When a parent has no constraints at the root, the direct rule in §6.3 applies.

A retired port first loses all incidents other than its designated cap, including its null attachments. That cap remains a live isolated column of the receiving chart through the receiving epoch. Its coefficient in a conserved receiving vector is zero by the receiving row equation. The cap is still an eligible named lift for the receiving chart's frozen geometry. Distinct receiving caps over the same parent row have distinct receiving lifetimes and row equations. At the receiving epoch fence, discard the old geometry and flush its own descendant and client maps using their old mappings, then purge pending caps and freeze the new input. At a parent epoch fence, tear down and flush the receiving instance before retiring any parent row or parent cap that it refers to, including the cap used by its permanent row. Charge each pending cap purge to the corresponding retirement. The holding count is at most the live count at the beginning of the receiving epoch plus the \(O(m_{i+1}/k_{i+1})\) admissions during it.

Allocate slot indices from a structural arena, reusing one only after its old maps have been flushed and the snapshots or epochs using it have ended. A generation tag for the identity can have \(O(\log\operatorname{poly}(M,B))\) bits and is separate from the local trie index. In a parent epoch the total distinct issued records before such reuse is \(O((m_i/k)b_i\mathcal L)\), within the next capacity with slack under §7. Incoming price, score, incidence or endpoint changes use deletion and birth of a fresh lifetime. Fixed-size row admissions and delete-first operations keep transient degree and slot bounds. The procedure yields
\[
 z=\sum_j c_j C_j+P_i z',\quad A_{i+1}z'=0,\qquad
 \sum_j |c_j|L(C_j)+\|\ell_{i+1}z'\|_1\le \kappa_i\|\ell_i z\|_1,\quad
 \kappa_i=\mathcal L ,
\tag{10}
\]
for every conserved \(z\), including any local zero-image structural lifts among the \(C_j\). The relative price and short-list bounds are independent of \(B\). The forest promotion and deletion argument is a direct version of the structural part of Lemma 8.4. \href{https://raw.githubusercontent.com/openai/math/adc7f1241b42e322a6451854ab7e4b4c146bf78a/preprints/Almost-Linear-Time-Maximum-Cardinality-Matching-in-Sparse-General-Graphs-September-24-2026/main.pdf\#page=67}{OpenAI, Lemma 8.4}

\hypertarget{recursive-queries-path-lifting-and-lifetime-logs}{%
\section{Recursive queries, path lifting, and lifetime logs}\label{recursive-queries-path-lifting-and-lifetime-logs}}

\hypertarget{exact-path-lifts-of-compensated-columns}{%
\subsection{Exact path lifts of compensated columns}\label{exact-path-lifts-of-compensated-columns}}

Use public forests of oriented binary edges and separate individual pieces for one-entry columns and cycle factors. Besides copies of real binaries and unit-transfer nulls, a public forest may use an edge between two copies of the same row with gain \(\sigma>0\). With unit input its physical lift is \((\sigma-1)Z_v\), in the inherited units, using the designated true cap. It pays its absolute protected price and has zero score. For a rounded segment (9) with both ends nonzero the parent realization has a first such edge of gain \(a/\widehat a\) and a last of gain \(\widehat b/b\), with the old path between them. A coefficient \(c\) in the child supplies mass \(c\widehat a\) in the forward direction. The first extra edge gives mass \(ca\), the old path receives coefficient \(c\) and outputs \(cb\), and the last restores \(c\widehat b\). Their cap motion is exactly \(c(a-\widehat a)Z_u+c(\widehat b-b)Z_v\). For negative moves use the original path with a negative amount, or the corresponding signed reverse traversal with its true reciprocal gain. If one rounded entry vanishes, the full lift is an individual piece consisting of the parent path and its cap corrections. No such unary is used as a binary path hop. The analogous zero lift is a local list. A literal cap-gain edge of the child becomes a single parent edge of the same gain, since the protected unit column maps to the parent's own protected unit at the appropriate row. Numeric descriptors of these inherited edges can be forwarded by reference.

In addition to the full frozen forest copies with current holes, provide the following occurrence representation. For each active child label allocate one slot per allowed child public occurrence over the label and one permanent slot for individual lifts. In each slot give each incident binary with nonzero immediate lift an independent arm of its exact parent route outward from the slot center, adding the two cap-gain edges where needed. The far tip and the center are indexed. For a child public edge representing such a binary, choose one endpoint side in advance and connect that arm's far tip to the assigned slot center of the other copy by a zero-lift connector. A binary whose lift uses a live null coordinate of the immediate parent still uses its arm. That coordinate has an incidence column in the parent chart. A child binary with zero immediate lift has equal endpoint magnitudes at two rows mapping to the same parent row, by \(A_iP_i=I_\phi A_{i+1}\). Represent it by a direct unit-transfer connector, and omit an individual bottom piece with this zero lift. The child's auxiliary unit-transfer connectors also lift by direct connectors. For an inherited cap-gain edge, connect the designated centers with the same gain and the appropriate mapped protected unit. Slots for unchanged labels keep their arms when empty or reassigned. Forest arms of different slots are disjoint. Contracting each slot tree gives the current child public forest, so their union with the connectors is a forest. The permanent individual-arm slots provide endpoints for a receiving bottom LP's individual binary with nonzero lift as well. When the visible recorded child graph contains isolated retired-cap rows their conserved coefficients are zero. A forest component over an inactive row consisting entirely of zero-lift unit connectors can be omitted. Retain its passive identity until the receiving epoch or activation discipline permits reuse.

Expand a child path conceptually by these arms and connectors. The signed lift agrees with its successive transported amounts. An adjacent reversal uses the delivery of the first traversal, takes the opposite signed coefficient, and returns the initial amount. For a cap-gain edge the motions are \(x(\sigma-1)Z_v\) and \(\sigma x(\sigma^{-1}-1)Z_v\), which sum to zero. Repeating this erases every backtrack of a forest walk and gives its unique simple path with the same signed lift and no higher priced absolute traversal. This holds for every subpath between two child occurrence vertices, including vertices lying inside arms built by the receiving chart. Such vertices also have indexed slots over their chart-row mappings. A path piece records an ordered pair of public occurrence IDs and its exact signed amount at the first ID. To lift it, look up the two assigned centers and pass the same start amount. An individual binary uses the indexed permanent arm. Lift each individual unary or zero piece by its parent list, which is a literal or a segment path and at most two individual caps. At the root, each remaining individual piece names a root column and gives its exact signed coefficient. The local list size is logarithmic-power, the bottom LP gives at most \(O(m_d)\) pieces, and lifting a binary path keeps one piece. Any individual piece can branch by at most a fixed constant at each level. Since the depth \(d=O(a\log a)\) is fixed on each branch, a query at the root returns at most \(m_d\log^{O_a(1)}M\) pieces. It chooses a single best witness at each node. \href{https://raw.githubusercontent.com/openai/math/adc7f1241b42e322a6451854ab7e4b4c146bf78a/preprints/Almost-Linear-Time-Maximum-Cardinality-Matching-in-Sparse-General-Graphs-September-24-2026/main.pdf\#page=69}{OpenAI, Lemma 8.5 and construction}

\hypertarget{public-multiplicity-and-old-to-final-transactions}{%
\subsection{Public multiplicity and old-to-final transactions}\label{public-multiplicity-and-old-to-final-transactions}}

Let \(D_i^{\rm pub}\) bound the number of simultaneous public vertices and input-edge copies per input row or edge, and let \(F_i=\log^{O_a(1)}M\) absorb the frozen copy count, port multiplicity and incidence degree of chart \(i\). In one compressed raw forest a physical interior edge lies in one segment and occurs at most twice per arm-slot index, once from each port. At a port its own centers and the incident arm tips have a degree-factor bound. Each cap-gain adjustment spends only a constant number of extra edges and vertices per arm. Connectors use only existing slot centers or tips. The same induced-forest degree count below bounds all connectors over one label. Including full copies and permanent individual slots, take \(D_i^{\rm pub}=O(F_i(1+D_{i+1}^{\rm pub}))\), with the bottom public forest empty. This is logarithmic-power for a fixed branch. A tracked unary has its base as an individual value. Only protected caps occur in cap-gain edges. The number of public records is the chart capacity times a logarithmic power. These counts involve no numerical bands or grid states. \href{https://raw.githubusercontent.com/openai/math/adc7f1241b42e322a6451854ab7e4b4c146bf78a/preprints/Almost-Linear-Time-Maximum-Cardinality-Matching-in-Sparse-General-Graphs-September-24-2026/main.pdf\#page=71}{OpenAI, Lemma 8.6}

Keep old and final passive dictionaries of the descendant's explicit forest logs and incidence lists by immutable lifetime. The changes to the \textbf{abstract child graph} determine the set \(Q\) of child labels whose arm pools must be rebuilt, including endpoints of changed segment lifts, assignments, and literal changes. Process the child edits without asking a path query, bookkeeping each intermediate log in a dictionary. At \(Q\) sweep the old and final incident edges, and build final pools. For one label \(q\) every incident forest edge has both ends in the union of occurrence sets over \(q\) and its current distinct abstract neighbors, including same-label connectors. The induced graph there is a forest, so the number of such edges is at most \(D_{i+1}^{\rm pub}(1+\deg(q))\). This bound is for the entire sweep, even if one public occurrence has large degree. All the arms have \(O(k_i)\) constituents. At unaffected labels retain maps and assignments for retained occurrences, free deleted assignments, and assign the new occurrences to the free slots. A descendant-internal log at these labels causes at most a constant number of connector operations and dictionary changes. Do not rebuild a pool on each inherited cap-gain log. A transient child lifetime that was neither old nor final is charged for the log processing and need not be exposed to the parent client.

This gives a prefix charge \(O(|Q|D_{i+1}^{\rm pub}(k_i+1)\mathcal L+L_{\rm log}\mathcal L)\), plus local full-copy changes. Here \(L_{\rm log}\) counts all processed child logs. To publish, remove old connectors first, then removed arms and full-copy edges, and all incident edges before any old vertex. Preserve the old input maps and prices for every removal. After these are flushed, install the final nodes, arms and connectors. The retained edges are a subforest of the final forest, so links join different components. Treat a root reprice as an old-lifetime death and a birth with the same original box coordinate. Its accumulated history follows the physical box coordinate, not the chart slot identifier. On an epoch reset release all dead maps after their own logs have been processed. Allocation can use reusable arenas only after the last live reference, including a frozen snapshot, has ended. The same construction applies to the cap placeholders since their eventual retirement follows their incident life ends. This is an explicit old-to-final construction with local degree and path scans charged to abstract changes, and logarithmic work charged to generic descendant logs. \href{https://raw.githubusercontent.com/openai/math/adc7f1241b42e322a6451854ab7e4b4c146bf78a/preprints/Almost-Linear-Time-Maximum-Cardinality-Matching-in-Sparse-General-Graphs-September-24-2026/main.pdf\#page=71}{OpenAI, Lemma 8.7}

\hypertarget{scores-and-exact-query-certificates}{%
\subsection{Scores and exact query certificates}\label{scores-and-exact-query-certificates}}

Use descent as the score convention, \(-g_i^\top z\). For a local kernel list \(C\) with certificate \(V=L(C)>0\), put \(|g_i^\top C|/V\) in the heap and store the sign giving nonnegative descent. A zero-price list has zero score, since every constituent chart coordinate has positive price, and can be dropped before forming a ratio. The empty answer has an empty piece list, score zero and price zero. Use it when no positive score is available, without evaluating a ratio for it. The root uses the mesh \(h_g\) of §6.4. For all receiving charts choose a second dyadic mesh at most \(\epsilon_g\ell_{\rm low}/2\), where \(\ell_{\rm low}=\min(1/2,\ell_{\min},1/(8^{d+2}\Xi_0))\) is a common positive lower price bound. Here \(d\) is the recursion depth and \(\epsilon_g\) is set using the root mesh below. For a child input compute \(g_i^\top P_i(:,e)\) with the immediate exact \(O(k_i)\)-constituent lift and round it to that receiving mesh. This gives error at most \(\epsilon_g\ell_{i+1,e}\). Protected and null scores are zero. All rule changes to a score use immutable lifetime replacements. The maximum absolute gradient-to-price ratio increases by at most \(\epsilon_g\) per level, since \(P_i\) is priced nonexpansive. Exposed price words are one, two, original dyadic literals, or nulls from the prescribed dyadic interval. Exposed incident magnitudes are the rounded ones in (9) or inherited exact words. The bound \(\Xi_i\le8^i\Xi_0\) and the chosen score meshes give height \(O_a(1)\) times the root bound plus a logarithmic power in \(M\) for prices, magnitudes and scores. Exact segment and repair data are local only. Products of a bounded-length list give bit height linear in entering height up to polynomial list factors. Section 6 approximates longer lifted path products locally.

At the bottom optimize the signed ratio by the LP on \(z^+,z^-\ge0\) with \(A_d(z^+-z^-)=0\), \(\sum_e\ell_{d,e}(z^+_e+z^-_e)\le1\), and linear objective \(-g_d^\top(z^+-z^-)\). It is rational, feasible at zero, and bounded by positive prices. See §7.4 for the exact vertex optimizer. If the optimum is positive, return the signed individual columns with price certificate \(V=1\) and advertise the exact objective as the score. An optimal comparison vector of raw price one is feasible for this LP by taking its positive and negative parts. Every LP solution represents a vector of raw price at most one by the triangle inequality. Hence the LP value equals the optimal ratio, with zero included. For a zero optimum use the empty answer. At a higher level compare its best local list certificate with a single child answer. For a nonempty child answer with advertised score \(s\) and certificate \(V>0\) for signed vector \(z'\), its lift obeys
\[
  -g_i^\top P_i z'\ \ge\ -g_{i+1}^\top z'-\epsilon_g\sum_e\ell_{i+1,e}|z'_e|
  \ \ge\ (s-\epsilon_g)V .
\]
The price of the displayed parent pieces is at most the same \(V\), including cap-gain motions. The exact physical image is zero. When \(s-\epsilon_g>0\), advertise \(s-\epsilon_g\) with the same \(V\) for that lift. Otherwise treat the child option as empty. Use the best positive option, or the empty answer if there is none.

Let \(\mathcal R_i\) be the true optimal nonnegative score per unit raw price on current chart \(i\), \(J_i\) its best local certificate, and \(K_i=\prod_{j\ge i}^{d-1}\kappa_j\). Equation (10) gives
\[
 \mathcal R_i\le \kappa_i\max\{J_i,\mathcal R_{i+1}+\epsilon_g\}.
\]
Take \(J_i=0\) when the local option is empty, \(K_d=1\), and all \(\kappa_i\ge1\). At the bottom the returned score is \(\mathcal R_d\). If the child score has the claimed lower bound \(\mathcal R_{i+1}/K_{i+1}-2(d-i-1)\epsilon_g\), the parent returns at least the maximum of \(J_i\) and that bound minus \(\epsilon_g\), even when the child option is empty. Since \(K_{i+1}\ge1\), (10) and the ratio inequality above give a score at least \(\mathcal R_i/K_i-2(d-i)\epsilon_g\). For a positive root optimum at least one real root score word is nonzero. If its fixed mesh is \(h_g\), that magnitude is at least \(h_g\). Complete the one column by its protected unit caps. The score stays that of the original and the price is at most \((1+2\Xi_0)\ell_{\max}\), including the unary and zero cases. Thus \(r_{\min}=h_g/((1+2\Xi_0)\ell_{\max})\) is a lower bound for positive root optimum. For \(d>0\), take the largest dyadic power with \(2d\epsilon_g\le r_{\min}/(2K_0)\), using a declared upper bound for \(K_0\). If \(d=0\) no receiving mesh is needed. This gives a root factor \(2K_0\), with exact conservation and certified raw price. The returned start amounts of explicit pieces use local rational arithmetic. Lifting through the fixed depth only passes those amounts and multiplies by bounded-length local coefficients for individual factors. Their bit heights add along the chain, up to logarithmic powers. No recursive multiplication of bit polynomial degrees is needed.

\hypertarget{a-two-level-reprice-with-an-actual-endpoint-correction}{%
\subsection{A two-level reprice with an actual endpoint correction}\label{a-two-level-reprice-with-an-actual-endpoint-correction}}

For a local trace take an old raw forest segment from \(a\) through \(x\) to \(b\), with \(f:a\to x\) of gain \(1/2\), raw price \(H_b=2^{30}\), and score 3, and \(q:x\to b\) of gain \(\lambda=3/4+2^{-64}\), price 2 and score 8. Use root row multiplier \(U_b=2^{40}\), so every real incidence in this example is multiplied by \(U_b\), and protected coordinates have incidence one in those units. Assume old ports at \(a,b\), with \(x\) still in their segment. This is permitted by the portal rule by padding the corresponding branches, or can be used directly as an interior two-port piece of an existing chart. Keep genuine cap resources live throughout.

The old child segment is the \((2H_b)^{-1}\)-scaled lift of \(f+\tfrac12 q\). Its true old certificate lies in \((1/2,1]\) before rounding. Its image is \((-512,192+2^{-56})\). On the \(1/128\) mesh it emits image \((-512,192)\), price 2, raw lift \(S=(f+\tfrac12 q)/(2H_b)-2^{-56}Z_b\), with score \(7/(2H_b)\) on a fine enough score mesh. Its head cap-gain edge has gain \(192/(192+2^{-56})\). For example an inherited local list containing \(S+512 Z_A-192 Z_B\) at the receiving level is exactly conserved with \(Z_A,Z_B\) the live units at its two child ends. The formula also holds when those units are genuine alternative cap lifts with zero score. The cut rule applies to this list and to every other hit indexed kernel.

After any settled query choose \(f\)'s reprice to price \(\epsilon=2^{-30}\) and score 9. Remove \(f\) everywhere and add \(f^+\) as a fresh literal. The two cut endpoints become logical marks and \(x\)'s relevant copies become ports. The old segment \(S\) retires. The surviving portion \(q\) compiles as \(\widehat Q= q/2-2^{-25}Z_b\), with image \((-U_b/2,3U_b/8)\) and score 4. There is a new child row for \(x\) with its protected cap before the births \(\widehat Q,f^+\). The old lifetime of \(q\) in the input is live, although the compiled segment name and its old coefficient change. On receipt of the death of \(S\), the example old kernel clips to \(512Z_A-192Z_B\). This full image is repairable by the same live units, so its old score is absent. Projected matching on the newly marked rows may leave residual at the old logical marks and uses (8) there. The child new literals remain at their actual recorded prices, including the newly cheap \(f^+\). Their assembly never uses the removed \(f\).

One conserved receiving witness is \(f^++\widehat Q+U_b Z_A-(3U_b/8) Z_B\). Its lift has real score 13 and raw physical cap coefficient at \(b\) equal to \(-3U_b/8-2^{-25}=-U_b\lambda/2\). Both rows and the newly exposed \(x\) balance exactly. Its price includes both true cap motions, the live half of \(q\) and the new price \(\epsilon\). Hubs, when present, merely add nulls to this calculation. The final parent public path maps use \(\widehat Q\)'s new endpoint adjustment. Remove connectors that mentioned \(S\)'s old map before retiring its slot, flush real histories on any removed \(f,q\) copies, and use the old and final incoming graph states as in §5.2. Intermediate receiving logical snapshots are processed before any query. A later support-window reset retires selected-only ports and reassembles the full catalogue with this death and all these logical marks still in force. Thus old protected lifetimes and any retained dormant arm histories are governed by their actual maps, and a new price word is never written over a live old path.

\hypertarget{finite-word-motion-and-additive-optimization}{%
\section{Finite-word motion and additive optimization}\label{finite-word-motion-and-additive-optimization}}

\hypertarget{a-reference-trajectory-and-exposed-path-algebra}{%
\subsection{A reference trajectory and exposed path algebra}\label{a-reference-trajectory-and-exposed-path-algebra}}

The true box coordinates are driven by a reference sequence. It starts at the exact dyadic point specified below. Each accepted step applies the nominal exact signed physical gains and start amounts of its displayed witness. In this \emph{definition} apply only the true box coordinates, omitting the protected and null coordinates. The algorithm keeps approximate histories of that same sequence. Point reads can influence future directions, leaving prior movements as defined. Queries, edit choices, phase changes and stopping decisions may use the reads.

The represented roots of the link-cut forest are occurrence vertices. An edge node stores whether its represented parent side is its fixed tail or head. This determines its transfer and unit coefficient in the down-path order of an auxiliary tree. Even an auxiliary segment beginning or ending at an edge node uses this local factor once, in the position occupied by that node. Occurrence vertices have transfer one and no motion. Exposing a path changes preferred edges and leaves represented parent directions intact. Everting an occurrence vertex reverses the whole exposed root path and flips the local edge orientations on it, with lazy propagation alongside path reversal. Every edge node on that root path has both of its adjacent occurrence vertices on the path.

To insert a subdivider, attach its isolated node below the fixed tail with the forward direction, evert the head in its distinct component, and attach the head below the new node. To retire an edge, flush its history while its map is still live. Read its orientation by point access, cut its child occurrence vertex from it, and then cut the subdivider from its parent. There is no evert at an edge node and no motion or query between the two half-edge operations. Other point operations access an edge node without everting it. These are the insertion and deletion orientations of the matching preprint's §5.4. The gain and coefficient fields below use the indicated direction. \href{https://raw.githubusercontent.com/openai/math/adc7f1241b42e322a6451854ab7e4b4c146bf78a/preprints/Almost-Linear-Time-Maximum-Cardinality-Matching-in-Sparse-General-Graphs-September-24-2026/main.pdf\#page=37}{OpenAI, §5.4 and Lemma 5.3}

Use a link-cut forest with subdivided edge nodes and positional lazy tags. For an exposed ordered path \(P\) let \(\tau_P>0\) be its transfer and \(c_{e,P}>0\) the absolute physical coefficient at a real mapped edge \(e\) for unit input including its prefix transfer. An edge with image \((-a,b)\) has forward and reverse unit coefficients \(1/a,1/b\). A null or a protected gain edge contributes its transfer but no tracked real motion. Maintain both ordered transfer summaries by multiplication along disjoint children and the local edge. Evert swaps the two summaries, without taking a reciprocal of a rounded path product. At a real copy keep two nonnegative physical histories and a gross age \(A_e\). A pending path tag consists of nonnegative starting amounts \(h^+,h^-\) and an independent gross amount \(b\). On a current real edge its signed effect is the signed traversal with amount \(h^+-h^-\), and the gross increment is \(b c_{e,P}\). Push tags to a suffix by multiplying its three fields by the prefix transfer. At an edge allocate the two signed channels according to its direction, never negating existing physical histories. Under reversal the tag becomes
\[
   (\tau_P h^-,\,\tau_P h^+,\,\tau_P b).
\]
This is because \(d_{P^{\rm rev}}=-d_P/\tau_P\) for the unit-input signed vector. Tags compose by nonnegative addition. Before changing auxiliary children, push their positional tags. Signed and gross fields are kept separate even if physical signed movements can cancel.

Each enabled real copy has an exact positive dyadic threshold \(\theta_e\). On an oriented path maintain a witnessed lower bound for
\[
     J_P=\min_{\text{enabled }e}(\theta_e-A_e)/c_{e,P}.
\]
Omit null and protected copies. Represent an empty enabled set by a separate \(+\infty\) sentinel with no witness. Finite scaling and subtraction leave that sentinel unchanged, and the minimum skips it if a finite candidate exists. The exact join formula for \(PQ\) is \(\min(J_P,J_Q/\tau_P)\) with a corresponding formula for the opposite order. A path gross update \(b\) decreases the forward minimum by \(b\) and the reverse one by \(\tau_P b\), because \(c_{e,P^{\rm rev}}=c_{e,P}/\tau_P\). Disabled copies still receive increments. The underlying edge-node forest operations are those of OpenAI Lemma 5.3. That lemma's splay potential counts pointer changes and assumes constant-many aggregate and tag transformations per primitive. Its proof uses path-parent pointers and preferred-path aggregates only. The fields and identities above meet that assumption, add no subtree operation, and preserve the legal link, cut, evert and point operation order in §5.2. Thus \(O(\log(2+N_{\rm f}))\) amortized tree primitives suffice per such operation or reported witness, where \(N_{\rm f}\) bounds live subdivided nodes. The hypotheses on word operations are implemented with charged rational or rounded arithmetic below. \href{https://raw.githubusercontent.com/openai/math/adc7f1241b42e322a6451854ab7e4b4c146bf78a/preprints/Almost-Linear-Time-Maximum-Cardinality-Matching-in-Sparse-General-Graphs-September-24-2026/main.pdf\#page=37}{OpenAI, Lemma 5.3}

\hypertarget{directed-arithmetic-and-witnessed-alerts}{%
\subsection{Directed arithmetic and witnessed alerts}\label{directed-arithmetic-and-witnessed-alerts}}

Use signed binary floats of mantissa length \(p\), signed integral exponents and exact zero. Products for both directions use outward primitive intervals from the exact local descriptors. For positive scales in \([\lambda^-,\lambda^+]\), a lower signed \(j\) multiplied by the scale uses \(\lambda^+\) if \(j<0\) and \(\lambda^-\) otherwise. For dividing a lower \(j\) by a positive denominator interval \([\tau^-,\tau^+]\), use \(\tau^-\) if \(j<0\) and \(\tau^+\) otherwise. Round the selected result downward. For upward signed bounds interchange endpoints and round upward. In particular the minimum-join division uses the correct sign at every lineage step. To subtract the gross increment, first bound its product upward, then subtract with downward rounding. Leaf tests start using the upper age, exact \(\theta_e\), and interval coefficient. Carry the actual selected witness through joins and updates. After a path update, while its lower minimum is nonpositive, record its witness and disable it by a point operation until reanchoring, re-exposing the path after each disable. Every reached true threshold is detected because the stored minimum is a lower bound. A path operation without a report cannot have increased an age outside that path.

The following bounds justify that false early alerts still have a motion charge. Fix a coarse polynomial guard \(H_\ast\) on all per-run operations and simultaneous live records, including a partly attempted operation, as specified in §7. On a guarded prefix let \(u=O(2^{-p})\) denote unit relative roundoff. Each transfer product uses its constituent leaf intervals once in a parenthesization computed on the current path, not a product updated from a previous rounded quotient. This gives relative interval width \(O(H_\ast u)\) by multiplying the per-leaf and per-join factors. For any exact nonnegative history, tag or age, its upper estimate is at most \(1+O(H_\ast^2u)\) times its value. Nonnegative sums preserve the largest relative error up to roundoff, each scaling adds \(O(H_\ast u)\), and any contribution crosses at most \(H_\ast\) operations. This includes upward sums on copy flushing and point reads. Exact zeros stay zero.

For a minimum's witness in its current oriented units, write \(\Theta=\theta_e/c_{e,P}\), \(B_e=A_e/c_{e,P}\). The stored lower bound is \(\Theta-B_e-\zeta_e\) for \(\zeta_e\ge0\) along that witness's lineage. A positive rescaling scales both nonnegative terms together, adding at most its exact scale times \(O(H_\ast u)(\Theta+B_e+\zeta_e)\) to the carried error. A subtraction of an upper gross estimate adds at most \(O(H_\ast^2u)\) times the new \(\Theta+B_e+\zeta_e\). Here the new \(B_e\) has only increased under the subtract operation. On a join follow the chosen witness's own rescaled pair. At a leaf recompute with its certified age. Across a guarded lineage these inequalities give \(\zeta_e\le O(H_\ast^3u)(\Theta+B_e)\). The normalization does not measure relative error against a nearly zero difference. Taking \(\epsilon_{\rm f}\ge O(H_\ast^3u)\) smaller than \(1/10\), a nonpositive returned bound then implies \(A_e\ge(1-\epsilon_{\rm f})\theta_e/(1+\epsilon_{\rm f})\) at that witness. Witnesses remain enabled mapped copies until the report or physical removal. All inequalities hold for adaptive histories as they are deterministic conditional on the tracked exact sequence.

The signed exponent guard can be explicit. All primitive transfer descriptors and exact starting amounts have bit heights bounded by a fixed polynomial of current small-list dimensions and exposed word height. Hence individually their logarithmic magnitudes are bounded by a fixed polynomial of \(M,B\). A path length is at most the guarded record count, so any path-product exponent and any one upward multiplication shift in a tag is also polynomial in \(M,B\). Positive sums can grow a magnitude by at most a factor two of their maximum per addition. Signed cancellations of two \(p\)-bit aligned floats can lose at most \(O(p)\) additional exponent positions when their magnitudes are comparable. For distant mantissas use a sticky bit with directed rounding. This bounds exponents over at most \(H_\ast\) operations by a fixed polynomial in \(M,B,p\), which needs only \(O(\log(2+M+B+p))\) exponent bits on an adequate branch. The exact descriptors can be multiplied or divided by entity-local rational arithmetic, then rounded once at their creation. Reverse transfer intervals are determined there separately. Inherited descriptors are forwarded without a fresh rounded reciprocal or a rescanning of the flattened arm. Word-range and zero tests are exact. The exponent range preserves finite positive values.

Each true physical variable has two base histories, a gross bucket for explicit movements and flushed copy ages, and a list of its live real copies. For a scaled individual piece with coefficient \(q\) on a real root variable, add \(\max(q,0)\) and \(\max(-q,0)\) to its two base channels and add \(|q|\) to its gross bucket, with the directed arithmetic above. This applies also to real unary and zero-image columns. Null and protected individual pieces contribute no tracked movement. For a scaled path piece of signed starting amount \(q\), apply channels \(\max(q,0),\max(-q,0)\) and gross amount \(|q|\) on the ordered path. A point read adds the initial value and the difference of the two channel sums, each including the base and all live copies.

Remove a copy by flushing its two histories and gross age to the correct old base and bucket before its map dies. Check the bucket's upper estimate against the current threshold after each direct addition or flush. Add the variable to the pending queue if the estimate reaches that threshold. A new real copy starts with zero histories and zero gross age, enabled at its variable's current threshold. When an unaffected dormant slot is reassigned, retained copies preserve their histories, ages and enabled status. Whenever an enabled minimum or the bucket flags a variable, coalesce a pending notification until all pieces of the current step finish. At service read its reference value to the certified absolute accuracy, reset the copies' gross ages and the bucket, reenable them, and compute its new stale words. Use the largest dyadic power at most
\[
     \frac{\alpha}{10^8(D_{\rm pub}+1)K\ell_e},
\tag{11}
\]
for the new \(\theta_e\), where \(\ell_e\) is its new price and \(D_{\rm pub}\) is the declared simultaneous root copy bound. Install this threshold on its surviving and new copies. All its own removals during the ensuing reprice flush zero age, while other removals may trigger a bucket. Signed histories always retain their contributions, including when transferred from copies to bases. At settlement each bucket or copy is below threshold by the conservative witness and bucket checks. The true displacement since the anchor is at most their sum. A report or disable charges a constant fraction of \(\ell_e\theta_e\) to the exact age of its witness. For a bucket, its upper estimate is at most \(1+O(H_\ast^2u)\) times its exact age, so a bucket report has the same constant-fraction charge. Motion spent on a killed copy can be charged once at it and once when it enters the bucket. Reanchoring erases that age. Thus every prefix of \(T'\) steps of root raw movement at most \(v_\ast=\alpha/(1000K)\) per step generates \(O((D_{\rm pub}+1)T')\) total charged disables and reports, besides explicit sweeps. A service takes \(O((D_{\rm pub}+1)\log(2+N_{\rm f}))\) amortized primitives plus the graph edits, and ordinary log flushing costs one point operation per affected copy.

Drain the pending queue after the step and before the next query, including entries created by log flushing while the queue is being drained. A copy or bucket's debt carries the variable's price from its current anchor. That price is fixed until its ages are cleared. After service the outgoing reprice flushes the variable's signed history and zero gross ages using its old map. For every other variable the reprice can move debt into its bucket at the same anchor price. This charges the positive debt of each actual movement at most once to a copy alert and once to a bucket alert. Phase sweeps clear gross debt after the old maps are flushed and reconstruct all anchors, with a separate sweep charge.

For accepted steps bounded by \(T\) with prices at least \(\ell_{\min}>0\), total exact gross tracked motion is at most \(T v_\ast/\ell_{\min}\). Choose \(p\) so \(O(H_\ast^4 2^{-p})\) times this bound plus the initial-value and final summation rounding is at most a desired \(\epsilon_x\). In particular
\[
 p=O\!\left(1+\log H_\ast+\log(2+T v_\ast/\ell_{\min})+\log(2+R)+\log(2+1/\epsilon_x)\right)
\tag{12}
\]
with a sufficiently large absolute coefficient and \(R\) bounding box values. If the two signed-channel totals at a read are \(H^+,H^-\), the absolute error of the signed result is \(O(H_\ast^4 2^{-p})(R+H^++H^-)\), including the final rounded subtraction. This uses the sum of the channel magnitudes as the error scale. Formula (12) ensures all reads refer to one shared exact trajectory through arbitrary reprices, old-map flushes and phase resets. This generalizes the threshold ledger in the proof of OpenAI Lemma 5.4. The gain summaries, rounding and reference trajectory were proved here. \href{https://raw.githubusercontent.com/openai/math/adc7f1241b42e322a6451854ab7e4b4c146bf78a/preprints/Almost-Linear-Time-Maximum-Cardinality-Matching-in-Sparse-General-Graphs-September-24-2026/main.pdf\#page=39}{OpenAI, Lemma 5.4}

\hypertarget{initialization-and-choices-of-barrier-parameters}{%
\subsection{Initialization and choices of barrier parameters}\label{initialization-and-choices-of-barrier-parameters}}

Consider an exactly feasible dyadic instance \(Ax=b,\ 0\le x\le u\) with linear cost and tolerance \(0<\eta\le1\) from §1. Remove zero widths. An empty variable set gives the empty solution, and with no constraints take the independent minimizing box corners. In this subroutine \(n\) denotes the assembled LP row count and \(m_{\rm box}\) its variable count before adding the artificials. Otherwise let \(Q_0\) be the least power of two bounding \(N,m_{\rm box}+n+1,1/\eta\), all absolute input magnitudes and all inverse positive widths. It has \(O(1+\widehat B)\) bits. On the original box, the cost magnitude is bounded by \(m_{\rm box}Q_0^2\le Q_0^3\), and \(\|b-Ax\|_\infty\le Q_0+m_{\rm box}Q_0^2\le2Q_0^3\). Put \(R=16Q_0^3\) and take \(P_{\rm art}\) to be the least power of two at least \(64Q_0^4/\eta\). Add \(z_v^+,z_v^-\in[0,R]\) at cost \(P_{\rm art}\) with columns \(+e_v,-e_v\). Write \(X\), \(A_+\), \(C(X)\), \(C_*\) for the enlarged actual box vector, matrix, cost and exact optimum, and \(m_+\) for its number of variables. A strictly interior start is
\[
 x=u/2,\qquad z_v^\pm=R/2\ \pm(b-Ax)_v/2 .
\]
It is exactly feasible and the artificial slacks are at least \(R/4\), since each original box point has maximum residual at most \(2Q_0^3\). The exact original optimum gives an enlarged feasible witness with zero artificials and \(C_*\le\mathrm{OPT}_{\rm LP}\).

Let \(C_{\max}=Q_0^3+2nP_{\rm art}R\), take \(t_0\) to be the largest power of two at most \(1/(1+C_{\max})\), and double until the first \(t_J\ge512m_+/\eta\). Take
\[
\begin{split}
 \alpha&=\big(100[J+\lceil\log_2(2RQ_0)\rceil+1]\big)^{-1},\\
 \rho&\le\min\{1,\eta/(100P_{\rm art}(n+1)),1/(100(J+1)t_JP_{\rm art}(n+1))\},
\end{split}
\]
with \(\rho\) the largest power of two at most that minimum. Take master size \(M=\Theta(N)\) rounded up to cover the root slots below, choose \(K=2^{\lfloor\log_2 M/a\rfloor}\), and check \(K\ge 2K_0\kappa_{\rm deg}\) from §5.3 and the root degree lift. It holds eventually for each sufficiently large fixed \(a\). Set \(T=\lceil 10^7(J+1)m_+K^2/\alpha\rceil\) as an accepted-step guard, and use the least power of two \(U_{\rm buf}\ge\max\{1,8T v_\ast/\rho\}\). The protected row units multiply physical balance by \(U_{\rm buf}\). A protected unit coordinate of price one and score zero in these units represents physical unit \(1/U_{\rm buf}\) in the direction chart.

Split each high-degree row permanently into a hub and a bounded-degree binary tree of slots for its real incidences, using \(O(N)\) slots in total and depth \(O(\log M)\). Null links have unit transfer and physical lift zero. Each slot, including hubs, has its own protected unit coordinate with the same original row image in physical units. Choose null price from \([1/(8\Xi_{\rm true}),1/(4\Xi_{\rm true})]\), where \(\Xi_{\rm true}\ge1\) bounds the actual root incidences in cap units divided by their smallest stale price. Routing each required incidence to the hub extends any old conserved comparator exactly, with cost at most \(\kappa_{\rm deg}\) times its true price, where \(\kappa_{\rm deg}=O(\log(2M))\) by this price. The entering chart envelope is at most \(8\Xi_{\rm true}\). The prices and images are positive dyadics on immutable lifetimes, with zero scores for null and protected columns. Physical reprices are at the same slots but are new lifetime IDs. Chart size including the padding caps and links remains \(O(N)\).

\hypertarget{descent-reports-and-recovery}{%
\subsection{Descent, reports, and recovery}\label{descent-reports-and-recovery}}

For each box slack \(s_j^-=X_j,\ s_j^+=u_j-X_j\) (with the artificial widths used when appropriate), set
\[
 \Psi(X)=\sum_{j,\pm}(s_j^\pm)^{-\alpha},\quad
 F_t=tC+\Psi,\quad E_t=F_t-t(C_*-P_{\rm art}n\rho).
\]
Suppose the reference trajectory currently has \(A_+X=b+\xi,\ \|\xi\|_\infty\le\rho\). Replace its artificials by positive and negative parts of \(b-Ax\) to get a point in the old exact slice. The cost increase is at most \(P_{\rm art}n\rho\) because the old sum already pays \(\|b+\xi-Ax\|_1\). Conversely take an optimal old comparator and replace its artificials relative to \(b+\xi\). Its new cost is at most \(C_*+P_{\rm art}n\rho\). Both replacements fit in \([0,R]\) by the residual bound and \(\rho\le1\). In particular \(C(X)\ge C_*-P_{\rm art}n\rho\) and \(E_t\ge\Psi(X)\).

Write \(w_j=\alpha[(s_j^-)^{-1-\alpha}+(s_j^+)^{-1-\alpha}]\) and \(g=t\nabla C+\nabla\Psi\). For the exact displacement \(d=Y-X\) to the new-slice comparator, use \(|\Delta s|\le2s+\Delta s\) for any slack displacement \(\Delta s\ge-s\). Summing gives \(\|w d\|_1\le2\alpha\Psi-\nabla\Psi^\top d\). Thus if \(t(C(X)-C_*-P_{\rm art}n\rho)\ge2\alpha\Psi\), the true optimal descent ratio on the box variables is at least one whenever the objective test is needed. To see this at equality with a zero comparator displacement, use a strict test, which ensures \(d\ne0\). We shall only require that a ratio below one implies \(t(C(X)-C_*-P_{\rm art}n\rho)\le2\alpha\Psi\).

Maintain root stale price and gradient words with
\[
       \tfrac12 w\le\widetilde w\le 2 w,\qquad |\widetilde g-g|\le \widetilde w/(1000K)
\tag{13}
\]
coordinatewise. Declare the slack floor \(\underline s=(16m_+)^{-1/\alpha}\), and choose \(\epsilon_x\) to be the largest power of two at most one sixty-fourth of the minimum of
\[
 \frac{\rho}{32(1+\|A_+\|_{\infty\to\infty})},\qquad
 \frac{\underline s}{10^8K},\qquad
 \frac{\eta}{1024(1+m_+)(1+\|A_+\|_{\infty\to\infty}+\|\nabla C\|_\infty)}.
\tag{14}
\]
The exponent \(1/\alpha\) is an integer \(O(B)\), so \(\underline s\) can be computed as an exact rational with bit height \(O(B\log M)\). At every phase start flush old maps, read and anchor all variables, and compute the stale words. At subsequent anchors evaluate at the read point. Take \(\ell_{\min}\) as the largest power of two at most \(\alpha/(16R)\), and \(\ell_{\max}\) as the least power of two at least \(16(1+\underline s^{-2})\). For slacks in \([\underline s/2,2R]\), these bounds include the rounded prices since \(s^\alpha\le2\) on this range and \(\alpha\le1\). Use the largest power of two at most \(\ell_{\min}/(10^6 K)\) as \(h_g\). Evaluate each term \(\alpha s^{-1-\alpha}\) to absolute error at most \(h_g/8\). For \(p_\alpha=1/\alpha\) and a nonnegative trial \(z\), compare \(z^{p_\alpha}s^{p_\alpha+1}\) with \(\alpha^{p_\alpha}\) to bisect a dyadic interval for the term. The interval \([0,\ell_{\max}]\) contains it. Sum the two terms for the price and use their signed difference plus the exact linear term for the score. Round the score to a nearest multiple of \(h_g\) using a fixed tie rule. Each combined numerical approximation has absolute error at most \(h_g\) at the read point. All queried magnitudes have logarithmic height \(O(B\log M)\). Integer powering and the bisection costs are fixed polynomials in bit height.

At an anchor write \(w^a\) for the true price there. The read changes each true anchor slack by at most \(1/(10^8 K)\) relatively by (14). For a relative slack change \(r\le1/10\), the mean value theorem and \(\alpha\le1\) bound the relative change of each weight term \(\alpha s^{-1-\alpha}\) by \(3r\). Thus the stored price at the anchor is at least \(w^a/2\). At settlement (11) gives displacement at most \((D_{\rm pub}+1)\theta_e\le\alpha/(10^8K\ell_e)\). Each anchor slack \(s\) satisfies \(s w^a\ge\alpha/2\), so the relative displacement from that slack is at most \(4/(10^8K)\). The errors from the read and the subsequent displacement in the price and barrier gradient are therefore each bounded in total by \(15 w^a/(10^8K)\). The linear gradient is unchanged within the phase. Adding \(h_g\) to these bounds proves (13), using \(h_g\le w^a/(10^6 K)\) and \(\ell_e\ge w^a/2\). Use the same price bounds in \(\Xi_{\rm true}\) and the score deduction of §5.3.

Accept a query only if its advertised descent score is at least \(1/(4K)\). If the true ratio is at least one, (13) gives stale ordinary ratio at least \(1/2-1/(1000K)\). Extend through degree links with zero protected motion and apply the ratio factor \(2K_0\). The check \(K\ge2K_0\kappa_{\rm deg}\) then forces acceptance. A rejection is a valid phase stop. For an accepted answer of certified price \(V>0\), take \(s\) to be the largest power of two at most \(v_\ast/V\), using integer magnitude comparisons. Then \(v_\ast/2\le sV\le v_\ast\). Use this scaled exact witness to define the next nominal step. Its real displacement \(h\) has \(\|\widetilde w h\|_1\le sV\), and (13) gives \(\|wh\|_1\le2v_\ast\). The nonreal coordinates have zero score. By the advertised score and (13),
\[
 -g^\top h\ge -\widetilde g^\top h-\frac{\|\widetilde w h\|_1}{1000K}
 \ge \frac{sV}{K}\left(\frac14-\frac1{1000}\right)
 \ge \frac{\alpha}{10^4 K^2}.
\]
The maximum relative slack change is at most \(2\|wh\|_1/\alpha<1/10\), from \(s_j^\pm w_j\ge\alpha/2\). For a slack \(b_s\) and its displacement \(\Delta b_s\), Taylor's remainder for \(b_s^{-\alpha}\) is at most \(2\alpha b_s^{-2-\alpha}(\Delta b_s)^2\) on this relative interval. Since \(b_s^\alpha\le2\), it is at most \(4(\alpha b_s^{-1-\alpha}|\Delta b_s|)^2/\alpha\). The sum of these term weights is \(\|wh\|_1\), so the total remainder is at most \(4\|wh\|_1^2/\alpha\le 16\alpha/(10^6K^2)\). Hence \(F_t\) decreases by at least \(\alpha/(2\cdot10^4 K^2)\). Lifting and compilation are exact on the enlarged chart. Omitting the protected coordinates changes real balance by at most \(v_\ast/U_{\rm buf}\) per step because their total absolute price bounds their physical contribution after all copies map to their original rows. Thus until the guard \(T\) the reference residual is at most \(\rho/8\), including every intermediate phase. This includes the cap-gain edges. The reanchoring and representation updates leave the reference trajectory fixed.

At a phase stop the preceding comparator gives \(t(C(X)-C_*+P_{\rm art}n\rho)\le2\alpha\Psi+2tP_{\rm art}n\rho\). At this same point \(E_{2t}=E_t+t(C(X)-C_*+P_{\rm art}n\rho)\). As \(\Psi\le E_t\), it follows that \(E_{2t}\le(1+2\alpha)E_t+2tP_{\rm art}n\rho\). At initialization the smallest slack is at least \(1/(2Q_0)\), so \(\Psi\le4m_+\) by \((2Q_0)^\alpha\le2\). Also \(C(X)\le C_{\max}\), \(C_*\ge-Q_0^3\), and \(t_0P_{\rm art}n\rho\le1/100\). Hence \(E_{t_0}\le7m_+\). The phase factors have product at most \(\exp(2\alpha J)\), and the sum of their additive terms is at most \(2(J+1)t_JP_{\rm art}n\rho\le1/50\). Together with the per-step decrease, this bounds every \(E_t\) by \((7m_++1/50)\exp(2\alpha J)\le16m_+\). In the induction an accepted step first preserves positive slacks by its relative change bound and the residual bound by its priced protected motion. The resulting bound \(E_t\ge\Psi\) then gives the declared slack floor for the next computation. Thus (12)--(14) and the Taylor estimate apply throughout the guarded prefix. Each phase has at most \(16m_+\cdot 2\cdot10^4 K^2/\alpha<T/(J+1)\) descents, since \(E_t\ge0\). Since \(J,1/\alpha=O(B)\), the global step bound is \(O(MK^2 B^2)\). Each chart epoch uses the structural counts in §4.

At the last phase stop \(C(X)\le C_*+P_{\rm art}n\rho+16m_+/t_J\le C_*+\eta/8\). Since \(C_*\le\mathrm{OPT}_{\rm LP}\le Q_0^3\) and the original cost is at least \(-Q_0^3\), the total artificial value is at most \((2Q_0^3+1)/P_{\rm art}\le\eta/8\). Omitting these artificials increases the balance residual by at most their total value, to at most \(\rho/8+\eta/8<\eta/4\), and leaves original cost at most \(\mathrm{OPT}_{\rm LP}+\eta/8\). Read all original coordinates with absolute accuracy (14) and project into their dyadic boxes. Projection cannot increase their distance to the nominal box point. Round downward on the largest power-of-two mesh at most \(\eta/[64(1+m_+)(1+\|A_+\|_{\infty\to\infty}+\|\nabla C\|_\infty)]\). The projection and rounding keep every value in its exact box. The coordinate errors from reading and rounding together are at most \(\eta/[32(1+m_+)(1+\|A_+\|_{\infty\to\infty}+\|\nabla C\|_\infty)]\), by (14). Multiplying by the maximum absolute row sum or by \(m_+\|\nabla C\|_\infty\), respectively, gives extra errors at most \(\eta/32\) for the balance and cost. This produces the exact-box dyadic additive answer for §1. All reads and final flushes concern the same reference trajectory.

\hypertarget{work-guards-and-the-uniform-algorithm}{%
\section{Work, guards, and the uniform algorithm}\label{work-guards-and-the-uniform-algorithm}}

\hypertarget{the-structural-recursion}{%
\subsection{The structural recursion}\label{the-structural-recursion}}

For each fixed sufficiently large integer \(a\), take \(m_0=M\) a power of two of size \(\Theta(N)\), large enough for the initial degree and cap padding. Set \(u_i=\log_2 m_i\) and
\[
 u_{i+1}=\left\lceil(1-1/(3a))u_i\right\rceil .
\]
Stop at the first \(m_d\le M^{2/a}\), tested as \(a u_d\le2u_0\), and use the exact bottom query of §5.3. At nonbottom levels take \(k_i=2^{\lfloor u_i/a\rfloor}\) and \(b_i=2^{\lfloor u_i/a^2\rfloor}\). Suspend a branch if the schedule fails to decrease at a nonbottom level, an inner spacing base is less than two, or the chart, comparison, or recursion capacity guards fail. For fixed \(a\) and large \(M\) the depth is \(O(a\log a)\). The geometric sum gives \(\sum_i\log m_i\le(3a+o_a(1))\log M\), since the ceiling contributes only \(O_a(1)\) to the sum before this stop. An epoch of level \(i\) starts with \(O((m_i/k_i)\log^{O_a(1)}M)\) child records. Even the larger total slot and lifetime issuance bound \(O((m_i/k_i)b_i\log^{O_a(1)}M)\) is below \(m_{i+1}/4\) eventually, since \(1/a-1/a^2>1/(3a)\). Its active bound includes the isolated-cap placeholders and deletion-first slack. Reserve capacity for chart and public records and, where needed, temporary slots proportional to these bounds. Allocate and index empty placeholders only within the charged capacity. The repair list dimension and the number of distinct resources fit \(M^2\) eventually by \(m_i k_i^{O(1)}\log^{O_a(1)}M\) and a sufficiently large fixed \(a\). These are structural checks. The temporary graph and numeric guards of §2.4 count only toward local work and storage.

Here is why the size guards do not require assuming an answer about the oracle or solving for \(H\). Root and emitted chart dyadics have height
\[
                    O_a((1+\widehat B)\log M+\log^{O_a(1)} M)
\tag{17}
\]
by the envelope, score meshes and the parameters of the barrier. Frozen gauge sweeps use \(O(k^3(1+\log Z))\) affine operations along the lineage of a stored value. Each operation adds \(O((1+\widehat B)\log M+\log^{O_a(1)}M)\) to its bit height. A stable-ball LP vertex has height at most a fixed polynomial in its dimension times this bound, by clearing denominators and using minors. The same additive-height rules for products and common denominators apply to frozen path and calendar summaries. Thus one can bound the heights of all immutable local summaries by
\[
 F_i=(1+\widehat B)^2(k_i+1)^{C_3}\log^{O_a(1)}(2M)
\tag{18}
\]
for an absolute \(C_3\). Coefficients in the repair lists have only logarithmic-length macro lists as input and output. In clipping, compute the prefix product and subpath price as summaries of immutable chart words before scaling by the inherited coefficient. This adds the immutable summary height to the old coefficient height. It does not multiply the old height by the number of physical edges on a subpath. A conformal split or a selector replacement then multiplies old height plus \(F_i\) by a fixed polynomial in the macro count and \(H\), as follows from determinant heights and products of at most \(O(H)\) current factors. There are \(O_a(1)\) snapshot ancestors. For any declared polylogarithmic \(H\) in §2.4, their product of these list factors is also a logarithmic power for fixed \(a\). New literal chart words use (17). This proves
\[
 \beta\le (1+\widehat B)^2(k_i+1)^{C_4}\log^{O_a(1)}(2M),
\tag{19}
\]
with absolute \(C_4\), on all prefixes that meet the structural prerequisites. It is sufficient to use \(M^5\) as a bit-height guard for \(\beta\). For all large \(M\) at fixed \(a\ge C_4\) and \(\widehat B\le N=O(M)\), (19) is below that guard. In particular a logarithmic exponent in (19) can depend on the declared \(e(a)\). This does not affect the fixed polynomial size envelope.

The original grid has \(G_0=O((1+\beta)(H+1)^2)\) states per row. Each sampled name generates \(O(G_0)\) links. The state counts satisfy \(n'=rG_0\) and \(m'=O(|S|G_0)\), where \(S\) is the stage sample, including all initial state labels. Under the guards on \(H^2\) and \(\beta\), and \(|S|\le M^2\), \(n'+m'=O(M^8)\). The cover's \(W_C\) is at most \(n'+m'\), and its reference \(\widehat W\) is at most \(2^{2a}\max(2^{16},16W_C)\). Both values are below \(X_*\) eventually for fixed \(a\). This proves adequacy of the guard used to choose the path bound. Small or late branches may suspend on their exact guard tests, as elsewhere in the uniform schedule. Parameter construction itself can stop at a guard using capped integer powering.

The larger \(H\) changes the price, occurrence and list-length estimates by logarithmic powers. Equations (3)--(4) and the token bound keep the primary support at \(O(r(1+J\log(2M))\log(2M))\) names. The unary completion adds \(O(r\log(2M))\) pivot groups and replaces each demand by at most \(2H+1\) original factors at price factor \(1+O(HL)\). Each cut repair scans these actual lists. Across the fixed number of snapshot ancestors the extra factors are absorbed in \(\mathcal L\). Selected names promote only actual chart endpoints, so §§4.3--4.4 still give \(O((m_i/k_i)b_i\mathcal L)\) issued child records and the same degree bounds. The stretch in (10) and the score factor \(2K_0\) remain logarithmic powers, giving \(K\ge2K_0\kappa_{\rm deg}\) eventually. The gate's \(M^{12c_1/a}\) factor and the polynomial in \(\beta\) in (5) are local work.

Use \(F=1+\widehat B\) in the work analysis of a call. Let \(P(F)\) be one fixed polynomial dominating all bit-work factors for a local operation, and let \(L_{\rm loc}=M^{C/a+o_a(1)}\) dominate the remaining local per-record work for an absolute \(C\). Equation (19) gives a gate factor \(\beta^{p}\le F^{2p}(k_i+1)^{C_4p}\log^{O_a(1)}(2M)\) for the absolute exponent \(p\) in (5). There are \(O(\log(2M))\) successful selector stages per batch and \(O_a(\log(2+M+\widehat B))\) attempts per stage or other stochastic call on a fixed branch. This also counts every aborted trial and its temporary records. A rejected draw publishes no chart, port or public-forest records. Charge validation and disposal with the attempted build. All these factors multiply local work only. Generic public logs incur dictionary and forwarding work from §5.2 with numeric descriptors passed by reference. Chart arithmetic and rational LP work have absolute polynomial degrees. In (12), \(\log H_\ast=O(\log(2+M+\widehat B))\), \(\log(2+Tv_\ast/\ell_{\min})=O_a(F+\log(2M))\), and \(\log(2+1/\epsilon_x)=O_a(F\log(2M)+\log^{O_a(1)}(2M))\). Hence tracker word and primitive costs also contribute a fixed power of \(F\). Integer identifiers for graph and dictionary instructions have at most logarithmic-power bit height. The work estimate uses the normalized \(F\) even when \(B\) is much larger.

For a prefix of \(s\) operations at level \(i\), counting edits and queries, write its bit work bound as \(m_i A_i+s U_i\). Include an operation partly attempted and then suspended. Initialization includes the frozen data, all-mark selection, rebuilding pools, and the initialized child. Local window work is charged by §4.3. The generic public-log recurrence separately has the same bounds with the bit-price factor omitted. A prefix pays full epoch initializations \(O(1+s k_i/m_i)\) times. Each abstract child edit is charged at the parent to §4.3 and its pool work. Each forwarded child log has only the separate logarithmic charge from §5.2. There is one child query per parent query. Consequently one can take
\[
\begin{aligned}
 A_i&\le L_{\rm loc} P(F)+(m_{i+1}/m_i)\log^{O_a(1)}M\,A_{i+1},\\
 U_i&\le L_{\rm loc} P(F)+O(k_i)A_i+b_i\log^{O_a(1)}M\,U_{i+1}.
\end{aligned}
\tag{15}
\]
The child term includes its work and the processing of its public logs. Taking one total upper coefficient for both is valid since each has the separate logarithmic transfer charge. Destruction of arenas and old map flushing enumerate charged current records or charge prior allocations. The bottom coefficients are bounded by \(L_{\rm loc}P(F)\), using the rational LP on \(m_d\le M^{2/a}\) records for a query and the absolute polynomial bit bound for explicit rational optimization. The first multiplier in (15) is at most \(2m_i^{-1/(3a)}\log^{O_a(1)}M=o_a(1)\), since \(m_i>M^{2/a}\) at every nonbottom level. This gives \(A_i\le2L_{\rm loc}P(F)\) eventually at every level. For the second line, \(k_i=O(M^{1/a})\) and \(\prod_i b_i\le M^{3/a+o_a(1)}\), while products of structural logarithmic factors have bounded exponent for fixed \(a\). Substitution yields \(M^{1+O(1/a)+o_a(1)}P(F)\) initialization and \(M^{O(1/a)+o_a(1)}P(F)\) amortized prefix work per operation. Since \(\widehat B\le N=O(M)\), any logarithmic power of \(2+M+\widehat B\) in this case is a logarithmic power of \(2+M\). The per-handle work \(k_i^{O(1)}\) on a true changed segment is part of the local term. The prefix principle also applies to the explicit partial states of the old-to-final transactions. \href{https://raw.githubusercontent.com/openai/math/adc7f1241b42e322a6451854ab7e4b4c146bf78a/preprints/Almost-Linear-Time-Maximum-Cardinality-Matching-in-Sparse-General-Graphs-September-24-2026/main.pdf\#page=73}{OpenAI, Lemma 8.8 and its prefix recurrence}

\hypertarget{total-phases-structural-guards-and-conditional-probability}{%
\subsection{Total phases, structural guards and conditional probability}\label{total-phases-structural-guards-and-conditional-probability}}

There are \(J+1=O(F)\) phase starts with full initialization and tracker sweep, and \(T=O(MK^2 F^2)\) steps. The tracker charges at most \(T\log^{O_a(1)}M\) root reports and elementary root reprices. There are at most \(T+J+1\) root queries, including the rejecting query at each phase stop. Each query reports at most \(M^{2/a}\log^{O_a(1)}M\) pieces. Let \(U_{\rm log}\) be the root log count bound per operation obtained by the structural recurrence accompanying (15), and let \(P_{\rm tr}(F)\) be the tracker word factor. A work bound including all phase starts, root operations and tracker costs is
\[
 O\!\left(Fm_0 A_0+(T+F)\log^{O_a(1)}M\,U_0+
 \big[FM+(T+F)(M^{2/a}+U_{\rm log})\big]M^{C_{\rm tr}/a}\log^{O_a(1)}M\,P_{\rm tr}(F)\right).
\]
The factor with absolute constant \(C_{\rm tr}\) pays for arithmetic on the explicit exact words at the interface. These include choosing \(s\) by comparing powers of two with \(v_\ast/V\), multiplying each exact signed start or individual coefficient by \(s\), and converting the result outward to the mantissa format of §6.2. They also include producing transfer and coefficient intervals from an exact descriptor when it first enters the tracker. At the bottom the LP has at most \(2M^{2/a}\) columns. Clearing denominators in its inequalities and taking minors bounds vertex and certificate heights by a fixed polynomial in that dimension times the entering word height (17). At a nonbottom level (18)--(19) give \(F^{O(1)}(k_i+1)^{O(1)}\log^{O_a(1)}M\) bits for each exact local descriptor, start, or certificate. An inherited path start is passed unchanged. For an individual lift a bounded-arity expansion multiplies by local exact factors, which adds their bit heights along each fixed-depth chain. Converting a rational word to a directed float uses integer division after a power-of-two shift determined by integer lengths and comparisons. Therefore each conversion, including the dyadic step calculation, costs at most \(M^{C_{\rm tr}/a}\log^{O_a(1)}M\,P_{\rm tr}(F)\), with a fixed polynomial \(P_{\rm tr}\). The same bound includes the old-map flushes, final recovery and point services. The splay bound in §6.1 adds only logarithmic structural amortization per piece or log. Since \(F\le1+O(B)\) and \(K=O(M^{1/a})\), the additive algorithm on each adequate branch costs at most
\[
       M^{1+C'/a+o_a(1)}(1+B)^{C''},
\tag{16}
\]
with absolute constants \(C',C''\). In particular for \(B=N^\epsilon\) the \(B\)-contribution to the exponent is at most \(C''\epsilon\) for every \(\epsilon>0\). Each assembled call with \(\widehat B>N\) uses the separate exact-LP bound from §1. Neither the number of vertex-scale bands in the forest cover nor the number of numerical states in a local path gate appears in the multiplier of any descendant term in (15). The fixed-depth logarithmic structural factors are functions of \(M\) only.

Use explicit two sorts of guards. Structural capacities, degree and occurrence bounds, allowed snapshot support size, and all input-name lists are tested in terms of \(M,k_i,b_i,\log M\) and their proven logarithmic factors. A structural guard is separate from the numerical grids and word arrays. For the whole run of a branch take \(H_\ast=(2+M+\widehat B)^{C_{\rm guard}}\) for an absolute sufficiently large exponent, where \(\widehat B=O(B)\) is the actual word bound declared in §1. Use a power of the same quantity for finite exponent magnitude, allocated bits and LP nodes where needed. Count stochastic calls and attempted arithmetic and tree operations before a next operation could exceed the corresponding guard. The mantissa choice (12) then uses \(\log H_\ast=O(\log(2+M+\widehat B))\). Under the guard its correctness proof and the descent induction apply to every prior prefix. Applying (15), the bit bounds, the notification bound and the step estimate to the prefix up to a purported first guard violation, including the attempted operation, gives values below these guards for all sufficiently large \(M\) at fixed adequate \(a\), uniformly in \(\widehat B\). Increasing the absolute guard exponent affects the mantissa only through a constant factor on its logarithm, so the argument is noncircular. Comparisons, emission counts and exact scope checks can suspend an inadequate branch.

For each fixed history at the start of a stochastic routine, a constant probability attempt succeeds for the padded cover, priority table, each selector stage, and the unary pivot verification. Trials have explicit dimension and fair-bit caps. Their accepted output properties have been checked exactly. On a branch of index \(j\) cap fresh attempts per stochastic call by \(C(j+\log(2+M+\widehat B))\) with the constant large enough to give probability at most
\[
     2^{-j}(2+M+\widehat B)^{-C_{\rm guard}-c-10}
\]
of exhaustion, conditional on its assembled input. The exponent \(c\) is any fixed desired confidence. In particular a large \(\widehat B\) affects trial work through its logarithm, while the accepted tables, physical copy counts, and label budget use the original structural \(M\). Each rejected trial is counted. A union bound over all guarded calls in all phases bounds detected stochastic failure on branch \(j\) by \(2^{-j}M^{-c-5}\). The same union bound for sequentially assembled wrapper inputs bounds failure over \(O(\log N)\) calls. Adaptivity is permitted here since the trials are independent of the past conditional on the actual call input. The common priority table remains valid through all future edits by its deterministic tests, so no additional union over possible edit sequences is needed. \href{https://raw.githubusercontent.com/openai/math/adc7f1241b42e322a6451854ab7e4b4c146bf78a/preprints/Almost-Linear-Time-Maximum-Cardinality-Matching-in-Sparse-General-Graphs-September-24-2026/main.pdf\#page=74}{OpenAI, Lemma 8.9 for the guarded-prefix method}

\hypertarget{uniformity-and-total-sparse-time}{%
\subsection{Uniformity and total sparse time}\label{uniformity-and-total-sparse-time}}

For each call with \(\widehat B\le N\), run a deterministic rational-LP fallback and the branches for successive fixed integer \(a\) with geometric instruction sharing. For example in round \(t\) give branch \(j\le t\) up to \(2^{t-j}\) charged instructions, resuming its arrays. Let the first successful completion return, and suspend branches on a failed guard or exhausted trials. The fallback solves the feasible dyadic LP exactly and rounds inside the box to additive tolerance, so the schedule is total. A fixed-index branch finishing after \(W_j\) instructions finishes under the schedule within \(O(2^j W_j)\) total issued instructions, up to logarithmic control and copying costs. This follows by summing the geometric budgets up to the round when it receives \(W_j\). Use fair independent random bits and charged digit arrays for large words so big operations can be interrupted between bit operations. Large allocations and control instructions consume their usual charged time. Input pointers and the normalization are shared among branches. A branch can only finish successfully after its certificates and guards have passed. For each fixed adequate \(a\), (16) and the probability bound apply to its own run uniformly in the numerical precision. With high probability the whole schedule is at most the time until this branch succeeds. Taking arbitrarily large fixed \(a\) in the bound gives one uniform \(N^{1+o(1)}(1+B)^{C'''}\) high-probability envelope with a fixed \(C'''\), including logarithmic control costs and the wrapper's \(O(\log N)\) calls. On the rare event of an exhausted stochastic branch the fallback also bounds overall running time by a polynomial, without invalidating any returned box vector. This uniformity argument uses the instruction sharing idea of OpenAI §9.4 and proves the bounds with the additional arbitrary-precision factors. \href{https://raw.githubusercontent.com/openai/math/adc7f1241b42e322a6451854ab7e4b4c146bf78a/preprints/Almost-Linear-Time-Maximum-Cardinality-Matching-in-Sparse-General-Graphs-September-24-2026/main.pdf\#page=78}{OpenAI, Lemmas 9.4--9.5 and §9.4}

These estimates and the exact-LP branch for large precision give the stated sparse time in terms of \(m+n\). Section 1 uses \(B=O(\log N+\log(2+NW/\delta))\), with the extra \(\log N\) factors absorbed in \(N^{o(1)}\). The same input-access model works for negative or zero costs, near-unit gains, arbitrary capacity scales, self-loops, and \(\delta>1\). The returned affine representation has exact capacities and the true unrounded imbalance and cost bounds derived in §1.

\hypertarget{exact-primitive-accounting}{%
\subsection{Exact primitive accounting}\label{exact-primitive-accounting}}

\begin{itemize}
\tightlist
\item
  \textbf{Rational optimization.} Grötschel, Lovász and Schrijver, Theorem 6.4.12, solves an explicitly listed rational linear program in polynomial bit time. All our equalities can be written as paired inequalities. Following Remark 6.5.2, find an optimum vertex by fixing the attained optimum and then lexicographically optimizing individual coordinates on the successive faces, using at most dimension plus one elementary LP calls. All such faces are of the original bounded polyhedron, so their attained coordinate values have the original determinant height bounds. In §3.3 there is at most one such optimization on \(O(Dk)\) columns per frozen stable row, with feasibility from its true cap and compactness from the price sublevel. At the bottom there is one feasible bounded signed LP per reached query with at most \(2m_d\) variables, feasible at zero and bounded by the sum constraint. The high-precision branch and the scheduled fallback use the entire bounded dyadic input and exact feasibility from §1. Total costs and bit heights are the polynomials charged in (15)--(16) and the fallback's own instruction share. These optimization calls have explicit rational inputs. \href{https://www.zib.de/userpage/groetschel/pubnew/paper/groetschellovaszschrijver1988.pdf\#page=192}{GLS, Theorem 6.4.12 and Remark 6.5.2}
\item
  \textbf{Dynamic distances and forests.} Section 2.4 uses the short-threshold portion of Chuzhoy Theorem 1.2, with the unit-length degree-two subdivision and the Theorem 3.3 bound detailed in §9. Each accepted sample is built once and its links undergo only deletions. Section 9 checks all invoked graph and pruning routines in the finite-word model. Section 6.1 uses the amortization of OpenAI Lemma 5.3 on the subdivided forests, with the directed fields and bounded arithmetic in §6.2. \href{https://arxiv.org/pdf/2109.05621v1}{Chuzhoy v1, Theorems 1.2 and 3.3} · \href{https://raw.githubusercontent.com/openai/math/adc7f1241b42e322a6451854ab7e4b4c146bf78a/preprints/Almost-Linear-Time-Maximum-Cardinality-Matching-in-Sparse-General-Graphs-September-24-2026/main.pdf\#page=37}{OpenAI, Lemma 5.3}
\item
  \textbf{Proof arguments from the matching preprint.} Sections 3--7 derive the needed versions of the integer padding in Lemma 7.2, portal-grid Lemma 7.6, clock and compilation Lemmas 8.3--8.4, public lifting and recourse Lemmas 8.5--8.7, and the prefix and schedule arguments of Lemmas 8.8--8.9 and §9.4. The resource selection here is §2, with its own finite lists, verified capacity guards and unit-length oracle. Its use on the two actual chart inputs is established in §4. The disjoint integral resource alphabet and master divisibility of the ordinary helper in matching Lemma 6.1 are specific to that helper. \href{https://raw.githubusercontent.com/openai/math/adc7f1241b42e322a6451854ab7e4b4c146bf78a/preprints/Almost-Linear-Time-Maximum-Cardinality-Matching-in-Sparse-General-Graphs-September-24-2026/main.pdf\#page=41}{OpenAI, specified preprint, Lemma 6.1}
\end{itemize}

The moderately dense two-sparse LP theorem of Jiang, Kapralov, Li and Sidford has an \(n^{3/2}\) term and additive output errors. Its Lemma 8.14 gives a diagonal initialization. Van den Brand et al.~Theorem 5.3 gives an adaptive randomized minimum-ratio cut structure for an ordinary incidence matrix, and Appendix B reduces to sparse transshipment. Daitch and Spielman Lemma 3.2 uses least-lossy-path balancing and shrinkage in a source--sink problem, Theorem 4.5 analyzes the \(M\)-matrix solver, and Appendix C supplies the interior-point analysis. These results are prior context. \href{https://arxiv.org/abs/2510.17740v1}{JKLS, Theorem 1.2 and §8} · \href{https://arxiv.org/abs/2407.10830v1}{van den Brand et al., Theorem 5.3 and Appendix B} · \href{https://arxiv.org/abs/0803.0988v2}{Daitch--Spielman, §§3--4 and Appendix C}

\hypertarget{reproducible-exact-checks}{%
\section{Reproducible exact checks}\label{reproducible-exact-checks}}

The following Python uses only rational arithmetic. It checks the rounded levels in §5.4, the nonzero physical compensation on the surviving segment, sign-aware rounding and barrier constants, and representative subdivision, radius and grid inequalities for the cover in §9.

\begin{lstlisting}[language=Python]
from fractions import Fraction as F
from collections import defaultdict
from random import Random

old_price = 2**30
row_scale = 2**40
gain = F(3, 4) + F(1, 2**64)
mesh = F(1, 128)

def round_down(amount):
    return (amount // mesh) * mesh

assert F(row_scale, 2 * old_price) == 512
old_head = row_scale * gain / (4 * old_price)
assert old_head == 192 + F(1, 2**56)
assert round_down(old_head) == 192
assert old_head * (F(192) / old_head) == 192
surviving_head = row_scale * gain / 2
assert round_down(surviving_head) == F(3 * row_scale, 8)
assert surviving_head - round_down(surviving_head) == F(1, 2**25)

child_witness = [
    (-row_scale, F(row_scale, 2), 0),
    (0, -F(row_scale, 2), round_down(surviving_head)),
    (row_scale, 0, 0),
    (0, 0, -round_down(surviving_head)),
]
assert not any(map(sum, zip(*child_witness)))
physical_witness = [
    child_witness[0],
    (0, -F(row_scale, 2), surviving_head),
    child_witness[2],
    (0, 0, -surviving_head),
]
assert not any(map(sum, zip(*physical_witness)))

for exponent in (120, 4096):
    defect = F(1, 2**exponent)
    cycle_transfer = 1 - defect
    unit_price = 3 / (1 - cycle_transfer)
    assert defect * unit_price == 3

routes = [
    {"copy0-tail": F(3, 7), "copy1-tail": F(2, 7)},
    {"copy1-tail": F(5, 11), "copy0-mid": F(1, 11)},
    {"copy0-mid": F(8, 13), "copy2-head": F(1, 13)},
]
images = [
    (-1, F(1, 2), 0),
    (0, -1, F(3, 4)),
    (-1, 0, F(3, 8)),
]
direction = (F(1), F(1, 2), F(-1))
assert not any(sum(coef * column[v] for coef, column in zip(direction, images))
               for v in range(3))
literal = {f"edge {i}": x for i, x in enumerate(direction)}
expected_raw_sum = defaultdict(F, literal)
for x, column in zip(direction, images):
    for incidence, route in zip(column, routes):
        for name, amount in route.items():
            expected_raw_sum[name] += x * incidence * amount
assert {name: val for name, val in expected_raw_sum.items() if val} == literal
hole_copies = {"copy0-tail", "copy1-tail", "edge 1"}
assert {
    name: val for name, val in expected_raw_sum.items() if val and name not in hole_copies
} == {name: val for name, val in literal.items() if name not in hole_copies}

rng = Random(20261120)
for _ in range(2000):
    lower = F(rng.randrange(-20, 21), 32)
    actual = lower + F(rng.randrange(128), 128)
    scale_low = F(rng.randrange(1, 8), 16)
    scale_high = scale_low + F(rng.randrange(8), 32)
    scale = (scale_low + scale_high) / 2
    div_bound = lower / (scale_low if lower < 0 else scale_high)
    mul_bound = lower * (scale_high if lower < 0 else scale_low)
    assert div_bound <= actual / scale
    assert mul_bound <= actual * scale

for variables in (5, 11, 128, 1024):
    for input_bits in (32, 128, 512, 2048):
        alpha = F(1, 100 * (2 * input_bits + 21))
        for ratio_parameter in (4, 16, 256):
            budget = alpha / (1000 * ratio_parameter)
            descent = alpha / (10000 * ratio_parameter**2)
            remainder = 4 * (2 * budget)**2 / alpha
            progress = alpha / (20000 * ratio_parameter**2)
            assert descent - remainder >= progress
            assert 16 * variables / progress < (
                10**7 * variables * ratio_parameter**2 / alpha
            )

hop_lag = 10 * (2 + 1)**2
gain = 1 + F(1, 4096)
long_hops = 16384
assert gain**(-hop_lag) > F(7, 8)
assert F(long_hops, 2**30) + gain**(-long_hops) < F(7, 8)

radius = 14
log_reference = 16
for distance in range(1, radius + 1):
    assert 2 * distance + 2 <= 3 * radius
largest_scale = 1 << (2048 * radius * log_reference**4 - 1).bit_length()
child_radius = max(8 * radius, largest_scale)
assert largest_scale <= child_radius <= 4096 * radius * log_reference**4
assert 4 * 1024 * radius * log_reference**4 < 1024 * child_radius * log_reference**4
path_budget = 47
Q = 1 << (10**5 * path_budget - 1).bit_length()
assert (1 + F(1, Q))**(4 * path_budget + 1) < 2

for shift in range(2, 9):
    phi = F(1, 2**shift)
    for m in (2, 16, 121):
        b = max(2, (2*m-1).bit_length())
        budget = 128*b / phi
        assert (1 + phi/64)**int(budget) > 2*m
    assert F(15, 4)/phi > F(18, 5)/phi
    assert F(16, 15)*F(15, 2)/phi == 8/phi
\end{lstlisting}

\hypertarget{finite-word-short-threshold-neighborhood-cover}{%
\section{Finite-word short-threshold neighborhood cover}\label{finite-word-short-threshold-neighborhood-cover}}

\hypertarget{the-short-threshold-cover}{%
\subsection{The short-threshold cover}\label{the-short-threshold-cover}}

Consider a unit-length undirected graph \(G_s\), with fixed labels and only online edge deletions after initialization. The target radius is an integer \(L\ge2\). Split the initial graph into its connected components. Set aside initially isolated labels and answer queries between distinct initial components negatively. For each other component \(G_C\), let \(W_C=|V(G_C)|+|E(G_C)|\). Use the following implementation with parameter \(\varepsilon=1/a\).

\begin{itemize}
\tightlist
\item
  Each state vertex \(x\) is a supernode. Subdivide each edge with a regular vertex, and add a regular leaf \(x'\) at every state vertex \(x\). Every auxiliary edge has length one. Use target radius \(D_C=3L\). Each state-edge deletion deletes both links at its subdivision vertex, then deletes that isolated regular vertex, with a query fence. This is a valid input structure with dynamic degree bound two. The number of regular vertices initially is \(W_C\). The maintained covering cluster for \(x'\) contains \(y'\) whenever the current state distance is at most \(L\), since the auxiliary distance is then at most \(2L+2\le3L\).
\item
  Store the \texttt{CoveringCluster} pointer for each regular leaf. A cluster's ordered membership list from the RecDynNC interface can use a balanced search tree with lookup and update overhead \(O(\log(2+W_C))\) per affected membership. A positive test issues one \texttt{short-path-query} in that cluster. Take a simple auxiliary path by loop erasure if the returned object is a walk, discard the end leaf links and suppress the subdivision vertices. This gives a walk of live state edges, of length at most \(3L(\log(2+W_C))^{2^{O(a)}}\). The recursive unfolding and conversion use \(L\log^{O_a(1)}(2+W_C)\) graph operations. In the selector \(L=O(\log M)\).
\item
  When \(\lfloor\log_2\lfloor\log_2 W_C\rfloor\rfloor<C_{\rm sm}a\), use ordinary deterministic BFS in that initial component at query time, stopping at radius \(L\). Here \(C_{\rm sm}\) is a fixed constant exceeding the admissibility constant for Theorem 3.3. This test uses bit lengths. The component has size bounded in terms of \(a\) only. The same small-size convention inside the cover uses AlgSlow for a MaintainCluster call whose initial regular count is too small for the required parameter. Its quadratic regular-count factor then costs only a constant depending on \(a\). Singleton and empty graphs have their direct answers. If \(D\) is large relative to the reference size, use the large-\(D\) AlgSlow case of §8. These cases cover the small internal instances without assuming a large component or a large radius at every call.
\end{itemize}

For a large component, Theorem 3.3 gives total initialization and update work
\[
 O_a\!\left(W_C^{1+c_1/a}L^3\log^{c_2(a)}(2+W_C)\right)
\tag{A}
\]
with an absolute \(c_1\). This includes the first cover construction and any flag-lowering splits before the first query. The \(D_C^3\) factor is acceptable since \(L=O(\log M)\). The gate uses just one radius. The minimum-over-scales distance interface of §3.4.2 and the general large-radius reduction in Appendix C, including its dynamic minimum spanning forest routine, are unnecessary.

\hypertarget{topology-radii-and-thresholds}{%
\subsection{Topology, radii and thresholds}\label{topology-radii-and-thresholds}}

Use the algorithms in Chuzhoy §§3.1--3.3, 4--9 and Appendices A, B, D and E, with the following conventions. A current graph length is integral. Induced subgraphs, all edge copies under supernode splitting, all memberships and all witness embeddings have explicit identifiers. Graph changes use adjacency lists, balanced dictionaries, heaps and lists. Keep adjacency in length order for the output-sensitive ball-growing searches of §§4.2 and 8.3. A settled vertex's temporary cursor waits at its cheapest outward incidence. Advance it when that incidence enters the settled set, using inverse lists of pending incidences. An explored prefix costs one initial outward pointer per settled vertex and at most one visit to each internal incidence, up to heap and dictionary overhead. This avoids scanning all edges of a high-degree supernode upon discovering it. In ProcCut' run the two searches with equal work credits and stop at the first eligible split. All edges removed from the old cluster meet its inner ball and, because their lengths are at most \(D\), lie in the explored larger ball. Thus short searches and emitted splits charge \(O(1+|V(C')|+|E(C')|)\) operations with logarithmic overhead. A generalized Even--Shiloach tree uses integral distance keys \(\lambda(v)+\ell(v,w)\), integer levels, and a sentinel for infinity. Fixed edge lengths and all rational parameters come from the particular invocation. Addition, comparison, minimum, integer product, and quotient with remainder suffice.

The regular-vertex weights are one and the supernode weights are zero. The code keeps membership counts and their integer dyadic classes. For ProcCut put \(b=\max(16,\lceil\log_2 W\rceil)\), where \(W\ge1\) is the initial number of regular vertices in that cover. Use \(64b^2\) in the C2--C3 layer tests, and \(b^4\) in their radius envelopes. The first layer includes the center at distance zero. Choose regular centers, with isolated regular vertices eligible, and drop isolated supernodes. An exhausted search has empty further layers. Take class cutoff \(r=2\lceil b/\lfloor\log_2 b\rfloor\rceil\). The D.1 growth proof compares \(W\le 2^b\) with \((1+1/(64b^2))^{128b^3-2}>2^b\). For example, \((1+1/u)^u>2\) at the integer \(u=64b^2\) and \(128b^3-2\ge bu\). The \(r+1<b\) ineligibility types leave an eligible layer before \(128b^4\) if the ball of radius \(256Db^4\) has at most half the cluster weight. At the outer cover and every scale, \(b\le b_*+1\), so \(2\cdot256Db^4<1024Db_*^4\). Here \(D\) is the split's own scale. This ensures the two explored balls of a flag pair are disjoint and one has weight at most half.

For the potential of Chuzhoy §4.1 and Appendix D.3, let a live membership of a regular vertex in cluster \(C\) contribute \(1+\log_2 w(C)/b^2\), where \(w(C)\) is its current live regular weight. The logarithm is evaluated only for a membership, so \(w(C)\ge1\). At global deletion freeze that vertex's membership count and each of its membership budgets in the analysis. Such a deleted vertex still contributes these frozen budgets to its dyadic class and never participates in a live cluster size, split, or layer test. The algorithm need only store the live records.

Each regular vertex moved to the new cluster off the overlap saves at least \(1/b^2\) by the halving condition. Each live vertex copied in the overlap adds at most two to the sum. For the total budget, and separately for the vertices already in any class of membership count at least \(2^j\), C2 and C3 bound the number copied by the corresponding number moved divided by \(64b^2\). Thus the budget for any such fixed old set cannot increase at a split. Frozen budgets are unchanged, and the live budgets can also only decrease under graph deletions. Initially the total is at most \(W(1+1/b)\), and it stays at least \(W\) since every live or deleted regular vertex accounts for at least one membership. Each split producing \(C'\) uses at least a constant times \(w(C')/b^2\) of the decrease, so the total initial regular weight of all clusters, including the original, is \(O(Wb)\).

For the multiplicity bound, let \(S_{\ge j}\) contain all vertices whose live or frozen count is at least \(2^j\). Counts grow by at most one per split and stay fixed after global deletion. A vertex entering \(S_{\ge j}\) contributes at most \(2^j(1+1/b)\) at that time. For \(j=0\) charge that bound at initialization. No other step increases the aggregate for the old set. If \(p_j=|S_{\ge j}|\), its budget is at least \(2^j(p_j+p_{j+1})\) and at most \(2^j(1+1/b)p_j\). Hence \(p_{j+1}\le p_j/b\) and \(p_j\le W/b^j\). In particular \(S_{\ge r}\) is empty. This gives the \(W^{O(1/\log\log W)}\) membership bound of Claims 4.8 and 4.12 with the bounded-size convention. Counting clusters a vertex ever enters gives the same bound: a split that moves a vertex out of the old cluster into the new one without copying it can occur along a membership lineage at most \(\log_2 W\) times, because the live weight at least halves at each such move, and the number of lineages can grow only when its counted memberships are copied. The extra factor \(1+\log_2 W\) is absorbed in \(W^{O(1/\log\log W)}\).

Use the least power of two \(\widehat W\) with exponent divisible by \(2a\) and value at least \(\max(2^{16},16W_C)\) at the start of the component algorithm. Retain this reference along the recursion of Chuzhoy §8. Let \(b_*=\log_2\widehat W\) and \(\rho=\widehat W^{1/a}\). All \(\widehat W^{z/(2a)}\), \(\rho\), and \(\widehat W^{j/a}\) size and phase thresholds are integral powers of two. The large-case check takes \(b_*\) large enough as a function of \(a\) for \(\rho\ge16\), the explicit logarithmic-power prerequisites of §9.3, and \(\widehat W\ge2^{C a^2}\). It uses AlgSlow on an internal connected graph if \(\lfloor\log_2\lfloor\log_2 N^0(C)\rfloor\rfloor<C a\), increasing the absolute \(C\) to cover the parameter \(2/z\) in Theorem 8.1 for \(1\le z\le2a\). This costs \(O_a(N^0(C)\mu D\log^{O_a(1)}\widehat W)\) in that case. Log tests are used only for counts above two. In Alg3 take \(\mu'=\mu\widehat W^{1/(2a)}\). Indeed, §7.2 gives at most \((1+\mu)\lambda\,W^{O(1/\log\log W)}\) incidences per regular vertex, including its original and fake copies, where \(W\le2\widehat W\) and \(\lambda=O(b_*)\) when \(D\le\widehat W\). For \(\log_2\log_2\widehat W\ge Ca\) with sufficiently large absolute \(C\), this is bounded by \(\mu\widehat W^{1/(2a)}\), using an absolute constant for the bounded \(W\) convention. A smaller reference at fixed \(a\) likewise uses AlgSlow. In Alg3, \(N^0(\widehat H)\le N^0(C)/\rho\) gives both \(N^0(\widehat H)\mu'\le N^0(C)\mu\) and \(N^0(\widehat H)(\mu')^2\le N^0(C)\mu^2\). Phase 4 and the smaller scale cover calls retain the parent degree bound.

Set \(\widehat D=2^{30}D b_*^{10}\) and \(\eta=\widehat D\rho b_*^{C a}\) as in Chuzhoy §9, with an absolute exponent constant sufficiently large for the weighted routing and separator bounds. For the certified expansion of an expander constructed by the cut-matching proofs, use the greatest reciprocal power of two at most \(b_*^{-C'a}\) for an absolute sufficiently large \(C'\). Constants include both the last matching and the losses in pruning. The selected bounds can also be compared with the exact bounds of the particular cut-matching calls. When operating on actual terminal counts, use integer upper binary logarithms \(\max(2,\lceil\log_2 k\rceil)\). Positive integer loop and depth bounds are rounded up when they represent size or work allowances. Distance predicates compare integer distances with the floor of their threshold. Counts that are required to stay below a rational budget are compared by cross multiplication. Integer powers with exponents depending on \(a\) use repeated squaring. A ceiling of \(n^{1/a}\), as in Appendix E.2, uses integer-root comparisons. Phase budgets and restart tests use the certified rational expansion and congestion bounds. Use AlgSlow at a failed large-case size or \(D\) test, including \(D>\widehat W\).

There are two radius conventions needed by the short-threshold implementation:

\begin{enumerate}
\def\labelenumi{\arabic{enumi}.}
\tightlist
\item
  In §8.2.2 build the \(2^i\)-cover on the restriction of \(\widetilde H\) to edges of length at most \(2^i\). Mirror edge deletions and filter a supernode-splitting list by this same length predicate and by cluster incidence. Process a split in a restriction or cluster only if its filtered list is nonempty. Obtain these lists by scanning inverse memberships of the copied edges. An isolated regular label in this restriction can be kept as a singleton. The restriction preserves every closed \(2^i\)-ball, because every edge used by a path of length at most \(2^i\) has length at most \(2^i\). In InitNC at that scale the original regular vertices of each created cluster have distances at most \(1024\cdot 2^i b^4\) to one another in the induced cluster. Each fake regular vertex at initialization has degree at most one. Removing them preserves the paths between original regular vertices. Since \(2^i\le 4096D b_*^4\) and \(1024\cdot4096D b_*^4(b_*+1)^4 <\widehat D\), the original regular vertices fit in a single pseudocut ball. Adding at most \(|\widehat E|\le N^0(C)/\rho\) fake vertices gives the size bound \(2N^0(C)/\rho\) of Observation 8.2, including a cluster with no original regular vertex. Every later cluster starts as a subcluster. The proof of Claim 7.2 uses each positive subpath at a scale covering its entire length, including its fake boundary edges. The portions of length zero join copies of the same original terminal and need no arc. A positive integral length \(q\) has an incidence scale \(2^i\le2q\) even when the first scale is 2. Summing the two incidences per portion gives \(d_{\widehat H}\le4d_C\) for \(d_C\le2048D b_*^4\).
\item
  The contracted graph in §8.2.3 has incidence lengths up to \(2^\lambda\) where \(\lambda=\lceil\log_2(2048D b_*^4)\rceil\). Use recursive radius \(D'=\max(8D,2^\lambda)\) in Alg3, preserving all its incidences. Now \(D'\) bounds all child lengths as required by §3.1. A child flag with distance \(>1024D'b_*^4\) still implies a parent distance \(>1024D b_*^4\), by the preceding distance comparison. A child answer has length at most \(\alpha_{z-1}D'\), and replacing its incidences by cluster paths costs at most another \(\alpha_{z-1}\) factor. The extra factors from \(D'/D=O(b_*^4)\) are \(O(b_*^{12})\) for the \(D^3\) cost in Theorem 8.1 and \(O(b_*^4)\) in path length. Absorb them in the logarithmic factor per induction level of Theorem 8.1. Explicitly, the work uses an exponent \(c_0(a)+c z\) on \(b_*\) with \(c\) large enough for Alg2 and this extra factor. The path exponents can obey \(e_z=2e_{z-1}+C\), with \(e_0=O(a^2)\), so \(e_z=2^{O(a)}\) for \(z\le 2a\). Regular counts, dynamic degree bounds and all four phase charges are unaffected by the larger radius.
\end{enumerate}

At any disconnected initialization of MaintainCluster, including Alg3 and Phase 4, run one instance per initial component with the shared reference and parameters. The sum of their initial regular counts is the original initial count. Components containing no regular vertex can be omitted. A split stays within its initial component, since every copied edge was already incident to the old vertex. If two initial components still contain regular vertices, use representatives as an infinite-distance flag. Keep one representative per nonempty component on deletion lists. Within a component, its normal flags handle later disconnections. At a §6.2 restart, Theorems 5.1--5.2 can run on the disconnected current graph. Their proofs partition by weight, route between explicit terminals or report separated sets, and take bounded balls. These operations use infinite distances between components. The §5.5 dynamic expander routine is on the separately constructed witness, which is connected whenever it has at least two vertices. In disconnected or isolated cases charge graph scans to \(|V(C)|+|E(C)|\) and ever-issued records, bounded by \(O(N^0(C)\mu)\), including one record per isolated regular vertex. If the search for a pseudocut reaches the empty set, return that certified bad witness immediately and start §8.3 without Alg3 or Alg4.

For every current flag, store its two live regular IDs and hold the corresponding queries. At the cover interface apply ProcCut' to that pair. Its disjoint search balls and \(2D\)-wide overlap give the split of §4.2. Keep its new subcluster, and send the old cluster all edge deletions incident to its removed rows before sending the isolated vertex deletions. This removes one advertised endpoint. A regular row's covering pointer moves only under the \(2iD-D\) inner-ball test of §4.4. Its entire current \(D\)-ball is then in the chosen side, because each shortest path to that ball is inside the old covering cluster. Under a supernode split, an old covering cluster containing a \(D\)-ball also contains any copied path of that radius, by folding the path and mirroring its incidences. This proves the pointer invariant at settlement. At a recursive interface process all induced subcluster membership changes in order, copying contracted incidences before the old cluster's deletions. If Alg3 raises a flag then its two regular IDs are current original vertices in \(S\); \(1024D' b_*^4\ge 8192D b_*^4\) and Claim 7.2 make the pair a valid parent flag. Propagate that pair unchanged. Deletion of an original flagged vertex first removes its relevant incident edges, so it removes the terminal from \(S\) and from the contracted graph. Thus the mirrored parent sequence lowers this child flag unless that witness instance has already ended. Flags for disconnected components, Alg4 or another pending subinstance can queue their pairs. Drop a pending pair when an endpoint disappears or its instance retires. Otherwise deliver all further deletions while holding queries, and keep the pair until an endpoint is deleted. A subsequent flag of a surviving instance names a current pair. Process flags within the current cover before passing its settled contraction and before taking another external update or query. This needs only event lists and an endpoint-to-pending-pair index in the corresponding instance, with constant work per event apart from its dictionary operations. Each serviced split charges its own new cluster work as in Claim 4.8. Each served flag in an instance removes a regular vertex of that instance.

For the empty-witness phase of §8.3, valid edits preserve the bound \(N^0(C)/\rho\) on regular vertices in a \(\widehat D\)-ball in the graph with the current pseudocut removed. Use Lemma 8.4 with \(D^*=2048D b_*^4\), including the seed at distance zero in the first layer, and at most \(4b_*\) layers to find a good one. Failure of either growth test doubles its count once that count is positive. A zero edge count can first become positive only once, and both cumulative counts are at most \(\widehat W\). The charged work per accepted layer is proportional, up to logarithmic overhead, to its interior edges and vertices, including the seed when it is isolated. The proof's buffered layer ensures that two surviving undiscarded vertices at distance at least \(D^*\) in the working graph were at distance at least \(D^*\) before the temporary removal. Each group is inside a ball of radius \(8b_*D^*<\widehat D\) and has weight at most \(N^0(C)/\rho\). The groups cover at least half of the regular rows at the start of the iteration. Until all but one group empty, flag any two live representatives of different groups. Distances remain greater than \(1024D b_*^4\), and when the iteration ends at least a quarter of the rows present at its start have been removed if its start count was at least \(4N^0(C)/\rho\). There are \(O(\log(2+N^0(C)))\) iterations, including their first builds. Phase 4 starts with fewer than \(4N^0(C)/\rho\le \widehat W^{(z-1)/(2a)}\) regular rows, using \(\sqrt\rho\ge4\).

In the path-lifting step of §8.2.4, maintain for each incidence \((s,u(C'))\) the nonempty bucket of current original and fake copies of \(s\) in \(C'\). Enter or remove a membership using its cluster and original label. Copy the contracted incidence when splitting a subcluster with a nonempty bucket and remove the incidence when its bucket empties or its label leaves \(S\). This adds dictionary work per membership or edge scanned in Alg2. To query, take a representative of each of the two buckets at the child-path incidences and unfold the corresponding in-cluster path. Each fake edge maps to its current original edge. Its copy ID is preserved by Alg1 even after a supernode split. Recursive path work includes each child edge even when a recursive leg is empty. For the mixed scales, \(\sum_{\rm legs}D_{\rm leg}\le\alpha_{z-1}D'\) and \(D_{\rm leg}\ge1\). Induction with this sum bounds query work by \(D b_*^{2^{O(a)}}\). Serve queries after the current flag-lowering sequence is fully settled.

For the work induction, the total initial regular mass of Alg2 at each scale is \(O(N^0(C)\log(2+\widehat W))\) and \(2^i=O(Db_*^4)\). Summing the cubes of those scales adds at most \(O(D^3 b_*^{12})\). Alg3 obeys the squared-degree product bound above and \(D'/D=O(b_*^4)\). The remaining calls in Alg1 and Alg4 scan \(O(N^0(C)\mu)\) ever-issued records with logarithmic overhead and ES depth at most \(2\widehat D\). Phase 1 is charged in §9.3. The extra scale power at each induction step is absolute. Constants and a factor \(b_*^{c z}\) for sufficiently large absolute \(c\) absorb the sum of the children and their bookkeeping, giving \(O_a(N^0(C)\mu^2 D^3\widehat W^{O(1/a)} b_*^{O(a^2)+O(z)})\) over \(z\le2a\). Here the exponent coefficient in \(O(1/a)\) is absolute. This proves (A) after summing the cover cluster masses at \(\mu=2,D=3L\). \href{https://arxiv.org/pdf/2109.05621v1}{Chuzhoy v1, §§6--8 and Appendix D}

\hypertarget{transitive-graph-routines}{%
\subsection{Transitive graph routines}\label{transitive-graph-routines}}

The following are the imports reached by this short-threshold route. Their numerical inputs have lengths polynomially bounded in the cover size and the original threshold, with logarithmic word overhead for fixed \(a\).

\begin{itemize}
\item
  \textbf{Integer ES with splits.} For integer cutoff \(D^*\), define \(\mathrm{cap}(x)=\min(x,D^*+1)\). Keep labels initially exact by Dijkstra and heaps of keys \(\mathrm{cap}(\lambda(y)+\ell(x,y))\) at \(x\). The minimum of an empty heap is the sentinel. On an edge deletion drop its keys, enqueue affected endpoints, and repeatedly raise any nonsource label below its heap minimum to that minimum, updating keys and enqueuing affected vertices. Each label is a lower bound on its true capped distance, since every incident key is a lower bound on the corresponding route length and the minimum of the true incident keys is the true capped distance at a nonsource. Also \(\lambda(x)\) is at most its heap minimum throughout. These inequalities persist under a deletion and each raise. At settlement every finite nonsource label equals a key with a strictly smaller neighbor label, and following current minimizing keys gives an actual shortest source path. Taking the current key suffices even if a parent changed without changing the label. A removed isolated source ends that source's tree. No view using that source is queried again. For a nonempty split of \(u\) into a new \(u'\), paths between old vertices fold back to walks through \(u\), so their distances are unchanged. A shortest path to \(u'\) ends in one of its copied incidences and its prefix has the old shortest distance in the old graph. Initialize its label to the minimum actual copied key, even when \(u\) is the source, then insert the copied keys in the old heaps. All labels are exact after the split, so the new keys cannot lower any old label. Each vertex has at most \(D^*+1\) increases charged to all its incident edge IDs. Initialization, isolated deletion, key changes, and splits total \(\widetilde O((n_0+m_{\rm ever})(1+D^*))\), with any new vertex paid by a copied edge. This implementation applies also to source-augmented graphs and to ordinary ES with no splits. \href{https://arxiv.org/pdf/2109.05621v1}{Chuzhoy v1, Theorem 3.2 and Appendix B}
\item
  \textbf{Witness subdivision and that tree's topology.} In a §6.2 phase each witness terminal \(t_e\) corresponds to the specific edge ID \(e=(u,v)\) of length \(\ell(e)\) that was present at initialization, with \(u\) super and \(v\) regular. Its descendant copies inherit the same subdivision label. In the source-free search graph use a half-edge \((u',t_e)\) of length \(\ell(e)\) for each live copy \((u',v)\), and one half-edge \((t_e,v)\) of that length until the last copy dies. Nonwitness edges retain their own IDs and lengths. A supernode split is exactly an incidence copy in this graph, including when it copies several such half-edges. A deletion removes the named copy's half-edge and the shared half-edge only if it was the last copy. A split introduces no new shared half-edge. Add a fixed source with length-one edges to the retained witness terminals. Those links only disappear during the phase.

  Embeddings were computed on the initial subdivision graph. Preserve both half-edge IDs of each initial witness edge in the inverse lists, and remove every expander edge whose path used either half when the initial edge dies. Removing this union invalidates at most \(2\eta\) expander edges per base deletion. In particular every incident expander path of \(t_e\) uses one of its two initial half-edges, so the initial edge's death isolates it in the expander before pruning. The residual conductance invariant below excludes that isolate while a nonsingleton witness is in use. The separate singleton case has a zero base-deletion budget and rebuilds after the first charged deletion, before any query. Descendant copies may still be used by the search tree, and the fixed expander embedding uses only its surviving initial edge IDs. Process all pruning and spoke and base-edge deletions before the next query or flag. Under the map \(t_e\mapsto v\), discard collapsed half-edges and replace every other half-edge by its current base copy. Every unfolded walk then lies in the current \(C\). A live base path lifts into the source-free search graph with length at most twice its original length, and the regular end of a retained witness edge is adjacent to its terminal at distance \(\ell(e)\le D\). It follows that a regular vertex outside search depth \(\widehat D+1\) has base distance \(>(\widehat D-D)/2>1024D b_*^4\) from that end. There are \(O(N^0(C)\mu)\) initial labels and ever-issued edge IDs in the search graph. The bound used in the proof of Theorem 5.5 is the global incidence and edge-ID bound \(N^0(C)\mu\), which holds even at a high-degree supernode. \href{https://arxiv.org/pdf/2109.05621v1}{Chuzhoy v1, Theorem 5.5 and §6.2}
\item
  \textbf{Pseudocuts and weighted matching.} In Theorems 5.1--5.2 and Appendix E.1, all pseudocut weights are counts of regular vertices. The initial partition sorts these counts, and the graph modifications copy integral lengths. Before the large-case call, check the count prerequisites in E.5 and Theorem 5.2 with integer upper logarithms. The source's singleton-witness case applies when \(0<k\le\rho^2\). For fixed \(a\) the relevant logarithmic powers, including the constants in those prerequisites, are below \(\rho\) for sufficiently large \(\widehat W\). At smaller reference sizes the constant-size AlgSlow case applies. Subsequent pseudocut calls either give a witness or reduce the cut count by a fraction at least \(1/(\rho b_*^{O(a)})\). Thus there are \(O(\rho b_*^{O(a)}\log(2+|E|))\) iterations. Separation, weight bounds and witness sizes use counts and integer-distance searches. Appendix E.1's cut matching invokes its cut player and matching player \(O(\log(2+k))\) times on \(k\) terminals, followed by pruning. In E.11--E.12 use the inclusive path test at depth \(D'+2\), and remove an edge on reaching its integral allowance. For a matching call allowed congestion \(\eta\), use \(\lfloor\eta/\lceil\log_2(2+\kappa)\rceil\rfloor\) per halving round on \(\kappa\) source terminals. The positive allowance is at least half its unrounded value in the large-case calls. Charging each removed edge to its allowance increases the separator bound by at most a constant, paid by increasing the constant 16 in the congestion choice of Theorem 5.2. A stopping test with rational missing-match allowance \(z\) can use \(\lfloor z\rfloor\), since the counts are integers. The separator case still starts with more than \(z\) unmatched sources. \href{https://arxiv.org/pdf/2109.05621v1}{Chuzhoy v1, §5.3 and Appendix E.1}
\item
  \textbf{Cut player and the unweighted matching player.} Chuzhoy E.8 and E.17 invoke Theorems 4.1 and 3.2 of Chuzhoy--Gao--Li--Nanongkai--Peng--Saranurak v2. In the latter use a degree upper bound \(\Delta\ge1\) and \(b=\max(2,\lceil\log_2 n\rceil)\). At a routing round of Lemma 3.3, build its \(k\) ordinary ES structures with depth \(\ell+2\). Strip the two terminal spokes from a reported path and delete every edge of the resulting graph path from all routing copies. Remove the two selected spokes from this pair's tree. The matched vertices remain available for transit. Once a pair's search ends, subsequent deletions cannot recreate a qualifying path for that pair. Consequently its surviving terminal sets have mutual distance greater than \(\ell\) in the graph remaining after these edge deletions.

  Let \(U\) be the total number of sources at the start of the round, and suppose the round routes fewer than \(8bU/\ell^2\) paths. For any residual pair with both sides present, Claim 3.5 is implemented by exploring balls about its two surviving terminal sets. The balls of radius \(\lceil\ell/4\rceil\) are disjoint because \(2\lceil\ell/4\rceil\le\ell\). One has at most half the remaining vertices. If no earlier radius on that side has boundary less than \(8\Delta b/\ell\) times its size, the ball grows by at least \(1+8b/\ell\) each layer. For \(\ell\ge32\Delta b\), this growth over \(\lceil\ell/4\rceil\) layers exceeds \(2^b\), since \(\log(1+x)\ge x/(1+x)\) for the natural logarithm and \(x=8b/\ell\le1/4\). This is impossible. Remove an eligible ball, retaining it in the cut union \(Z\), and repeat in the remaining graph until the union has at least \(n/4\) vertices or each pair has lost an entire side. In the latter case, disjointness of the terminal sets implies \(|Z|\) is at least the number of sources left after routing, which is at least \(U/2\). In the former case, \(n/4\ge U/2\). Immediately before the last ball was added there were fewer than \(n/4\) vertices in the union, and the last ball had at most half the remaining vertices, so \(|Z|\le3n/4\). The boundary of \(Z\) in the residual graph is at most the sum of the removed balls' boundaries at removal, at most \(8\Delta b|Z|/\ell\). Restoring the routed edges adds fewer than \(\ell\cdot8bU/\ell^2\le16b|Z|/\ell\) edges to the boundary. Both sides have size at least \(U/2\), and the smaller side has size at least \(|Z|/3\). The original cut therefore has sparsity at most \(72\Delta b/\ell\). Each ball removes one complete residual terminal side, so there are at most \(k\) such searches. The argument includes degree zero and disconnected cases with the padded \(\Delta\).

  A routing round uses \(O(k(n+m)(1+\ell))\) times logarithmic overhead by the ES and scan bounds. Unmatched source counts fall by a factor at most \(1-8b/\ell^2\) on a successful round, so at most \(\ell^2\) rounds suffice even for \(z=0\). Congestion is at most \(\ell^2\), lengths at most \(\ell\), and total work is \(\widetilde O(k(n+m)(1+\ell)^3)\). In all our calls either \(n=O(m)\) or the padded bound \(n+m\) is charged and gives the same asymptotic estimate. For RouteOrCut choose \(\ell=\lceil144\Delta b/\psi\rceil\). With certified lower sparsity \(\psi\) the cut bound is strictly below \(\psi\), so the expander call in E.18 with \(z=0\) matches every source.

  The recursion for Theorem 4.1 uses this RouteOrCut, explicit constant-degree expanders from Theorem 2.4, fake-edge pruning from Lemma 2.9, grouping, and smaller cut-player calls. Its expander construction uses integer square root and modular arithmetic on pairs of residues. Drop loops and retain parallel edge IDs. When the composition has a single full block, retain its expander and route the extra vertices by the same one-pair matching. No edges between full blocks are required. Theorem 2.5's entropy is only a proof of the round count. For \(k\) vertices take \(q=a\) and the least power of two \(N_0'\) with \((N_0')^q\ge k\). Then \((N_0')^{q+1}\le 2^{q+1} k^{1+1/q}\), and \((q\log_2 N_0')^{-8q}\) is an exact rational bound. For sizes below the fixed-size prerequisites, a component of at least half the vertices is a \(1/k\)-expander. If all components are smaller, join components until at least \(k/4\) vertices are included, giving a balanced cut with zero boundary. The expansion bound used in E.9 can be reduced to at most \(1/k\) throughout that bounded range.

  In the base of Theorem 4.1 use \(z=\lfloor c n'/b^5\rfloor\) in the 17 matching calls. In the induction step use \(z=\lfloor c\psi_{q-1} n'/b^8\rfloor\) in the \(O(b)\) matching calls with the fixed cut threshold \(1/100\). Here \(b=\max(2,\lceil\log_2 n'\rceil)\), \(n'\) is the current cut-player subproblem size, and \(c>0\) is an absolute sufficiently small rational. For the base the full embedding has \(\eta=O(b^4)\), padded degree \(O(b)\) and constant sparsity. For the step the corresponding bounds are \(\eta=O(b^5)\), padded degree \(O(b)\) and sparsity at least an absolute constant times \(\psi_{q-1}/b\). Therefore \(17z\) or \(O(b)z\) satisfies the deletion budget of Lemma 2.9, including its padded degree, and \(4|F|\eta/\psi_{\rm emb}\le n'/3\), where \(\psi_{\rm emb}\) is the full embedding's certified sparsity. The separator side bound is at least half the unrounded allowance, since a cut is returned with more than \(\lfloor z_{\rm real}\rfloor\) unmatched sources. This also covers the case when the floor is zero. The extraction after dyadic pruning gives sparsity at least an absolute constant times \(1/b^5\) in the base and \(\psi_{q-1}/b^7\) in the step. With \(b\le q\log_2 N_0'\) and \(N_0'\) above a fixed absolute threshold these exceed the respective targets. Every recursive call's supported degree is bounded by the cut-matching round count as in Theorem 4.1, and the host graph for each round has a padded degree \(O(b)\), including an edgeless host. The work recurrence in §4 of that source is unchanged up to logarithmic factors, as each root call takes \(O(N_0'^{\,q+1}(q\log_2 N_0')^{O(q^2)})\) and the amount of extra work per graph record is logarithmic-power. \href{https://arxiv.org/pdf/1910.08025v2}{Balanced cut v2, Theorems 2.4--2.6, 3.2 and 4.1 and Lemma 2.9}
\item
  \textbf{Expander pruning.} Chuzhoy Theorem 2.2 and CGL+19 Theorem 2.6 use Saranurak--Wang Theorem 1.3, §A.4 and Algorithm 3. From a certified sparsity lower bound \(\varphi\le1\) and upper degree bound \(\Delta\ge4\), round \(\varphi/\Delta\) down to a reciprocal power of two \(\phi\) within a factor two. Thus \(\phi\le1/4\). The proof below maintains the original pruning bounds for \(k\le\phi m/10\). In CGL+19 Lemma 2.9 the graph with fake edges is connected by the embedding and \(m\ge n-1\). Its budget \(k\le\psi n/(32\Delta\eta)\) fits the rounded pruning input, since that input allows at least \(\psi(n-1)/(20\Delta\eta)\) for \(n\ge3\). The same volume bound gives \(|B|\le16k\Delta\eta/\psi\le n/2\), so its boundary argument gives the stated \(4k\eta/\psi\) bound on the number pruned. Use the direct small case for \(n<3\).

  In Chuzhoy E.1.1 use a round bound \(T_0=O(\log(2+k))\), a degree bound \(\Delta_*\ge\max(4,T_0)\), and missing-match allowance \(z=\widehat\varphi k/(100T_0\Delta_*)\) for the certified sparsity \(\widehat\varphi\le1/4\) of the full cut-matching graph. At most \(T_0z\) fake edges are accumulated. This is inside \(\phi m/10\) by \(m\ge k-1\) and retains at least \(k/2\) vertices by \(|P|\le 8T_0z/\phi\le 16k/100\). Its smaller constant in the separated sets' size is allowed in E.1. Other dynamic expander calls use an explicit vertex-based budget below.

  Fix the original loop-free multigraph \(G\), let \(A\) be the unpruned vertex set and let \(D_{\rm del}\) consist of actual deletions. Unit-Flow uses all original edges induced by \(A\), including those in \(D_{\rm del}\). The degree-preserving loops of the source proof are represented by using the original degree \(d_G(v)\) as the sink capacity. Each original internal edge has capacity \(2/\phi\) in either direction under the antisymmetric flow convention. An actual deletion does not remove a flow edge. Increase the source supply of each still-retained endpoint by \(15/(4\phi)\), which is integral for the chosen \(\phi\). Thus supply at \(v\in A\) is exactly
  \[
  \Delta(v)=\frac{2}{\phi}|E_G(v,V\setminus A)|+\frac{15}{4\phi}\deg_{D_{\rm del}}(v).
  \]
  Retain the warm-start flow and labels. The published A.4 displays \(4/\phi\) per endpoint but counts \(4/\phi\) per edge in its mass ledger. Use \(h=128b/\phi\), \(b=\max(2,\lceil\log_2(2m)\rceil)\), and test the growth of volume in original degrees. At termination of Unit-Flow, if there is excess then the top level has positive volume. A boundary edge of a level cut skipping a level carries exactly \(2/\phi\) out of the cut. If each of the top \(h-1\) layers fails the test of at most \((\phi/64)\operatorname{vol}_G(S)\) non-level-skipping boundary edges, growth makes \(\operatorname{vol}_G(S_1)\ge(1+\phi/64)^{h-1}>2m\). Indeed \(h-1\ge b(64/\phi)\) and \((1+\phi/64)^{64/\phi}>2\). Thus a qualifying set \(S\) exists and can be found by a downward bucket sweep, scanning vertices and their original incidences only as they enter the explored volume.

  To prune \(S\), restrict the flow and labels to \(A\setminus S\), restore prescribed boundary supply \(2/\phi\) per original edge to \(S\), and keep the deletion supply of all remaining vertices. The node mass at its retained end increases by \(2/\phi-f(u,v)\) on a boundary edge directed from \(u\in S\) to \(v\in A\setminus S\). This is nonnegative and is at most \(4/\phi\), with zero increase on saturated outgoing edges. Hence the prune creates at most \(\operatorname{vol}_G(S)/16\) and destroys at least \(\operatorname{vol}_G(S)\) of the existing node mass, since all vertices in \(S\) have saturated sinks. The warm-start invariants continue to hold. Up to deletion \(i\), total external injection \(I\le15i/(2\phi)\). Live mass never increases under a prune, so it is at most \(I\le3m/4\) when \(i\le\phi m/10\). If \(C_{\rm mass}\) is the total creation and \(V_{\rm pr}\) the total pruned volume, conservation and the bound on new mass give \(V_{\rm pr}\le C_{\rm mass}\le I+V_{\rm pr}/16\). Thus \(C_{\rm mass}\le 8i/\phi\) and \(V_{\rm pr}\le 8i/\phi\). Since live node mass equals total supply by antisymmetry, \(2|E_G(A,V\setminus A)|/\phi\le 8i/\phi\). Both volume and boundary estimates use original edges.

  For Proposition A.3, suppose a nontrivial cut of the actual retained graph \((G-D_{\rm del})[A]\) has residual boundary at most \((\phi/6)\) times the smaller original volume, attained by a side \(S\). Put \(Y=|E_G(S,A\setminus S)|\), \(J=|E_G(S,V\setminus A)|\), and \(d_S=\sum_{v\in S}\deg_{D_{\rm del}}(v)\). We have \(\operatorname{vol}_G(S)\le m\), so original conductance and the residual cut inequality give \(\phi\operatorname{vol}_G(S)\le Y+J\) and \(Y-d_S\le(Y+J)/6\). Therefore \(5Y\le J+6d_S\). Absorption plus outward capacity in the auxiliary graph is at most
  \[
  \operatorname{vol}_G(S)+2Y/\phi
     \le (3Y+J)/\phi\le(8J/5+18d_S/5)/\phi .
  \]
  The supplied \((2J+15d_S/4)/\phi\) exceeds this by at least \((2J/5+3d_S/20)/\phi\). The case \(J=d_S=0\) implies \(Y=0\), contrary to original connectivity. Thus a settled feasible preflow guarantees every actual retained cut has boundary at least \((\phi/6)\min(\operatorname{vol}_G(S),\operatorname{vol}_G(A\setminus S))\). Every original degree is positive in a nonsingleton connected input. This gives both ordinary conductance and edge expansion at least \(\phi/6\), and rules out a retained isolated vertex when there are two or more retained vertices.

  All pushes and warm starts use integers. Choose the minimum label among vertices with positive excess and label below \(h\), so an eligible receiving vertex has zero excess. A push sends at most the receiver's degree, the sender's excess and the edge's residual capacity, and a positive push sends at least one unit. For the A.2--A.3 work charge, write \(C_v\) for the cumulative newly created mass placed at \(v\). While its excess is positive no incoming push occurs, so its excess is at most \(d_G(v)+C_v\). The sum of original degrees of ever-relabelled vertices is bounded by \(C_{\rm mass}\), since a relabel requires a saturated sink and its absorbed units stay until the vertex is pruned. Each relabel raises the potential \(\sum_v\mathrm{ex}(v)\mathrm{label}(v)\) by at most \(d_G(v)+C_v\). Injection or the warm-start adjustment raises it by at most \(h\) times the newly created mass, and each push lowers it by at least its integer amount. Each vertex is relabelled at most \(h\) times. A cursor scan of a nonsuccessful incidence is charged once per sender label. An incidence whose neighbor is below the eligible level then has zero residual capacity, by the label invariant. A reverse push that makes its capacity positive requires the neighbor already to be above the sender's current level. Thus no skipped incidence becomes eligible again at that sender label. A successful repeated check is charged to its push. The sweep and prune scans cost \(O(V_{\rm pr})\) times dictionary overhead. Maintaining only nonempty label buckets avoids any array over the height range. The total updating work is \(\widetilde O(C_{\rm mass}h)=\widetilde O(k/\phi^2)\) after an \(O(n+m)\)-record initialization. \href{https://arxiv.org/pdf/1812.08958v4}{Saranurak--Wang v4, Theorem 1.3, §3.4.2 and Appendix A}
\item
  \textbf{Decremental paths inside witness expanders.} Use Chuzhoy Theorem 5.3 and Appendix E.2 with \(a\) levels, \(\gamma=\lceil n^{1/a}\rceil\), unweighted embeddings from E.13, bounded-depth ES trees and the preceding pruning routine. In E.13 route demands of the cut-matching graph on the predecessor expander by E.18, including the last matching from the complement of the returned large side into it. The odd-size case attaches one terminal by a shortest path in the predecessor. Use dyadic lower sparsities throughout. At stage \(i\) with initial count \(n_i=\lceil n/\gamma^i\rceil\), sparsity \(\varphi_i\) and degree bound \(\Delta_i\), allow at most \(\lfloor \varphi_i n_i/(64\Delta_i)\rfloor\) edge deletions before rebuilding. For \(n_i\ge2\) this fits the pruning budget and preserves at least half its vertices, by \(m_i\ge n_i-1\) and \(|S_i|\le16k\Delta_i/\varphi_i\). Use the certified retained sparsity \(\psi_i=\phi_i/8\) after pruning, where \(\phi_i\) is its dyadic conductance input. Thus one may take an integral routing length bound \(\ell_{i+1}=\lceil300\Delta_i b_*/\psi_i\rceil\), embedding congestion a constant times \(b_*\ell_{i+1}^2\), and tree depth \(4\ell_{i+1}\). These use the post-pruning sparsity for both initial builds and rebuilds. The stage leaves enough vertices for level \(i+1\) because \(\lceil n_i/2\rceil\ge n_{i+1}\). In Theorem 5.5 choose the phase-budget divisor large enough that the top expander stays inside this count. The degree bound of that top expander is \(O(b_*)\), as required by the divisor. When the next parent edge deletion would exceed a phase budget, consume the edit in the parent and rebuild its witness before queries. This also handles a budget of zero.

  At an intermediate level settle the predecessor's actual deletions and pruning first, and enumerate the set of edges that cease being live in the predecessor. This set contains both kinds of removal, deduplicated by edge ID, and drives the inverse embedding lists and ES trees. For the current embedding generation accumulate its distinct invalid edges. An invalid edge already counted does not count twice. If this batch would exceed the deletion budget, retire level \(i\) and its descendants and rebuild them from the settled predecessor, including the ES tree in the predecessor with the new terminal spokes. The retired generation's batch is discarded. If the budget permits the batch, deliver those deletions to the auxiliary pruning instance including actual deletions of edges whose endpoints were already pruned. Use only newly lost live edges and vertices for the next level. Remove spokes for any newly pruned terminals from the predecessor's tree. If a predecessor vertex is pruned, all paths incident to its terminal in a nonsingleton stage were invalidated. When the stage survives, its residual expansion excludes that isolated terminal. There are at least two retained vertices in this case: any positive deletion allowance requires \(n_i\ge64\Delta_i/\varphi_i\), and at least half are retained. Rebuilds choose all terminals from the settled predecessor. Stop at the first singleton level, which exists by \(\gamma^a\ge n\). If its terminal is pruned from the predecessor, choose another live terminal and rebuild the tree for that level. There are at most \(n_{i-1}\) such replacements in a predecessor stage.

  Within one predecessor stage, at most \(n_{i-1}\Delta_{i-1}\) predecessor edge IDs cease being live, even including all pruned incidences. Each generates at most the certified congestion \(\eta_i\) notifications to the generation present at the start of its batch. No edge of that batch needs processing in a new generation, whose paths were chosen on the settled predecessor. A completed level-\(i\) stage consumes at least \(\lfloor\varphi_i n_i/(64\Delta_i)\rfloor+1\ge\max(1,\varphi_i n_i/(64\Delta_i))\) notifications, including those that triggered rebuilding. Thus the number of level-\(i\) stages per predecessor stage is at most
  \[
     1+\frac{n_{i-1}\Delta_{i-1}\eta_i}
                  {\lfloor\varphi_i n_i/(64\Delta_i)\rfloor+1}
     =O_a((n_{i-1}/n_i)b_*^{O(a)}).
  \]
  The singleton bound \(1+n_{i-1}\) has this form as well. The ratios telescope, including the first builds, to at most \((n/n_i)b_*^{O(a^2)}\) times a branch constant stages at level \(i\). E.13 with the corrected matching player costs \(O_a((n_{i-1}+n_i^{1+O(1/a)})b_*^{O(a^2)})\) per stage: each of \(O(\log(2+n_i))\) rounds uses the cut player on \(n_i\) labels and at most one unweighted matching call on \(n_{i-1}\) vertices of degree \(\Delta_{i-1}\), with \(\ell_i=b_*^{O(a)}\). In the bounded-size case explicit expansion and shortest paths cost \(O_a(n_{i-1}\Delta_{i-1})\). This bound also covers the first inverse lists, initializing the ES tree on the predecessor, its deletions to depth \(4\ell_i\), and pruning at level \(i\). Unprocessed notifications and teardown are covered by the enumeration and allocation charges. Since \(n_{i-1}/n_i\le\gamma\le2n^{1/a}\) and \(n_i\le n\), the sum at each level is at most \(O_a(n^{1+O(1/a)} b_*^{O(a^2)})\), with absolute coefficient in \(O(1/a)\). The query recurrence is two ES legs since the live predecessor of degree \(\Delta_{i-1}\) has diameter at most \(\ell_i\), plus the expansion of one deeper query by a factor at most \(\ell_i\). The diameter claim follows from growth by \(1+\psi_{i-1}/\Delta_{i-1}\) until half its vertices are reached from each endpoint, using the chosen \(\ell_i\). Induction over at most \(a\) embedding levels yields path length and unfolding overhead \(b_*^{O(a^2)}\). \href{https://arxiv.org/pdf/2109.05621v1}{Chuzhoy v1, Theorems 5.3 and 5.5 and Appendix E.2}
\end{itemize}

For the layer test of Chuzhoy E.2--E.3 let \(b=\max(2,\lceil\log_2 k\rceil)\), use width \(2D''\), threshold \(1/(8b)\) and at most \(16b^2\) layers. Set \(D'=64D''b^2\). Balls of radius \(32D''b^2\) about the separated witness sets are disjoint, so one side has a count at most half throughout the search. If every layer on that side failed the growth test, its count would grow by \((1+1/(8b))^{16b^2}>2^b\ge k\). This uses the layer cutoff from the proof of E.3 with the same radius. In §6.2 use the exact phase budget \(\lfloor\varphi^* k/(C\eta b_*)\rfloor\), with \(C\ge128\Delta_*/b_*\) where \(\Delta_*=O(b_*)\) bounds the top expander degree and \(k\) is its vertex count. Indeed, each base-edge deletion causes at most \(2\eta\) top-expander deletions, and \(2\eta\varphi^* k/(C\eta b_*)\le \varphi^* k/(64\Delta_*)\). Check the budget before delivering the trigger deletion to the old witness. Include the trigger in the phase's charge. A completed type-2 phase with \(w=W_j(C)\) regular vertices initially thus consumes at least \(\varphi^* k/(C\eta b_*)=\Omega(w\varphi^*/(\rho^3\eta b_*))\) base deletions. For phases starting in \(w\in[2^t,2^{t+1})\), the total number of remaining base edge deletions after the first such start is \(O(2^{t+1}\mu)\). Therefore each weight band has \(O(\mu\rho^3\eta b_*/\varphi^*)\) completed phases. With \(\eta=\widehat D\rho b_*^{O(a)}\) and \((\varphi^*)^{-1}=b_*^{O(a)}\), the number over all bands, including an unfinished phase and the other initializations of §6.2, is \(O(1+\mu D\widehat W^{O(1/a)}b_*^{O(a)})\). Each type-2 build, including pseudocut search, all witness-stage builds and updates, embedding lists and ES, costs \(O_a(N^0(C)\mu D^2\widehat W^{O(1/a)}b_*^{O(a^2)})\). This also bounds the work of identifying the last bad phase. The type-1 call to AlgSlow uses at most \(\widehat W^{10/a}\) regular vertices, hence costs \(O(N^0(C)\mu D\widehat W^{10/a} b_*^{O(1)})\). Multiplication gives the Phase 1 bound used in §9.2.

Thus the short invocation is a finite graph algorithm. Initial graph sizes, edge counts, list counters, distances, labels and integral pruning masses are polynomially bounded in the cover size and \(D_C\), allowing a logarithmic-power factor for fixed \(a\). Along the at most \(2a\) MaintainCluster levels a radius increases by at most \(O(b_*^4)\) per level, and the expander path depths use only inverse logarithmic-power expansion. With cover size at most a fixed polynomial in \(M\) and \(D_C=O(\log M)\), counters, addresses and distance labels take \(O_a(\log(2+M))\) bits. Dyadic and rational thresholds use the displayed fixed integer powers and exact comparisons on words of at most logarithmic-power bit length. The counts in (A) come from the proof of Theorem 3.3, including all internal builds. The integer operations and the internal corrections add only logarithmic-power factors and constants to those charges. The exponent \(c_1/a\) has an absolute coefficient.

\end{document}
